% concatenated from main.tex
\documentclass{amsart}
\oddsidemargin 0pt
\evensidemargin 0pt
\marginparwidth 20pt
\textwidth 480pt

\textheight 720pt
\topmargin -60pt

%\usepackage[notref,notcite]{showkeys}   
%\usepackage[notref,notcite]{showkeys}  
\usepackage{pgfplots}
\usepackage{amsmath}
\pgfplotsset{compat=1.18}
\usepackage{sansmath} % optional: math font consistent with sans serif text
\sansmath
%\usepackage[nottoc,notocchapter]{tocbibind}
\usepackage{graphicx}
\usepackage{slashed}
\usepackage{enumerate}
\usepackage{dsfont}
%\usepackage{amsmath}
\usepackage{mathtools,cancel}
\usepackage[bottom]{footmisc}
\usepackage[scr=rsfso]{mathalfa}
\usepackage{tikz-cd}
\usepackage[hidelinks]{hyperref}
{
	
	\newtheorem{assumption}{Assumption}
}
\usepackage{amsfonts}
\usepackage{psfrag}
\usepackage{color}
\usepackage{mathtools}
\usepackage{pgf,tikz}
\usepackage{mathrsfs}
\usetikzlibrary{arrows}
\usepackage{fancyhdr}
\usepackage{amssymb}
\usepackage{slashed}
\usepackage{enumitem}

%%%%%%%
\usepackage[T1]{fontenc}
\usepackage[utf8]{inputenc}
%%%%%%%
%\usepackage{ntheorem}
\newtheorem*{theorem-non}{Theorem}

\pagestyle{plain}
\newtheorem{definition}{Definition}
\newtheorem{remark}{Remark}
\newtheorem{proposition}{Proposition}[section]
\newtheorem{problem}{Problem}


\newtheorem{theorem}{Theorem}[section]
\newtheorem{corollary}{Corollary}[section]
\newtheorem{lemma}{Lemma}[section]
\newtheorem{nono-theorem}{Theorem}
\newtheorem{conjecture}{Conjecture}[section]
\newtheorem{question}{Question}
\newtheorem{claim}{Claim}
%\renewcommand{\theequation}{\thesection.\arabic{equation}}

%\newcommand{\auth}{\textsc}
%\newcommand{\tit}{\textrm}
%\newcommand{\jou}{\textit}


\newcommand{\Yub}{\mathcal{Y}_{\ubar}}
\newcommand{\Lbar}{\underline{L}}
\newcommand{\Hbar}{\underline{H}}
\newcommand{\psig}{\psi_g}
\newcommand{\aln}{(\al \nabla)^i}
\newcommand{\ali}{a^{\frac{i}{2}}}
\newcommand{\scaletwoHprime}[1]{\lVert #1 \rVert_{L^2_{(sc)}(H_{u^{\prime}}^{(0,\ubar)}) }}
\newcommand{\twoSuprime}[1]{\lVert #1 \rVert_{L^2(S_{u^{\prime},\ubar})}}
\newcommand{\scaletwoHbarprime}[1]{\lVert #1 \rVert_{L^2_{(sc)}(\Hbar_{\ubar^{\prime}}^{(u_{\infty},u)}) }}
\newcommand{\tbeta}{\beta^R}
\newcommand{\nablaF}{\nabla^{i_5}}
\newcommand{\hnablaF}{\hnabla^{i_5}}
\newcommand{\tbetabar}{\betabar^R}
\newcommand{\hsp}{\hspace{.5mm}}
\newcommand{\omegabar}{\underline{\omega}}
\newcommand{\gslash}{\slashed{g}}
\newcommand{\psibar}{\underline{\psi}}
\newcommand{\psiKl}{\psi_{\mathfrak{Kl}}}
\newcommand{\psibarKl}{\psibar_{\mathfrak{Kl}}}
\newcommand{\dubarprime}{\hspace{.5mm} \text{d} \ubar^{\prime}}
\newcommand{\hnabla}{\nabla}
\newcommand{\intu}{\int_{u_{\infty}}^u}
\newcommand{\intubar}{\int_0^{\ubar}}
\newcommand{\del}{\partial}
\newcommand{\F}{{{F}^P}_Q}
\newcommand{\chihat}{\hat{\chi}}
\newcommand{\g}{\gamma}
\newcommand{\upr}{\lvert u^\prime \rvert}
\newcommand{\al}{a^{\frac{1}{2}}}
\newcommand{\chibar}{\underline{\chi}}
\newcommand{\chibarhat}{\underline{\hat{\chi}}}
\newcommand{\ubar}{\underline{u}}
\newcommand{\e}{\mathrm{e}}
\newcommand{\sumit}{\sum_{i_1+i_2+i_3=i}}
\newcommand{\sumitm}{\sum_{i_1+i_2+i_3=i-1}}
\newcommand{\sumiF}{\sum_{i_1+i_2+i_3+i_4+i_5=i}}
\newcommand{\sumiFi}{\sum_{i_1+i_2+i_3+i_4+i_5=i-1}}
\newcommand{\sumif}{\sum_{i_1+i_2+i_3+i_4=i}}
\newcommand{\sumifi}{\sum_{i_1+i_2+i_3+i_4=i-1}}
\newcommand{\sumifim}{\sum_{i_1+i_2+i_3+i_4=i-2}}
\newcommand{\sumjf}{\sum_{J_1+J_2+J_3+J_4=J}}
\newcommand{\sumjfj}{\sum_{J_1+J_2+J_3+J_4=J-1}}
\newcommand{\sumjfjm}{\sum_{J_1+J_2+J_3+J_4=J-2}}
%\newcommand{\tr}{\text{tr}}
\newcommand{\be}{\begin{equation}}
\newcommand{\ee}{\end{equation}}
\newcommand{\dtext}{\text{d}}
\newcommand{\bm}{\begin{align}*}
\newcommand{\enm}{\end{align}*}
\newcommand{\hess}{\operatorname{Hess}}
\newcommand{\supp}{\operatorname{supp}}
\newcommand{\scalep}[1]{\lVert{#1} \rVert_{L^p_{(sc)}(u,\underline{u})}}
\newcommand{\scaleinfinityS}[1]{\lVert{#1} \rVert_{\mathcal{L}^\infty_{(sc)}(S)}}
\newcommand{\scaleinfinitySu}[1]{\lVert{#1} \rVert_{L^\infty_{(sc)}(S_{u,\ubar})}}
\newcommand{\scaleinfinitySuprime}[1]{\lVert{#1} \rVert_{L^\infty_{(sc)}(S_{u^{\prime},\ubar})}}
\newcommand{\ScaleinfinitySuprime}[1]{\big\lVert{#1} \big\rVert_{L^\infty_{(sc)}(S_{u^{\prime},\ubar})}}
\newcommand{\scaleinfinitySuprimeubarprime}[1]{\lVert{#1} \rVert_{L^\infty_{(sc)}(S_{u^{\prime},\ubar^\prime})}}
\newcommand{\ScaleinfinitySuprimeubarprime}[1]{\big\lVert{#1} \big\rVert_{L^\infty_{(sc)}(S_{u^{\prime},\ubar^\prime})}}
\newcommand{\modulo}[1]{{\left|#1\right|}}
\newcommand{\norma}[1]{{\left\|#1\right\|}}
\newcommand{\Dist}{\operatorname{dist}}
\newcommand{\const}{\operatorname{const}}
\newcommand{\snr}[1]{\lvert #1\rvert}
\newcommand{\dv}{\operatorname{div}}
\newcommand{\vj}{v^{j}}
\newcommand{\tildetr}{\widetilde{\tr \chibar}}
\newcommand{\bx}{B_{r}(x_{0})}
\newcommand{\da}{D_{\alpha}}
\newcommand{\nr}[1]{\lVert #1 \rVert}
\newcommand{\gr}[1]{\mathbf{#1}}
\newcommand{\Lc}{\mathcal{L}}
\newcommand{\Y}{\mathrm{\Upsilon}}
\newcommand{\ubaruS}[3]{{L_{\ubar}^{#1} L_u^{#2} L_{\vphantom{\ubar}}^{#3}(S)}}
\newcommand{\nablap}{\nabla^{i_1}\psi_g^{i_2}}
\newcommand{\nablapp}{\nabla^{i_1}\psi_g^{i_2+1}}
\newcommand{\uubarS}[3]{{L_{u}^{#1} L_{\ubar}^{#2} L_{\vphantom{\ubar}}^{#3}(S_{u,\ubar})}}
\newcommand{\uubarSu}[3]{{L_{u}^{#1} L_{\ubar}^{#2} L_{\vphantom{\ubar}}^{#3}(S_{u,\ubar})}}
\newcommand{\scalefourS}[1]{\lVert{#1} \rVert_{L^4_{(sc)}(S)}}
\newcommand{\scaletwoS}[1]{\lVert{#1} \rVert_{L^2_{(sc)}(S)}}
\newcommand{\twoSu}[1]{\lVert{#1} \rVert_{L^2(S_{u,\ubar})}}
\newcommand{\pSu}[1]{\lVert{#1} \rVert_{L^p(S_{u,\ubar})}}
\newcommand{\inftySu}[1]{\lVert{#1} \rVert_{L^{\infty}(S_{u,\ubar})}}
\newcommand{\oneSu}[1]{\lVert{#1}\rVert_{L^1(S_{u,\ubar})}}
\newcommand{\fourSu}[1]{\lVert{#1} \rVert_{L^4(S_{u,\ubar})}}
\newcommand{\bespeq}{\begin{equation}\begin{split}}
\newcommand{\espeq}{\end{split}\end{equation}}
\newcommand{\scaletwoSu}[1]{\lVert{#1} \rVert_{L^2_{(sc)}(S_{u,\ubar})}}
\newcommand{\scaleoneSu}[1]{\lVert{#1} \rVert_{L^1_{(sc)}(S_{u,\ubar})}}
\newcommand{\ScaletwoSu}[1]{\big\lVert{#1} \big\rVert_{L^2_{(sc)}(S_{u,\ubar})}}
\newcommand{\scaletwoSuprime}[1]{\lVert{#1} \rVert_{L^2_{(sc)}(S_{u^\prime,\ubar})}}
\newcommand{\scaletwoSuzprime}[1]{\lVert{#1} \rVert_{L^2_{(sc)}(S_{u^\prime,0})}}
\newcommand{\ScaletwoSuprime}[1]{\big\lVert{#1} \big\rVert_{L^2_{(sc)}(S_{u^\prime,\ubar})}}
\newcommand{\scaleoneSuprimeubarprime}[1]{\lVert{#1} \rVert_{L^1_{(sc)}(S_{u^\prime,\ubar^\prime})}}
\newcommand{\scaletwoSuprimeubarprime}[1]{\lVert{#1} \rVert_{L^2_{(sc)}(S_{u^\prime,\ubar^\prime})}}
\newcommand{\ScaletwoSuprimeubarprime}[1]{\big\lVert{#1} \big\rVert_{L^2_{(sc)}(S_{u^\prime,\ubar^\prime})}}
\newcommand{\scaletwoSuubarprime}[1]{\lVert{#1} \rVert_{L^2_{(sc)}(S_{u,\ubar^\prime})}}
\newcommand{\scaletwoSuzubarprime}[1]{\lVert{#1} \rVert_{L^2_{(sc)}(S_{u_{\infty},\ubar^\prime})}}
\newcommand{\ScaletwoSuubarprime}[1]{\Big \lVert{#1} \Big \rVert_{L^2_{(sc)}(S_{u,\ubar^\prime})}}
\newcommand{\LpSu}[1]{\lVert #1 \rVert_{L^p(\Suu)}}
\newcommand{\LinftySu}[1]{\lVert #1 \rVert_{L^{\infty}(\Suu)}}
\newcommand{\LpHu}[1]{\lVert #1 \rVert_{L^p(\Hu)}}
\newcommand{\LpHbaru}[1]{\lVert #1 \rVert_{L^p(\Hbu)}}
\newcommand{\scalefourH}[1]{\lVert{#1} \rVert_{\mathcal{L}^4_{(sc)}(H)}}
\newcommand{\alphabar}{\underline{\alpha}}
\newcommand{\betabar}{\underline{\beta}}
\newcommand{\etabar}{\underline{\eta}}
\newcommand{\scalefourHbar}[1]{\lVert{#1} \rVert_{L^4_{(sc)}(\underline{H})}}
\newcommand{\scaletwoH}[1]{\lVert{#1} \rVert_{L^2_{(sc)}(H)}}
\newcommand{\scaletwoHbar}[1]{\lVert{#1} \rVert_{L^2_{(sc)}(\underline{H})}}
\newcommand{\scalefourHu}[1]{\lVert{#1} \rVert_{L^4_{(sc)}(H_u^{(0,\underline{u})})}}
\newcommand{\scalefourHbaru}[1]{\lVert{#1} \rVert_{L^4_{(sc)}(\underline{H}_{\underline{u}}^{(0,u)})}}
\newcommand{\scaletwoHzero}[1]{\lVert{#1} \rVert_{L^2_{(sc)}(H_{u_{\infty}}^{(0,\underline{u})})}}
\newcommand{\scaletwoHu}[1]{\lVert{#1} \rVert_{L^2_{(sc)}(H_u^{(0,\underline{u})})}}
\newcommand{\scaletwoHbarzero}[1]{\lVert{#1} \rVert_{L^2_{(sc)}(\underline{H}_{0}^{(u_{\infty},u)})}}
\newcommand{\scaletwoHbaru}[1]{\lVert{#1} \rVert_{L^2_{(sc)}(\underline{H}_{\underline{u}}^{(u_{\infty},u)})}}
\newcommand{\dl}[1]{\epsilon #1}
\newcommand{\dt}[1]{\epsilon \tilde{#1}}
\newcommand{\scalepinitialHbar}[1]{\lVert{#1} \rVert_{L^p_{(sc)}(u,0)}}
\newcommand{\scalepinitialH}[1]{\lVert{#1} \rVert_{L^p_{(sc)}(0,\underline{u})}}
%\newcommand{\br}{B_{r}(a)}
%\newcommand{\dbr}{\partial B_{r}(a)}
\newcommand{\Lloc}[1]{\mathbf{L^{#1}_{loc}}}
\newcommand{\C}[1]{\mathbf{C^{#1}}}
\newcommand{\Cc}[1]{\mathbf{C_c^{#1}}}
\newcommand{\Hutwo}[1]{\lVert #1 \rVert_{L^2(H_u)}}
\newcommand{\Hbarubartwo}[1]{\lVert #1 \rVert_{L^2(\Hbar_{\ubar})}}
\newcommand{\diver}{\operatorname{div}}
\newcommand{\omegad}{\omega^{\dagger}}
\newcommand{\omegabard}{\omegabar^{\dagger}}
\newcommand{\duprime}{\hspace{.5mm} \text{d}u^{\prime}}
\newcommand{\mubar}{\underline{\mu}}
\newcommand{\kappabar}{\underline{\kappa}}
\newcommand{\sh}{\tau_{hs}}
\newcommand{\Suu}{S_{u,\ubar}}
\newcommand{\Suo}{S_{u,0}}
\newcommand{\Sou}{S_{0,\ubar}}
\newcommand{\hs}{\tau_{-hs}}
\newcommand{\Dt}{D_{t}}
\newcommand{\ep}{\sh^{+}}
\newcommand{\eme}{\hs^{-}}
\newcommand{\utne}{\tilde{u}^{n}_{\varepsilon}}
\newcommand{\une}{u^{n}_{\varepsilon}}
\newcommand{\un}{u^{n}}
\newcommand{\ute}{\tilde{u}_{\varepsilon}}
\newcommand{\wu}{W^{1,p}_{u}}
\newcommand{\br}{B_{r}(0)}
\newcommand{\brp}{B_{r'}(0)}
\newcommand{\brs}{B_{r''}(0)}
\newcommand{\bs}{B_{s}(0)}
\newcommand{\bt}{B_{t}(0)}
\newcommand{\aupr}{\frac{a}{\upr^2}}
\newcommand{\haln}{(\al \hnabla)^i}
\newcommand{\RN}{\mathbb{R}^{N}}
\newcommand{\SN}{\mathbb{S}^{N-1}}
\newcommand{\N}{\mathbb{N}}
\newcommand{\gn}{G^{n}}
\newcommand{\nablat}{\nabla^{i_3}}
\newcommand{\hnablat}{\hnabla^{i_3}}
\newcommand{\nablaf}{\nabla^{i_4}}
\newcommand{\hnablaf}{\hnabla^{i_4}}
\newcommand{\Hodge}[1]{\prescript{*}{}{#1}}
\newcommand{\gne}{G^{n}_{\varepsilon}}
\newcommand{\ff}[1]{\int_{\br}F(D #1) \ dx}
\newcommand{\ffp}[1]{\int_{\brp}F(D #1) \ dx}
\newcommand{\intscaletwoSuubarprime}[1]{\int_0^{\ubar} \scaletwoSuubarprime{#1} \hspace{.5mm} \text{d} \ubar^\prime}
\newcommand{\intscaleTwoSuubarprime}[1]{\int_0^{\ubar} \ScaletwoSuubarprime{#1} \hspace{.5mm} \text{d} \ubar^\prime}
\newcommand{\fn}[1]{\int_{\br}F(D #1)+\frac{n^{2}}{4}\left(\snr{#1}^{2}-1\right)^{2}+\frac{1}{p}\snr{#1 -u}^{p}  \ dx}
\newcommand{\fne}[1]{\int_{\brp}F(D #1)+\frac{n^{2}}{4}\left(\snr{#1}^{2}-1\right)^{2}+\frac{1}{p}\snr{#1 - \ute}^{p}+\frac{\tilde{\varepsilon}}{2}\snr{D^{j} #1}^{2}  \ dx}
\newcommand{\fnp}[1]{\int_{\brp}F(D #1)+\frac{n^{2}}{4}\left(\snr{#1}^{2}-1\right)^{2}+\frac{1}{p}\snr{#1 - \ute}^{p}  \ dx}
\newcommand{\ffnp}[1]{\int_{\brp}F(D #1)+\frac{n^{2}}{4}\left(\snr{#1}^{2}-1\right)^{2}+\frac{1}{p}\snr{#1 - u}^{p}  \ dx}
\newcommand{\brl}{B_{\tilde{r}}(0)}
\newcommand{\gu}[1]{\left(n^{2}\left(\snr{#1}^{2}-1\right)#1 + \snr{#1 - \ute}^{p-2}(#1 - \ute) \right)}
\newcommand{\X}{\xi}
\newcommand{\E}{\eta}
\newcommand{\vn}{v_{j}}
\newcommand{\alphabarF}{\alphabar^F}
\newcommand{\rhoF}{\rho^F}
\newcommand{\sigmaF}{\sigma^F}
\newcommand{\alphaF}{\alpha^F}
\newcommand{\YM}{\mathcal{Y}}
\newcommand{\ub}{\underline{u}}
\newcommand{\Cb}{\underline{C}}
\newcommand{\Lb}{\underline{L}}
\newcommand{\Lh}{\hat{L}}
\newcommand{\Lbh}{\hat{\Lb}}
\newcommand{\phib}{\underline{\phi}}
\newcommand{\Phib}{\underline{\Phi}}
\newcommand{\Db}{\underline{D}}
\newcommand{\Dh}{\widehat{D}}
\newcommand{\Dbh}{\hat{\Db}}
\newcommand{\omb}{\underline{\omega}}
\newcommand{\omh}{\hat{\omega}}
\newcommand{\ombh}{\hat{\omb}}
\newcommand{\Pb}{\underline{P}}
\newcommand{\chib}{\underline{\chi}}
\newcommand{\chih}{\hat{\chi}}
\newcommand{\chibh}{\hat{\chib}}




\newcommand{\alb}{\underline{\alpha}}
\newcommand{\zeb}{\underline{\zeta}}
\newcommand{\beb}{\underline{\beta}}
\newcommand{\etb}{\underline{\eta}}
\newcommand{\Mb}{\underline{M}}
\newcommand{\oth}{\hat{\otimes}}




\def\a {\alpha}
\def\b {\beta}
\def\ab {\alphab}
\def\bb {\betab}
\def\nab {\nabla}





\def\ub {\underline{u}}
\def\th {\theta}
\def\Lb {\underline{L}}
\def\Hb {\underline{H}}
\def\chib {\underline{\chi}}
\def\chih {\hat{\chi}}
\def\chibh {\hat{\underline{\chi}}}
\def\omegab {\underline{\omega}}
\def\etab {\underline{\eta}}
\def\betab {\underline{\beta}}
\def\alphab {\underline{\alpha}}
\def\Psib {\underline{\Psi}}
\def\hot{\widehat{\otimes}}
\def\Phib {\underline{\Phi}}
\def\thb {\underline{\theta}}
\def\t {\tilde}
\def\st {\tilde{s}}







%%%%%%%%
%%%%%%%%
\def\d {\epsilon}
\def\f {\frac}
\def\i {\infty}
\def\l {\bigg(}
\def\r {\bigg)}
\def\S {S_{u,\underline{u}}}
\def\o{\omega}
\def\O{\Omega}
%\def\be{\begin{equation}\begin{split}}
%\def\en{\end{split}\end{equation}}

\def\od{\omega^{\dagger}}
\def\ombd{\underline{\omega}^{\dagger}}
\def\K{K-\frac{1}{|u|^2}}
\def\ut{\frac{1}{|u|^2}}
\def\Kb{K-\frac{1}{(u+\underline{u})^2}}
\def\M{\mathcal}
\def\p{\psi}

\def\D{\epsilon}
\def\T{\Theta}
\def\s{S_{u',\underline{u}'}}
\def\Hu{H_u^{(0,\underline{u})}}
\def\Hbu{\underline{H}_{\underline{u}}^{(u_{\infty},u)}}
%\def\ee{(\eta,\underline{\eta})}

\def\at{a^{\f12}}
\def\sigmac{\check{\sigma}}
\def\p{\psi}
\def\q{\underline{\psi}}
\def\ls{\leq}
\def\de{\epsilon}
\def\ls{\lesssim}
\def\oo{\Omega\mbox{tr}\chib-\frac{2}{u}}
\def\om{\omega}
\def\Om{\Omega}
%%%%%%%%
%%%%%%%%
\renewcommand{\div}{\mbox{div }}
\newcommand{\curl}{\mbox{curl }}
\newcommand{\trchb}{\mbox{tr} \chib}
\def\trch{\mbox{tr}\chi}


\newcommand{\Ls}{{\mathcal L} \mkern-10mu /\,}
\newcommand{\eps}{{\epsilon} \mkern-8mu /\,}
%%%%%%%%%
%%%%%%%%%






\newcommand{\tr}{\mbox{tr}}


\newcommand{\xib}{\underline{\xi}}
\newcommand{\psib}{\underline{\psi}}
\newcommand{\rhob}{\underline{\rho}}
\newcommand{\thetab}{\underline{\theta}}
\newcommand{\gammab}{\underline{\gamma}}
\newcommand{\nub}{\underline{\nu}}
\newcommand{\lb}{\underline{l}}
\newcommand{\mub}{\underline{\mu}}
\newcommand{\Xib}{\underline{\Xi}}
\newcommand{\Thetab}{\underline{\Theta}}
\newcommand{\Lambdab}{\underline{\Lambda}}
\newcommand{\vphb}{\underline{\varphi}}

\newcommand{\ih}{\hat{i}}
\newcommand{\ui}{u_{\infty}}
\newcommand{\shb}{L^2_{sc}(\underline{H}_{\ub}^{(u_{\infty},u)})}
%\newcommand{\sh}{L^2_{sc}(H_{u}^{(0,\ub)})}
\newcommand{\Rb}{\underline{\mathcal{R}}}
\newcommand{\tc}{\widetilde{\tr\chib}}

\newcommand{\tcL}{\widetilde{\mathscr{L}}}


\newcommand{\sRic}{Ric\mkern-19mu /\,\,\,\,}
\newcommand{\sL}{{\cal L}\mkern-10mu /}
\newcommand{\sLh}{\hat{\sL}}
\newcommand{\sg}{g\mkern-9mu /}
\newcommand{\seps}{\epsilon\mkern-8mu /}
\newcommand{\sd}{d\mkern-10mu /}
\newcommand{\sR}{R\mkern-10mu /}
\newcommand{\snab}{\nabla\mkern-13mu /}
\newcommand{\sdiv}{\mbox{div}\mkern-19mu /\,\,\,\,}
\newcommand{\scurl}{\mbox{curl}\mkern-19mu /\,\,\,\,}
\newcommand{\slap}{\mbox{$\triangle  \mkern-13mu / \,$}}
\newcommand{\scaletwoHuprime}[1]{\lVert #1 \rVert_{L^2_{(sc)}(H_{u^{\prime}}^{(0,\ubar)}    ) }  }
\newcommand{\sGamma}{\Gamma\mkern-10mu /}
\newcommand{\somega}{\omega\mkern-10mu /}
\newcommand{\somb}{\omb\mkern-10mu /}
\newcommand{\spi}{\pi\mkern-10mu /}
\newcommand{\sJ}{J\mkern-10mu /}
\renewcommand{\sp}{p\mkern-9mu /}
\newcommand{\su}{u\mkern-8mu /}
\newcommand{\m}{\mu}
\newcommand{\n}{\nu}
%%%%	
\newcommand\restri[2]{{% we make the whole thing an ordinary symbol
		\left.\kern-\nulldelimiterspace % automatically resize the bar with \right
		#1 % the function
		%\vphantom{\big|} % pretend it's a little taller at normal size
		\right|_{#2} % this is the delimiter
}}
%\newcommand{\paola}[1]{\textcolor{blue}{#1}}

\definecolor{ffqqqq}{rgb}{1.,0.,0.}
\definecolor{uuuuuu}{rgb}{0.26666666666666666,0.26666666666666666,0.26666666666666666}

%\DeclareMathOperator{\argmin}{argmin}
\DeclareMathOperator{\sgn}{sgn}
\DeclareMathOperator{\dist}{dist}
\DeclareMathOperator{\diam}{diam}
%\DeclareMathOperator{\proj}{Proj}\includegraphics[]{../../../../Desktop/NiceMunich.pdf}


\delimitershortfall=-0.1pt
	
\makeatletter
\def\ps@pprintTitle{%
  \let\@oddhead\@empty
  \let\@evenhead\@empty
  \let\@oddfoot\@empty
  \let\@evenfoot\@oddfoot
}
\def\@author#1{\g@addto@macro\elsauthors{\normalsize%
    \def\baselinestretch{1}%
    \upshape\authorsep#1\unskip\textsuperscript{%
      \ifx\@fnmark\@empty\else\unskip\sep\@fnmark\let\sep=,\fi
      \ifx\@corref\@empty\else\unskip\sep\@corref\let\sep=,\fi
      }%
    \def\authorsep{\unskip,\space}%
    \global\let\@fnmark\@empty
    \global\let\@corref\@empty  %% Added
    \global\let\sep\@empty}%
    \@eadauthor={#1}
}
\makeatother
%%%%%%%%%%%%%%%%%%
\delimitershortfall=-0.1pt

\begin{document}

	
\title{A large data result for vacuum Einstein's equations}%/ Dynamical formation of the Einstein-Yang-Mills blackholes/trapped surfaces?}


\author{Puskar Mondal}\footnote{ e-mail:pushkarmondal@gmail.com}
%\author{Shing-Tung Yau}
%\author{Shing-Tung Yau}

%\date{October 2025}

\maketitle

	



\begin{abstract}
We prove a global well-posedness and asymptotic convergence theorem for the \((3+1)\)-dimensional vacuum Einstein equations with positive cosmological constant \(\Lambda\) on globally hyperbolic spacetimes \(\widetilde M \cong M \times \mathbb R\), where \(M\) is a closed three-manifold of negative Yamabe type. In constant-mean-curvature transported spatial coordinates, an open set of large initial data gives rise to future-global solutions whose renormalized spatial metrics converge smoothly to a limiting metric of constant negative scalar curvature. The key new ingredient is an integrable damping mechanism, induced by the cosmological constant in this gauge and absent in the \(\Lambda=0\) vacuum problem, which yields time-integrable decay for the nonlinear evolution. As a consequence, the Einstein--\(\Lambda\) flow does not in general canonically encode the Thurston geometrization of the underlying three-manifold. This confirms a conjecture of Ringstr\"om on the asymptotic topological indistinguishability of large-data Einstein--\(\Lambda\) dynamics. An analogous theorem is also proved for manifolds of positive Yamabe type, under an additional technical hypothesis.
\end{abstract}









\setcounter{tocdepth}{2}
{\hypersetup{linkcolor=black}
\small
\tableofcontents
}


\section{Introduction}
\noindent We will study the vacuum Einstein's equations including a positive cosmological constant $\Lambda$. We are given a $n+1$-dimensional ($n=3$ in this article) $C^{\infty}$ globally hyperbolic connected  Lorentzian manifold $(\widetilde{M},\widehat{g})$ with signature $-+++$. Einstein's equations on $(\widetilde{M},\widehat{g})$ reads 
\begin{align}
\label{eq:first}
 \text{Ric}[\widehat{g}]-\frac{1}{2}R[\widehat{g}]\widehat{g}+\Lambda \widehat{g}=0.   
\end{align}
Global hyperbolicity implies the existence of a Cauchy hypersurface and in particular $\widetilde{M}= M\times \mathbb{R}$ with $M$ being diffeomorphic to a Cauchy hypersurface. In this article, we will choose $M$ to be of closed negative Yamabe type (a topological invariant, the so-called $\sigma$ constant $\sigma(M)\leq 0$). By definition, these admit no Riemannian
metric $g$ having scalar curvature $R(g)\geq 0$ everywhere. A closed 3-manifold $M$ is of negative Yamabe type if and only if it lies in one of the
following three mutually exclusive subsets:
\begin{itemize}
\item[(1)] $M$ is hyperbolizable (that
admits a hyperbolic metric),
\item[(2)] $M$ is a non-hyperbolizable $K(\pi,1)$ manifold of non-flat type (the six flat $K(\pi,1)$ manifolds are of zero Yamabe type), here $K(\pi,1)$ is of Eilenberg-MacLane type and by definition $\pi_{1}(M)=\pi$ and $\pi_{i}=0,~i\geq 2$-the universal cover of $K(\pi,1)$ is contractible and known to be diffeomorphic to $\mathbb{R}^{3}$ \cite{porti2008geometrization}, and
\item[(3)] $M$ has a nontrivial connected sum decomposition (i.e., $M$ is composite)
\begin{align}
\label{eq:negative}
M\approx (\mathbf{S}^{3}/\Gamma)_{1}\#..\#(\mathbf{S}^{3}/\Gamma)_{k}\#(\mathbf{S}^{2}\times \mathbf{S}^{1})_{1}\#..\#(\mathbf{S}^{2}\times \mathbf{S}^{1})_{l}\#\mathbf{K}(\pi,1)_{1}\#..\#\mathbf{K}(\pi,1)_{m}  
\end{align}
in which at least one factor is a $K(\pi, 1)$ manifold i.e, $m\geq 1$ \cite{schoenyau}. In this case the $K(\pi,1)$
factor may be either of flat type or hyperbolizable or non-hyperbolizable non-flat type. 
\end{itemize}
The six flat manifolds
comprise by themselves the subset of zero Yamabe type. These admit metrics
having vanishing scalar curvature (the flat ones) but no metrics having strictly
positive scalar curvature. Here $\Gamma\subset SO(4)$ acts freely and properly discontinuously on $\mathbb{S}^{3}$. Also note that $M\#\mathbb{S}^{3}\approx M$ for any 3-manifold $M$. It is known that every prime $K(\pi,1)$ manifold is decomposable
into a (possibly trivial but always finite) collection of (complete, finite volume)
hyperbolic and graph manifold \footnote{A compact oriented $3-$manifold $M$ is a graph manifold if $\exists$ a finite disjoint collection of embedded $2-$tori $\{T_{j}\}\subset M$ such that each connected component of $M-\cup_{j}T_{j}$ is the total space of circle bundle over surfaces} components. 

\noindent Our goal is to study the long-term behavior of Einstein's equations \ref{eq:first} on spacetimes $\widetilde{M}=M\times \mathbb{R}$ with $M$ being general $3-$ manifolds of negative Yamabe type as in \ref{eq:negative}. An issue closely related to Penrose's weak cosmic censorship \cite{penrose1999question} in the cosmological context (compact spatial topology where the notion of null infinity is not defined) may be phrased as follows    

\begin{conjecture}[No-Naked Singularity Conjecture for Cosmological Spacetimes]
\label{conjecture1}
Let $(\widetilde{M}, \widehat{g})$ be a smooth, time-oriented, globally hyperbolic Lorentzian $(3+1)$-dimensional spacetime with $\widetilde{M} \cong M \times \mathbb{R}$, where $M$ is a closed, connected, oriented three-manifold. Suppose $\widehat{g}$ satisfies the Einstein field equations,
\[
\operatorname{Ric}_{\widehat{g}} - \tfrac{1}{2} R_{\widehat{g}} \widehat{g} + \Lambda \widehat{g} =\mathcal{T},
\]
where $\Lambda \in \mathbb{R}$ is the cosmological constant and $\mathcal{T}$ is a smooth energy-momentum tensor obeying the dominant energy condition and representing physically reasonable matter and radiation content.

\noindent Then, for an open and dense set of smooth initial data prescribed on a Cauchy hypersurface $M \times \{t_0\}$ satisfying the Einstein constraint equations, the maximal globally hyperbolic development $(\widetilde{M}, \widehat{g})$ does not develop a future naked singularity; that is, any future singularity (if it forms) is not visible to any future-directed timelike curve originating from the initial hypersurface.
\end{conjecture}

\begin{remark}
 Notice that it is very unlikely to obtain global control of a solution with arbitrary initial data in a fully generic framework since Einstein's equations are quasi-linear  \textit{hyperbolic}  
\end{remark}

\noindent In this article, we focus on the vacuum case with a positive cosmological constant, i.e., $\mathcal{T}=0$ and $\Lambda>0$.
\subsection{Background Literature}

\noindent The future stability problem for cosmological solutions to the Einstein equations has been extensively studied in the small-data regime, particularly when the spatial manifold $M$ lies in the negative Yamabe class and admits a hyperbolic metric.

\noindent The foundational result in this context is due to Andersson and Moncrief \cite{andersson2004future}, who established the nonlinear future stability of the Milne universe,
\[
\widehat{g} = -dt^2 + \frac{t^2}{9}\eta, \quad \text{Ric}[\eta] = -\frac{2}{9}\eta,
\]
on manifolds of the form $\mathbb{H}^3/\Gamma \times \mathbb{R}$, where $\Gamma \subset \mathrm{SO}^{+}(1,3)$ is cocompact and acts freely and properly discontinuously on $\mathbb{H}^3$. Their result applies to vacuum solutions with spatial topology restricted to eliminate moduli of flat perturbations. The Milne model arises as a quotient of the interior of the future light cone in Minkowski space, and the proof employs a generalized energy method under small perturbations of the hyperbolic metric.

\noindent This analysis was extended to higher-dimensional spacetimes ($n+1 \geq 4$) in \cite{andersson2011einstein}, assuming that the spatial manifold admits a strictly stable negative Einstein metric. The proof relies on the decay of an energy functional and the spectral gap of the Lichnerowicz Laplacian associated to the Einstein background.

\noindent Subsequent work has incorporated various matter models. Andersson and Fajman \cite{andersson2020nonlinear} established the nonlinear future stability of the Milne model within the Einstein–massive Vlasov system. In a higher-dimensional Kaluza–Klein setting, Branding, Fajman, and Kröncke \cite{KK} proved future stability for $(\mathbb{H}^3/\Gamma \times \mathbb{T}^d) \times \mathbb{R}$ with Kaluza–Klein reduction, reducing the system to a coupled Einstein–Maxwell–wave map system.

\noindent Further stability results in the small data regime exist for Einstein equations coupled with a positive cosmological constant \cite{fajman2020stable, fajman2018cmc, mondal2019attractors}, relativistic fluids \cite{oliynyk2016future, oliynyk2021future, hadvzic2015global, lefloch2021nonlinear, mondal2024nonlinear}, and scalar fields \cite{wang, fajman2021attractors}. These developments confirm the validity of Conjecture~\ref{conjecture1} for small perturbations of homogeneous background spacetimes with hyperbolic spatial sections.

\noindent In contrast, the asymptotically flat case has a distinct lineage, beginning with the monumental work of Christodoulou and Klainerman \cite{christodoulou1993global}, who proved the global nonlinear stability of Minkowski spacetime. This framework has since been extended in various directions \cite{lindblad2010global, lefloch2016global, fajman2021stability, taylor, mondal2024global}, focusing on vacuum and matter models under asymptotically flat conditions.

 


\subsection{Main Result}

\noindent We consider the Einstein vacuum equations with positive cosmological constant $\Lambda > 0$ on $(3+1)$-dimensional globally hyperbolic spacetimes $\widetilde{M} \cong M \times \mathbb{R}$, where $M$ is a closed, connected, oriented $3$-manifold of negative Yamabe type. Importantly, we do not assume that $M$ admits a negative Einstein metric; in particular, $M$ may contain graph manifold summands in its prime decomposition, cf.~\eqref{eq:negative}, for which no Einstein metric exists.

\noindent In contrast to previous work, we make no smallness assumption on the initial data. We fix the constant mean curvature (CMC) transported spatial gauge. The time parameter is taken to be the mean extrinsic curvature $\tau := \operatorname{tr}_{g}(k)$, and we consider solutions in the expanding direction, i.e., $\tau < 0$ and increasing. 

\noindent To capture the asymptotic geometry, we introduce a natural geometric rescaling by the function $\varphi^2 := \tau^2 - 3\Lambda > 0$ (see claim \ref{important}) and analyze the rescaled evolution system \eqref{eq:cmc1}--\eqref{eq:cmc2}. 

\noindent To quantify the deviation of the evolving geometry from hyperbolicity, we define the \emph{obstruction tensor}
\[
\mathfrak{T}[g] := \operatorname{Ric}[g] - \tfrac{1}{3} R(g) g,
\]
which measures the failure of $g$ to be Einstein. In three spatial dimensions, by Mostow rigidity, $\mathfrak{T}[g] = 0$ implies that $g$ is a hyperbolic metric, uniquely determined up to isometry. For generic $M$ of negative Yamabe type, $\mathfrak{T}[g] \neq 0$ for all $g$.

\noindent We impose the transported spatial gauge (see Definition~\ref{spatial}), in which the pair $(\Sigma, \mathfrak{T})$ satisfies a manifestly hyperbolic evolution system. This is coupled with an elliptic equation for the lapse function determined by the CMC condition. In this formulation, we establish the following global convergence theorem for the Einstein flow in the large data regime on general negative Yamabe slices.



\begin{theorem}[Global Well-posedness: $\Lambda > 0$, $\sigma(M) \leq 0$]
\label{main}
Let $(\widehat{M}^{3+1}, \widehat{g})$ be a globally hyperbolic Lorentzian spacetime satisfying the Einstein vacuum equations with positive cosmological constant $\Lambda > 0$, and suppose that $\widehat{M}$ admits a constant mean curvature (CMC) foliation by compact spacelike hypersurfaces diffeomorphic to a closed $3$-manifold $M$. Assume furthermore that $M$ is of negative Yamabe type, i.e., $\sigma(M) \leq 0$.

\noindent Fix a smooth background Riemannian metric $\xi_0$ on $M$ and a constant $C > 1$. For any initial energy quantity $\mathcal{I}^{0} > 0$, there exists a constant $a = a(\mathcal{I}^{0}) > 0$, so that $\mathcal{I}^{0} e^{-a/10} < 1$.

\noindent Let $(g_0, \Sigma_0)$ be an initial data set verifying the Einstein constraint equations at initial CMC time $T_0 = a>0$, written in CMC-transported spatial coordinates and satisfying:
\begin{align}
\label{eq:1}
C^{-1} \xi_0 \leq g_0 \leq C \xi_0,
\end{align}
\begin{align}
\label{eq:2}
\sum_{I = 0}^{3} \| \nabla^I \Sigma_0 \|_{L^2(M)} + \sum_{I = 0}^{2} \left( \| \nabla^I \mathfrak{T}[g_0] \|_{L^2(M)} + \| e^a \nabla^I \Sigma_0 \|_{L^2(M)} \right) \leq \mathcal{I}^0,
\end{align}
where $\mathfrak{T}[g_0]$ denotes the renormalized trace-free spatial Ricci curvature tensor of $g_0$.

\noindent Then, the Einstein-$\Lambda$ evolution equations admit a unique classical solution
\[
T \mapsto (g(T), \Sigma(T)) \in \mathcal{C}^\infty([T_0, \infty) \times M)
\]
in CMC-transported spatial coordinates, satisfying the constraint equations at each slice $T$ and obeying the following uniform a priori estimates for all $T \in [T_0, \infty)$:
\begin{align}
\sum_{I = 0}^{3} \| \nabla^I \Sigma(T) \|_{L^2(M)} + \sum_{I = 0}^{2} \left( \| \nabla^I \mathfrak{T}[g(T)] \|_{L^2(M)} + \| e^T \nabla^I \Sigma(T) \|_{L^2(M)} \right) \leq C_1(1 + \mathcal{I}^0),
\end{align}
\begin{align}
C_2^{-1} g_0 \leq g(T) \leq C_2 g_0.
\end{align}
Here, $C_1, C_2 > 0$ are numerical constants depending only on the universal geometric and analytic data of the problem (e.g., Sobolev constants of $(M, g_0)$ and the constants in the structure equations, dimension of $M$), but independent of $T$. The developed spacetime is future geodesically complete.

\noindent Moreover, the solution $(g(T), \Sigma(T))$ converges in the $C^\infty$ topology, as $T \to \infty$, to a limiting Riemannian metric $\widetilde{g}$ on $M$ of pointwise constant negative scalar curvature, in the sense that
\[
\Sigma(T) \to 0 \quad \text{and} \quad g(T) \to \widetilde{g}, \quad \text{as } T \to \infty,
\]
with convergence holding in all Sobolev norms. In particular, the spacetime $(\widehat{M}, \widehat{g})$ admits a future-complete CMC foliation asymptotic to a constant negative scalar curvature slice.
\end{theorem}



%\begin{remark}
 %The choice of $\gamma$ as stated in the theorem \ref{main} is not optimal. The optimal value of $\gamma$ is $1$. This is discussed in section \ref{energy}   
%\end{remark}

 
\begin{remark}
Note that the allowed initial data size $\mathcal{I}^{0}$ is modulated by the condition $\mathcal{I}^{0}e^{-a/10}<1$ for initial time $T_{0}=a>0$ i.e., when $a$ is small, then the allowed data size is $O(1)$ but when $a\gg 1$, $\mathcal{I}^{0}$ can be very large $O(e^{a/10})$.    
\end{remark}

\begin{remark}
The result of the convergence of an arbitrary metric to the constant negative scalar curvature metric can be compared to the \textbf{Yamabe flow} on negative Yamabe manifolds. Of course, the convergence of the Yamabe flow on manifolds with $\sigma(M)\leq 0$ is trivial compared to its positive counterpart. Contrary to the Einstein-$\Lambda$ flow, the Yamabe flow is of energy-critical parabolic nature and endowed with monotonic entities (e.g., similarly, Yang-Mills heat flow \cite{YM} is of parabolic nature). Moreover, the limit constant scalar curvature metric obtained in the framework of Einstein's equations is \textbf{not} conformal to the initial metric.  
\end{remark}




\noindent \subsection{Heuristics on the genericity of the initial data}\label{genericity}
We comment on the notion of ``genericity'' of the initial data from a heuristic yet technically relevant viewpoint. The compact spatial manifold $M$ under consideration is of negative Yamabe type and, in particular, admits no Riemannian metric with nontrivial global Killing fields. Indeed, any compact $3$--manifold that supports a continuous isometry group necessarily admits an effective $SO(2)$--action. The classification theorem of Fischer--Moncrief~\cite{fischer1996quantum} implies that no nontrivial connected sum of negative Yamabe type supports such an action. Among prime manifolds, only a distinguished subclass of Seifert fibered spaces---specifically those whose fundamental group has nontrivial infinite cyclic center---can admit an $SO(2)$ symmetry. Even in these exceptional geometries, such continuous symmetries arise only in a nongeneric subset of the moduli space of initial data. Thus, from the perspective of geometric constraints, the absence of continuous symmetries is the generic situation.

\medskip

\noindent Although the topology of \(M\) leaves substantial freedom in the choice of initial data, the hypotheses of Theorem~\ref{main} still impose a definite restriction on how such data may concentrate. To describe this heuristically, we associate to the pair \((\Sigma,\mathfrak T)\) amplitude--wavelength parameters
\[
(A_{\Sigma},L_{\Sigma}),\qquad (A_{\mathfrak T},L_{\mathfrak T}),
\]
where \(\Sigma\) is the trace-free part of the second fundamental form and \(\mathfrak T\) is the renormalized trace-free Ricci tensor appearing in Theorem~\ref{main}. The quantities \(L_{\Sigma}\) and \(L_{\mathfrak T}\) should be interpreted as characteristic spatial scales, while \(A_{\Sigma}\) and \(A_{\mathfrak T}\) represent dimensionally normalized amplitudes. This discussion is purely heuristic and is intended only to indicate the concentration regimes compatible with the norm bounds of the theorem.

\noindent In three spatial dimensions, a field of amplitude \(A\) concentrated at wavelength \(L\) has squared \(H^{k}\)-size of order
\[
|A|^{2}\sum_{j=0}^{k} L^{3-2j}.
\]
Consequently, the control required in Theorem~\ref{main} yields, at the highest order for \(\Sigma\),
\begin{equation}\label{eq:highercontrol_refined}
L_{\Sigma}^{-3}|A_{\Sigma}|^{2}
+L_{\Sigma}^{-1}|A_{\Sigma}|^{2}
+L_{\Sigma}|A_{\Sigma}|^{2}
+L_{\Sigma}^{3}|A_{\Sigma}|^{2}
\lesssim (\mathcal I^{0})^{2},
\end{equation}
whereas the lower-order weighted bounds for \(\Sigma\) give
\begin{equation}\label{eq:lowercontrol_refined}
L_{\Sigma}^{-1}|A_{\Sigma}|^{2}
+L_{\Sigma}|A_{\Sigma}|^{2}
+L_{\Sigma}^{3}|A_{\Sigma}|^{2}
\lesssim e^{-2a}(\mathcal I^{0})^{2}.
\end{equation}
Similarly, for \(\mathfrak T\) one obtains
\begin{equation}\label{eq:ampl_refined}
L_{\mathfrak T}^{-1}|A_{\mathfrak T}|^{2}
+L_{\mathfrak T}|A_{\mathfrak T}|^{2}
+L_{\mathfrak T}^{3}|A_{\mathfrak T}|^{2}
\lesssim (\mathcal I^{0})^{2}.
\end{equation}
Moreover, the smallness assumption in Theorem~\ref{main} allows \(\mathcal I^{0}\) to grow with \(a\), subject only to
\[
\mathcal I^{0}e^{-a/10}\ll 1.
\]
In particular, one may heuristically regard \(\mathcal I^{0}\) as being as large as \(e^{a/10}\), up to a fixed small factor.

\noindent The contrast between \eqref{eq:highercontrol_refined} and \eqref{eq:lowercontrol_refined} is important. The term
\[
L_{\Sigma}^{-3}|A_{\Sigma}|^{2}
\]
shows that \(\Sigma\) may carry nontrivial high-frequency concentration at the top order, but only under the severe lower-order restriction imposed by \eqref{eq:lowercontrol_refined}. Indeed, if one writes
\[
A_{\Sigma}\sim e^{-\alpha a},
\qquad
L_{\Sigma}\sim e^{-\beta a},
\]
and takes \(\mathcal I^{0}\sim e^{a/10}\), then compatibility of
\eqref{eq:highercontrol_refined}--\eqref{eq:lowercontrol_refined} requires, at the level of exponents,
\[
3\beta-2\alpha\le \frac15,
\qquad
\beta-2\alpha\le -\frac95.
\]
Thus \(\Sigma\) may oscillate on short scales, but only with correspondingly small amplitude. For example,
\[
A_{\Sigma}=O(e^{-a}),
\qquad
L_{\Sigma}=O(e^{-a/10})
\]
is admissible. In this sense, the theorem allows \emph{high-frequency, small-amplitude} concentration for \(\Sigma\), but excludes a genuine short-pulse regime for \(\Sigma\) itself.

\noindent By contrast, \eqref{eq:ampl_refined} is markedly less restrictive. At the heuristic threshold allowed by \(\mathcal I^{0}\sim e^{a/10}\), one may take
\[
|A_{\mathfrak T}|^{2}=O(e^{a/10}),
\qquad
L_{\mathfrak T}=O(e^{-a/10}),
\]
for which the left-hand side of \eqref{eq:ampl_refined} is of order
\[
e^{a/5}+1+e^{-a/5}.
\]
This is compatible with the bound \((\mathcal I^{0})^{2}\sim e^{a/5}\). Heuristically, the tensor \(\mathfrak T\) may therefore exhibit a genuine short-pulse-type concentration: large amplitude supported at very small spatial scale.

\noindent This has a natural interpretation in terms of the free gravitational field. Relative to the future-directed unit normal \(T\), let
\[
E:=W(T,\cdot,T,\cdot),\qquad
B:={}^{*}W(T,\cdot,T,\cdot)
\]
denote the electric and magnetic parts of the Weyl tensor. The trace-free Gauss--Codazzi relations show schematically that
\[
E=\mathfrak T+\mathrm{l.o.t.}(\Sigma),
\qquad
B=\operatorname{curl}\Sigma,
\]
where \(\mathrm{l.o.t.}(\Sigma)\) denotes expressions of lower differential order in \(\Sigma\) (depending on the precise normalization, these may include terms linear and quadratic in \(\Sigma\)). Since the lower-order norms of \(\Sigma\) are strongly suppressed, the electric part may inherit the short-pulse concentration carried by \(\mathfrak T\), while the magnetic part remains constrained by the smallness of \(\Sigma\) and its first derivative. Thus the class of data constructed here permits \emph{high-frequency, large-amplitude electric Weyl curvature}, but only \emph{high-frequency, small-amplitude magnetic Weyl curvature} on the initial slice. See Figure~\ref{fig:blue} for a schematic illustration, and Section~\ref{data} for the rigorous construction.

\noindent This asymmetry between the electric and magnetic components is specific to the present CMC/\(\Lambda>0\) framework and has no direct analogue in the standard \(\Lambda=0\) short-pulse constructions. It is therefore natural to ask whether a corresponding short-pulse concentration in the magnetic component could overcome the cosmological expansion and drive gravitational collapse. That question lies beyond the scope of the present paper. Our point here is only that the class of initial data considered in Theorem~\ref{main} is not a merely formal nonempty class: it already captures a genuinely nonlinear and physically meaningful concentration mechanism.

\subsection{Heuristics for the construction of the initial data}

Recall that in CMC gauge the re-scaled vacuum constraint equations reduce to
\begin{align}
\label{eq:cons1}
R(g)+\frac{n-1}{n} &= |\Sigma|_{g}^{2},\\
\label{eq:cons2}
\nabla^{j}\Sigma_{ij} &= 0,
\end{align}
where \(\Sigma\) is the trace-free part of the second fundamental form. Thus \(\Sigma\) is, by definition, trace-free, and \eqref{eq:cons2} is the momentum constraint.


\noindent For the class of data considered in Theorem~\ref{main}, one seeks
\[
\|\Sigma\|_{H^{2}}\lesssim e^{-a}\mathcal I^{0},
\qquad
\|\Sigma\|_{H^{3}}\lesssim \mathcal I^{0},
\qquad
\|\mathfrak T[g]\|_{H^{2}}\lesssim \mathcal I^{0},
\]
with
\[
\mathcal I^{0}e^{-a/10}< 1.
\]
At the heuristic upper threshold allowed by the theorem, one may think of
\[
\mathcal I^{0}\sim e^{a/10},
\qquad a\gg 1.
\]

\noindent The Hamiltonian constraint \eqref{eq:cons1} then forces the renormalized scalar-curvature defect to be much smaller than the \(H^{2}\)-size of \(\mathfrak T[g]\). Indeed, since \(H^{2}\) is an algebra in dimension three,
\[
\Bigl\|R(g)+\frac{n-1}{n}\Bigr\|_{H^{2}}
=
\||\Sigma|_{g}^{2}\|_{H^{2}}
\lesssim
\|\Sigma\|_{H^{2}}^{2}
\lesssim
e^{-2a}(\mathcal I^{0})^{2},
\]
while at one higher derivative one only obtains
\[
\Bigl\|R(g)+\frac{n-1}{n}\Bigr\|_{H^{3}}
=
\||\Sigma|_{g}^{2}\|_{H^{3}}
\lesssim
\|\Sigma\|_{H^{2}}\|\Sigma\|_{H^{3}}+\|\Sigma\|_{H^{3}}^{2}
\lesssim
(\mathcal I^{0})^{2}.
\]
If \(\mathcal I^{0}\sim e^{a/10}\), these bounds become
\[
\Bigl\|R(g)+\frac{n-1}{n}\Bigr\|_{H^{2}}
= O(e^{-9a/5}),
\qquad
\Bigl\|R(g)+\frac{n-1}{n}\Bigr\|_{H^{3}}
= O(e^{a/5}).
\]

\noindent Accordingly, the geometric core of the construction is the following. One must produce a metric \(g\) on a closed three-manifold of negative Yamabe type such that
\[
\|\mathfrak T[g]\|_{H^{2}} = O(e^{a/10}),
\qquad
\Bigl\|R(g)+\frac{n-1}{n}\Bigr\|_{H^{2}} = O(e^{-9a/5}),
\qquad
\Bigl\|R(g)+\frac{n-1}{n}\Bigr\|_{H^{3}} = O(e^{a/5}),
\]
while at the same time arranging for a trace-free tensor \(\Sigma\) satisfying
\[
\|\Sigma\|_{H^{2}} = O(e^{-9a/10}),
\qquad
\|\Sigma\|_{H^{3}} = O(e^{a/10}),
\]
and the momentum constraint.

\noindent The standard conformal method provides the natural framework for solving
\eqref{eq:cons1}--\eqref{eq:cons2}, but in the present setting the delicate point is not merely the solvability of the constraints. Rather, one must solve them in such a way that the conformal correction needed to enforce \eqref{eq:cons1} remains sufficiently small at low order so that the large \(H^{2}\)-size of the renormalized trace-free Ricci tensor is not destroyed.

\noindent The construction therefore proceeds heuristically in three steps. First, one builds a seed metric \(\widehat g\) for which the renormalized trace-free Ricci tensor is large in \(H^{2}\), whereas the scalar-curvature defect \(R(\widehat g)+\frac{n-1}{n}\) is very small in \(H^{2}\) and controlled in \(H^{3}\). Second, one constructs a \(\widehat g\)-transverse-traceless tensor \(\widehat\Sigma\) with the quantitative bounds required in Theorem~\ref{main}; in particular, this resolves the momentum constraint at the seed level. Third, one uses the conformal method to solve the Hamiltonian constraint exactly, producing a nearby pair \((g,\Sigma)\) satisfying the full system \eqref{eq:cons1}-\eqref{eq:cons2}. One then shows that the conformal correction is sufficiently small in the relevant norms, so that the large \(H^{2}\)-size of the renormalized trace-free Ricci tensor is retained.



\noindent This is the geometric mechanism underlying the initial-data construction. The following discussion here is only heuristic; the rigorous implementation is carried out in Section~\ref{data}.


\noindent Let $M$ be a closed connected $3$-manifold of negative Yamabe type. By definition, every conformal class on $M$ contains a smooth metric of constant negative scalar curvature. We therefore fix a smooth background metric $\gamma$ satisfying
\[
R(\gamma)=-\frac{n-1}{n}.
\]
Since throughout this discussion $n=3$, one may equivalently write $R(\gamma)=-\frac23$; for notational clarity, however, we retain the general $n$-dimensional form in the formulas below.

\noindent We begin with a smooth background metric $\gamma$ satisfying
\[
R[\gamma]=-\frac{n-1}{n}.
\]
We then consider a perturbation
\[
g_{1}=\gamma+\epsilon h,\qquad 0<\epsilon\ll 1,
\]
where $h$ is chosen to be transverse-traceless with respect to $\gamma$,
\[
\div_{\gamma}h=0,
\qquad
\tr_{\gamma}h=0,
\]
and normalized by
\[
\|h\|_{L^{2}(M,\gamma)}=1.
\]
The essential point is that $h$ is taken to be highly oscillatory, with characteristic frequency $\mu\gg1$. Concretely, one may choose $h$ from a high-frequency family of transverse-traceless tensors associated with any self-adjoint elliptic operator on the transverse-traceless bundle whose principal symbol is that of the rough Laplacian. Standard elliptic theory then yields the scale of estimates
\[
\|h\|_{H^{2}(M,\gamma)}\sim \mu,
\]
with constants depending only on the background geometry. If the background happens to be Einstein, one may take $h$ to be a high-frequency transverse-traceless eigentensor of the Lichnerowicz Laplacian; however, that special identification is not needed for the argument below.

\noindent We write
\[
\mathfrak{T}[g]:=\text{Ric}[g]-\frac1n R[g]\,g
\]
for the trace-free Ricci tensor. The linearization of $\mathfrak{T}$ at $\gamma$, restricted to transverse-traceless directions, is a second-order elliptic operator; we denote it schematically by $\mathcal{A}_{\gamma}$. Its principal part is the Laplace-type term, and therefore it carries the same high-frequency scaling as a second-order elliptic operator. Accordingly, one has the schematic expansion
\[
\mathfrak{T}[g_{1}]
=
\mathfrak{T}[\gamma]
+\epsilon\,\mathcal{A}_{\gamma}h
+\epsilon^{2}\mathcal{Q}_{\gamma}(h,\nabla h,\nabla^{2}h),
\]
where $\mathcal{Q}_{\gamma}$ is quadratic in $h$ and its derivatives. At the heuristic level relevant for the introduction, the ellipticity of $\mathcal{A}_{\gamma}$ implies
\[
\|\mathcal{A}_{\gamma}h\|_{H^{2}(M,\gamma)}\sim \mu^{2},
\]
whereas the quadratic remainder obeys the scale
\[
\|\mathcal{Q}_{\gamma}(h,\nabla h,\nabla^{2}h)\|_{H^{2}(M,\gamma)}\lesssim \mu^{3}.
\]
Thus
\[
\|\mathfrak{T}[g_{1}]\|_{H^{2}}
\gtrsim
\epsilon\mu^{2}
-
\epsilon^{2}\mu^{3}
-
O(1).
\]
This already exhibits the mechanism we exploit: the linear contribution to the trace-free Ricci tensor gains two powers of frequency, while the nonlinear error is smaller provided $\epsilon\mu$ remains sufficiently small.


\noindent The scalar curvature behaves differently. Since $h$ is transverse-traceless, its first variation is only zeroth order:
\[
D R_{\gamma}[h]
=
-\langle \text{Ric}^{\circ}[\gamma],h\rangle_{\gamma},
\]
where $\text{Ric}^{\circ}[\gamma]=\mathfrak{T}[\gamma]$. Consequently, the scalar curvature expansion has the schematic form
\[
R[g_{1}]
=
R[\gamma]
+
\epsilon\,\langle \text{Ric}^{\circ}[\gamma],h\rangle_{\gamma}
+
\epsilon^{2}\mathcal{R}_{\gamma}(h,\nabla h,\nabla^{2}h),
\]
with $\mathcal{R}_{\gamma}$ quadratic. Since the linear term involves no second derivatives of $h$, it is smaller by one power of frequency than the linear term in the trace-free Ricci tensor. Correspondingly,
\[
\Big\|R[g_{1}]+\frac{n-1}{n}\Big\|_{H^{2}}
\lesssim
\epsilon\mu+\epsilon^{2}\mu^{3}.
\]
The two expansions therefore, separate the relevant scales:
\[
\|\mathfrak{T}[g_{1}]\|_{H^{2}}
\sim
\epsilon\mu^{2},
\qquad
\Big\|R[g_{1}]+\frac{n-1}{n}\Big\|_{H^{2}}
\sim
\epsilon\mu+\epsilon^{2}\mu^{3},
\]
up to lower-order terms. In particular, we seek parameters for which
\[
\epsilon\mu^{2}\gg1,
\qquad
\epsilon\mu\ll1,
\qquad
\epsilon^{2}\mu^{3}\ll1.
\]

\noindent For the later quantitative argument, it is convenient to enforce the stronger bounds
\[
\epsilon^{2}\mu^{3}<e^{-a/10},
\qquad
\epsilon\mu^{2}>e^{a/10},
\]
for some large parameter $a\gg1$. Solving these inequalities yields
\[
e^{a/20}\epsilon^{-1/2}<\mu<\epsilon^{-2/3}e^{-a/30}.
\]
Equivalently,
\[
\epsilon^{1/2}e^{a/20}<\epsilon\mu<e^{-a/30}\epsilon^{1/3}.
\]
This interval is nonempty provided
\[
\epsilon^{1/2}e^{a/20}<e^{-a/30}\epsilon^{1/3},
\]
that is,
\[
\epsilon^{1/6}<e^{-a/12},
\qquad\text{equivalently}\qquad
\epsilon<e^{-a/2}.
\]
Thus, once $\epsilon$ is chosen sufficiently small relative to $a$, one may select a transverse-traceless tensor of frequency $\mu$ in the admissible range
\begin{equation}\label{eq:accum_intro}
e^{a/20}\epsilon^{-1/2}<\mu<\epsilon^{-2/3}e^{-a/30}.
\end{equation}
At the heuristic level, this produces a metric $g_{1}$ for which the trace-free Ricci tensor is large in $H^{2}$, while the scalar-curvature defect remains perturbative. This, for the appropriate choice of $\epsilon<e^{-a/2}$, yields the bounds 
\begin{align*}
 ||\mathfrak{T}[g_{1}]||_{H^{2}}=O(e^{a/10}),~||R(g)+\frac{2}{3}||_{H^{2}}=O(e^{-9a/5}).   
\end{align*}
However, the task is not completed since we need to construct the physical initial data $(g_{0},k_{0})$ that verifies the constraint equations. However, note that the momentum constraint is relatively easy to address. On the other hand, the Hamiltonian constraint deals with the scalar curvature, which, in principle, can be modified almost independently of the trace-free Ricci curvature by means of a conformal transformation. We do this in the second step.  

\noindent The second step is to solve the momentum and Hamiltonian constraint through a conformal deformation. Starting from the metric $g=g_{1}$ constructed above, together with a symmetric transverse-traceless tensor $\Sigma$ relative to $g$, we seek a positive conformal factor $\varphi$ and define
\[
g_{0}=\varphi^{\frac{4}{n-2}}g,
\qquad
\Sigma_{0}=\varphi^{-2}\Sigma.
\]
Here $(g_{0},\Sigma_{0})$ would be the physical initial data set that we desire.
The conformal covariance of the momentum constraint ensures that $\Sigma_{0}$ remains transverse-traceless with respect to $g_{0}$:
\[
\div_{g_{0}}\Sigma_{0}=0
\]
and therefore $\Sigma$ is the free data. Therefore, one needs to prove the existence of TT tensor (with respect to $g$) on $(M,g)$ that verifies $||\Sigma||_{H^{2}}=O(e^{-9a/10})$ and $||\Sigma||_{H^{3}}=O(e^{a/10})$. Then, our remaining task is to solve for the Hamiltonian constraint through the conformal transformation given $\Sigma$ and prove that the resulting entities are not deformed substantially i.e., the estimates $||\mathfrak{T}[g_{0}]||_{H^{2}}=O(e^{a/10})$ and $||\Sigma||_{H^{2}}=O(e^{-9a/10}),~||\Sigma||_{H^{3}}=O(e^{a/10})$ survives in $g_{0}$ norm. For this, we need to prove that the conformal factor $\varphi$ is close to $1$ in $H^{4}$ in $(g-)$ norm and hence in $H^{4}$ with respect to $g_{0}-$ norm. 
The scalar curvature transforms according to
\begin{equation}\label{eq:scal_intro}
R[g_{0}]
=
\varphi^{-\frac{n+2}{n-2}}
\left(
-\frac{4(n-1)}{n-2}\Delta_{g}\varphi
+
R[g]\varphi
\right).
\end{equation}
Hence, the Hamiltonian constraint
\[
R[g_{0}]+\frac{n-1}{n}=|\Sigma_{0}|_{g_{0}}^{2}
\]
reduces to the usual Lichnerowicz equation for $\varphi$. For the purposes of the present discussion, the important point is simply that the conformal factor is chosen so as to remove the scalar-curvature defect left over from the first step.

\noindent What must still be verified is that this scalar-curvature correction does not significantly alter the large trace-free Ricci component already present in $g$. The transformation law for the trace-free Ricci tensor is
\[
\mathfrak{T}[g_{0}]
=
\mathfrak{T}[g]
-(n-2)\varphi^{-1}
\left(
\nabla^{2}\varphi-\frac1n\Delta_{g}\varphi\,g
\right)
+\frac{2(n-1)}{n-2}\varphi^{-2}
\left(
d\varphi\otimes d\varphi-\frac1n|\nabla\varphi|_{g}^{2}g
\right).
\]
The detailed estimates established later show that the correction term on the right-hand side is negligible in $H^{2}$ in light of $H^{4}$ norm of $\varphi$ being close to 1:
\[
\left\|
-(n-2)\varphi^{-1}
\left(
\nabla^{2}\varphi-\frac1n\Delta_{g}\varphi\,g
\right)
+\frac{2(n-1)}{n-2}\varphi^{-2}
\left(
d\varphi\otimes d\varphi-\frac1n|\nabla\varphi|_{g}^{2}g
\right)
\right\|_{H^{2}}
=
O(e^{-a}).
\]
Therefore
\[
\|\mathfrak{T}[g_{0}]-\mathfrak{T}[g]\|_{H^{2}}=O(e^{-a}),
\]
so the large $H^{2}$-size of the trace-free Ricci tensor survives the conformal correction. In particular,
\[
\|\mathfrak{T}[g_{0}]\|_{H^{2}}
\sim
\|\mathfrak{T}[g]\|_{H^{2}},
\]
up to an error which is negligible on the scale relevant to the construction. Therefore, we end up with the desired open set of initial data set $(g_{0},\Sigma_{0})$ that verifies 
\begin{align*}
R[g_{0,a}]+\frac23&=|\Sigma_{0,a}|_{g_{0,a}}^{2},\\
\div_{g_{0,a}}\Sigma_{0,a}&=0,\\
\tr_{g_{0,a}}\Sigma_{0,a}&=0,
\end{align*}
and 
\begin{align*}
\|\mathfrak T[g_{0,a}]\|_{H^{2}}&=O(e^{a/10}),~||\Sigma_{0}||_{H^{2}}=O(e^{-9a/10}),~||\Sigma_{0}||_{H^{3}}=O(e^{a/10}).   
\end{align*}


\noindent In summary, the initial data are produced by combining two mechanisms with complementary effects. A high-frequency transverse-traceless perturbation creates a large trace-free Ricci component while keeping the scalar curvature nearly constant, and a subsequent conformal deformation restores the Hamiltonian constraint while modifying the trace-free Ricci tensor only perturbatively. Simultaneously one explicitly constructs the $TT$ tensor $\Sigma$ verifying large $H^{3}$ norm while small $H^{2}$ norm (in appropriate scale). This is the geometric content behind the construction; the rigorous implementation is carried out in Section~\ref{data}.






\begin{center}
\begin{figure}
\begin{center}
\includegraphics[width=15cm,height=48cm,keepaspectratio,keepaspectratio]{1.png}
\end{center}
\begin{center}
\caption{A rough heuristic depicting the characterization of the smooth initial data considered in the article in case the spatial slice $M$ admits a hyperbolic metric. This clearly depicts the difference between the current work and previous work in this framework. The red center is the isolated hyperbolic fixed point (Mostow rigidity). The internal yellow ball $D_{\delta}(\gamma,0)$ depicts the small data regime studied in \cite{fajman2020stable,mondal2019attractors} where $d_{0}$ is the $H^{s}\times H^{s-1}$ Sobolev distance for $s> 3/2+1$. In contrast current study can handle data lying within the ball $B_{a}((\gamma,0),\mathcal{I}^{0})$ where the ball $B_{a}$ is with respect to the distance as defined in (\ref{eq:2}) of the main theorem. This ball can be made arbitrarily large by increasing the parameter $a$.}
\label{fig:blue}
\end{center}
\end{figure}
\end{center}



\noindent
A corollary to the theorem \ref{main} is as follows      

\begin{corollary}
\label{cor:cosmic_censorship}
Conjecture~\ref{conjecture1} holds in the class of $(3+1)$-dimensional solutions to the vacuum Einstein equations with positive cosmological constant $\Lambda > 0$, defined on globally hyperbolic spacetimes $\widetilde{M} \cong M \times \mathbb{R}$ where $M$ is a closed $3$-manifold of negative Yamabe type. The result holds for a class of large initial data prescribed on a Cauchy hypersurface at a Newtonian-like time $T_{0}=a$. In particular, no future naked singularities can form within this setting.
\end{corollary}


%Note that 
%\begin{theorem}
%\label{4}
 %   Let $(\widehat{M},\widehat{g})$ be a $3+1$ dimensional globally hyperbolic spacetime. Suppose the Cauchy slice $(M,g)$ is negative Yamabe type i.e., $\sigma(M)\leq 0$.  Let the initial data set $(g_{0},\Sigma_{0})$ for the reduced Einstein-$\Lambda$ equations \ref{eq:cmc1}-\ref{eq:cmc2} in CMC transported spatial gauge verify the constraint equations together with the estimates 
%\begin{align}
%\label{eq:1}
%\nonumber C^{-1}\xi_{0}\leq g_{0}\leq C\xi_{0},
%\end{align}
%\begin{align}
%\label{eq:2}
%\sum_{0\leq I\leq 4}||\nabla^{I}\Sigma_{0}||_{L^{2}(M)}\leq \epsilon _{0}\\
%\sum_{0\leq I\leq 3}||\nabla^{I}\mathfrak{T}[g_{0}]||_{L^{2}(M)}\leq \mathcal{I}^{0}
%\end{align}
%for a constant $C>0$, a smooth background Riemannian metric $\xi_{0}$, and a large $\mathcal{I}^{0}$. Then a unique solution to the Einstein-$\Lambda$ equation with the initial data $(M,g_{0},\Sigma_{0})$ at CMC \textit{Newtonian} time $T_{0}=0$ exists for all $T\in [0,\infty)$ , and the following estimates hold 
%\begin{align}
%\sum_{0\leq I\leq 4}||\nabla^{I}\Sigma(T)||_{L^{2}(M)}\lesssim \epsilon _{0},\\
%\sum_{0\leq I\leq 3}||\nabla^{I}\mathfrak{T}[g(T)]||_{L^{2}(M)}\lesssim \mathcal{I}^{0}\\
%\sum_{0\leq I\leq 3}||e^{T}\nabla^{I}\Sigma||_{L^{2}(M)}\lesssim \epsilon_{0},
%\end{align}
%\begin{align}
%\nonumber C^{-1}_{2} g_{0}\leq g(T)\leq C_{2} g_{0}~\forall T\in [0,\infty)
%\end{align}
%for numerical constants $C_{1},C_{2}$. 
%Moreover, every solution $T\mapsto (g(T),\Sigma(T))$ with initial data solving constraints \ref{eq:HC11}-\ref{eq:cmc2} and verifying the estimates \ref{eq:1}-\ref{eq:2} converges to $(\widetilde{g},0)$ for a metric $\widetilde{g}$ with point-wise constant negative scalar curvature as $T\to\infty$.
%\end{theorem}

\subsection{Comparison with previous work}

The mathematical analysis of Einstein's equations with positive cosmological constant has a substantial history, but the regimes treated in the literature differ markedly from the one considered here.

\noindent At the level of spatially homogeneous cosmologies, Wald's cosmic no-hair theorem shows that initially expanding homogeneous solutions with \(\Lambda>0\) are driven toward de Sitter behavior at late times \cite{wald}.  This result is foundational, but it concerns homogeneous dynamics and does not address the large-data Cauchy problem on a general compact spatial manifold.  In a different direction, Ringström established future nonlinear stability results for expanding solutions in related Einstein--matter models; in particular, the Einstein equations with positive cosmological constant arise there as a distinguished special case of the more general framework \cite{ringstrom2008future}.  These works provide essential background for the subject, but they do not furnish the type of compact, large-data, CMC-based theorem proved here.

\noindent For the vacuum Einstein--\(\Lambda\) equations, Friedrich introduced the conformal field equations and recast the problem as a regular first-order symmetric hyperbolic system for suitably rescaled variables \cite{friedrich1986existence1,friedrich1986existence}.  In that framework one prescribes asymptotic data at conformal infinity \(\mathscr J^\pm\), rather than physical Cauchy data on an interior hypersurface, and the conformal formulation removes the scalar constraint at the hypersurface where the conformal factor vanishes.  Friedrich's theory yields asymptotically simple solutions for data sufficiently close to the de Sitter configuration.  This is a profound and geometrically global approach, but it is conceptually and analytically different from the present one: our theorem is formulated directly on a physical compact Cauchy slice at a large CMC time \(T_{0}=a\), and the initial data satisfy the full nonlinear Einstein constraint equations on that slice.

\noindent More closely related to the present work are the CMC-based perturbative results of Fajman--Kr\"oncke and the attractor analysis of Mondal.  In the work of Fajman--Kr\"oncke, one studies the Einstein--\(\Lambda\) flow near background solutions for which the spatial metric is Einstein with positive or negative Einstein constant; the analysis is perturbative around such fixed points in high Sobolev regularity \cite{fajman2020stable}.  Likewise, in Mondal's work, one proves global existence and convergence for sufficiently small perturbations of a family of Einstein background solutions in CMC spatial harmonic gauge, with a shadow-gauge mechanism in higher dimensions to control the Einstein moduli directions \cite{mondal2019attractors}.  In both cases the initial data are assumed to lie in a genuinely small neighborhood of an Einstein background (or of the associated finite-dimensional Einstein moduli space).

\noindent The present theorem operates in a different regime.  First, no Einstein background is fixed in advance.  The spatial manifold \(M\) is only assumed to be a closed \(3\)-manifold with
\[
\sigma(M)\le 0,
\]
and the theorem does not require that \(M\) admit a negative Einstein metric.  Second, the result is formulated for physical Cauchy data \((g_{0},k_{0})\) posed on a compact slice at time \(T_{0}=a\), and the constraint equations are solved explicitly for the class of data under consideration.  Third, the theorem is genuinely non-perturbative from the viewpoint of the evolution variables: the size parameter \(\mathcal I^{0}\) is arbitrary, and the only smallness requirement is the relative condition
\[
\mathcal I^{0}e^{-a/10}<1.
\]
In particular, the renormalized trace-free Ricci tensor \(\mathfrak T[g_{0}]\) is not assumed to be small in \(H^{2}\), and the unweighted top-order norm of the trace-free second fundamental form \(\Sigma_{0}\) is not assumed to be small in \(H^{3}\).  This places the theorem outside the perturbative framework of \cite{fajman2020stable,mondal2019attractors}.

It is also important to distinguish the present large-data statement from a late-time reformulation of a small-data theorem.  The point is not merely that the initial slice is chosen far out in the expanding regime, but that the argument exploits a genuinely new structural feature of the Einstein--\(\Lambda\) equations in CMC-transported spatial coordinates: the nonlinear terms that are critical in the \(\Lambda=0\) theory acquire an additional time-integrable \(e^{-T}\)-weight.  This makes it possible to control explicitly constructed physical data for which certain natural geometric Sobolev norms are arbitrarily large, provided only that their largeness is dominated by the expansion scale at the initial slice.  The theorem therefore, lies beyond perturbative stability near de Sitter or near a fixed compact Einstein background.

\noindent To the best of our knowledge, this is the first future-global well-posedness and convergence theorem for the \(3+1\)-dimensional vacuum Einstein equations with \(\Lambda>0\) on compact CMC slices, based on an explicitly constructed class of physical Cauchy data, that is not perturbative around a prescribed Einstein metric.  In particular, it yields global future control and smooth convergence for a nontrivial large-data class on compact \(3\)-manifolds with \(\sigma(M)\le 0\).
 

\subsection{Novelty and the main difficulty }
A superficial analogy might suggest that a small--data global existence theorem on
\([0,\infty)\) should automatically yield a large--data result on \([T_{0},\infty)\), provided
the initial slice is placed sufficiently far in the expanding regime, that is, \(T_{0}\gg1\).
In the present problem, this transference principle is false.  The reason is that the
available perturbative results for the Einstein--\(\Lambda\) flow are formulated in a genuinely
small neighbourhood of an Einstein background, whereas Theorem~\ref{main} applies to
an explicitly constructed family of initial data for which the renormalized trace--free
Ricci tensor and the unweighted top--order norm of the trace--free second fundamental
form may be arbitrarily large.

\noindent To make this precise, recall first the perturbative hypotheses in the work of
Fajman--Kr\"oncke \cite{fajman2020stable}.  In the expanding negative Einstein case, the background metric
\(\gamma\) is assumed to satisfy
\[
\text{Ric}(\gamma)=-\frac{(n-1)}{n^{2}}\gamma,
\]
and the initial data are required to obey a smallness condition of the form
\begin{equation}
\label{eq:FK-small}
\|g_{0}-\gamma\|_{H^{s'}}+\|k_{0}+\sqrt{2}\,\gamma\|_{H^{s'-1}}<\delta,
\end{equation}
for suitable Sobolev exponents \(s'>\frac n2+1\), with \(\delta>0\)
sufficiently small.  Likewise, in the perturbative Einstein--\(\Lambda\) attractor theory of
\cite{mondal2019attractors}, one assumes that the rescaled data lie in a small ball
\begin{equation}
\label{eq:Mondal-small}
(g(T_{0}),K^{TT}(T_{0}))\in B_{\delta}(\gamma_{0},0)\subset H^{s}\times H^{s-1},
\end{equation}
centered at an Einstein background \(\gamma_{0}\) (more precisely, at the associated
integrable deformation space), together with the shadow gauge condition.  In both
cases, the theory is perturbative in the strongest possible sense: the initial metric is
assumed to be close, in a high Sobolev topology, to a negative Einstein metric, and the
transverse--traceless part of the second fundamental form is required to be small in the
corresponding Sobolev norm.

\noindent The hypothesis of Theorem~\ref{main} is of a completely different character.  One assumes
only the uniform metric comparability
\[
C^{-1}\xi_{0}\leq g_{0}\leq C\xi_{0},
\]
together with the bound
\begin{equation}
\label{eq:main-large}
\sum_{I=0}^{3}\|\nabla^{I}\Sigma_{0}\|_{L^{2}(M)}
+\sum_{I=0}^{2}\Bigl(\|\nabla^{I}\mathfrak T[g_{0}]\|_{L^{2}(M)}
+\|e^{a}\nabla^{I}\Sigma_{0}\|_{L^{2}(M)}\Bigr)
\leq \mathcal I^{0},
\end{equation}
where \(\mathcal I^{0}>0\) is arbitrary and the only smallness requirement is the relative
condition
\[
\mathcal I^{0}e^{-a/10}<1.
\]
In particular, neither \(\|\mathfrak T[g_{0}]\|_{H^{2}}\) nor the unweighted quantity
\(\|\Sigma_{0}\|_{H^{3}}\) is required to be small.  Thus the present theorem allows a regime
in which
\[
\|\mathfrak T[g_{0}]\|_{H^{2}}\sim \mathcal I^{0},
\qquad
\|\Sigma_{0}\|_{H^{3}}\sim \mathcal I^{0},
\]
with \(\mathcal I^{0}\) as large as one pleases, provided that the initial CMC time \(a\) is
chosen sufficiently large in terms of \(\mathcal I^{0}\).

\noindent This difference is not merely linguistic; it is quantitative.  Indeed, if \(\gamma\) is an
Einstein metric, then \(\mathfrak T[\gamma]=0\), and by continuity of the nonlinear map
\(g\mapsto \mathfrak T[g]\) from high Sobolev metrics to curvature in a two--derivative lower
Sobolev space, any perturbative hypothesis of the form \eqref{eq:FK-small} forces
\(\|\mathfrak T[g_{0}]\|_{H^{s'-2}}\lesssim \delta\).  Hence an explicitly constructed family with
\(\|\mathfrak T[g_{0}]\|_{H^{2}}\) arbitrarily large eventually lies outside every perturbative
neighbourhood covered by the small--data theories.  The same is true for the momentum
variable: the smallness assumptions in \eqref{eq:FK-small} or \eqref{eq:Mondal-small}
are incompatible with an unweighted \(H^{3}\)-norm of \(\Sigma_{0}\) of size \(\mathcal I^{0}\gg1\).

\noindent A second distinction is geometric.  The perturbative works start from a negative Einstein
background and therefore presuppose the existence of such a metric on the spatial
manifold.  Theorem~\ref{main}, by contrast, is stated on an arbitrary closed
\(3\)-manifold of negative Yamabe type, under the geometric bounds
\eqref{eq:1}--\eqref{eq:2}, without assuming the existence of any Einstein metric and
without placing the data in a small neighbourhood of any center manifold or moduli
space.  The result therefore applies to explicitly constructed large data on compact
topologies that need not arise as small perturbations of a hyperbolic Einstein geometry.

\noindent It is also important to distinguish the present large--data regime from mere smallness of
the scalar curvature defect.  In dimension \(3\), the Hamiltonian constraint gives
\[
R(g_{0})+\frac{2}{3}=|\Sigma_{0}|^{2}.
\]
Since the \(H^{2}\)-norm of \(\Sigma_{0}\) is weighted by \(e^{a}\), one obtains
\[
\|R(g_{0})+\tfrac23\|_{H^{2}}
\lesssim
\|\Sigma_{0}\|_{H^{2}}^{2}
\lesssim
e^{-2a}(\mathcal I^{0})^{2}.
\]
\noindent Thus the scalar curvature defect may be very small even when
\(\|\mathfrak T[g_{0}]\|_{H^{2}}\) and \(\|\Sigma_{0}\|_{H^{3}}\) are very large.  This does not
place \(g_{0}\) in a perturbative neighborhood of an Einstein metric, because the scalar
curvature controls only the trace part of the Ricci tensor and does not control the full
trace-free Ricci tensor in any comparable topology.  The data considered here are
therefore genuinely large from the point of view of the full Einstein evolution.

\noindent In summary, the present theorem is not a rephrasing of a known small-data stability
result at late time.  It gives global future control for an explicitly constructed family of
initial data satisfying only the relative smallness condition
\(\mathcal I^{0}e^{-a/10}<1\), while permitting arbitrarily large
\(\|\mathfrak T[g_{0}]\|_{H^{2}}\) and arbitrarily large unweighted
\(\|\Sigma_{0}\|_{H^{3}}\).  To the best of our knowledge, this is the first rigorous global
large--data convergence theorem for the vacuum Einstein equations with
\(\Lambda>0\) in CMC gauge on compact spatial manifolds of negative Yamabe type that is
not perturbative around a fixed Einstein background.
  


\noindent We next record several additional features of the formulation and of the proof that are essential in the large--data regime.

\medskip

\noindent A first point concerns the choice of gauge.  Rather than working with the coupled hyperbolic--elliptic Einstein system in constant mean curvature and spatial harmonic gauge, we work in CMC--transported spatial coordinates, so that the shift vector field vanishes identically, and we formulate the evolution in terms of the pair
\[
(\Sigma,\mathfrak T),\qquad
\mathfrak T_{ij}:=\text{Ric}_{ij}[g]-\frac13 R[g]\,g_{ij},
\]
coupled only to the elliptic lapse equation \eqref{eq:lapse3}.  This point of view is closer in spirit to the classical work of Choquet--Bruhat and York, where one seeks to evolve curvature-type variables together with the second fundamental form.  In the present gauge, the metric is recovered from the transport equation \eqref{eq:metric_transport} along the future-directed normal flow, whereas the higher-order analysis is carried by the evolution system satisfied by \(\Sigma\) and \(\mathfrak T\).

\medskip

\noindent We now make precise why the spatial harmonic gauge is well suited to perturbative
theory, but not to the genuinely large--data setting considered here. Let \(\xi\) be a fixed smooth background metric on \(M\), and let
\[
V^{i}:=g^{jk}\bigl(\Gamma[g]^{i}_{jk}-\Gamma[\xi]^{i}_{jk}\bigr)
\]
denote the spatial harmonic gauge vector. Preserving the condition \(V=0\) under the
Einstein evolution leads to an elliptic equation for the shift vector field \(X\) of the form
\begin{equation}
\label{eq:shift_operator_intro}
\mathcal P_{g,\xi}X=\mathcal F[g,\Sigma,N],
\end{equation}
where, schematically,
\begin{equation}
\label{eq:P_g_xi_intro}
(\mathcal P_{g,\xi}X)^{i}
=
-\Delta_{g}X^{i}
-\text{Ric}[g]^{i}{}_{j}X^{j}
+
2\nabla^{j}X^{k}\bigl(\Gamma[g]^{i}_{jk}-\Gamma[\xi]^{i}_{jk}\bigr),
\end{equation}
and \(\mathcal F[g,\Sigma,N]\) depends on the lapse, the second fundamental form, and the
metric.  This is the operator that must be inverted in order to solve for the shift.

\noindent The perturbative theory rests on the fact that when \(g\) remains close to a fixed negative
Einstein metric \(\xi\), the operator \(\mathcal P_{g,\xi}\) is a small perturbation of the
background operator
\begin{equation}
\label{eq:P_background_intro}
\mathcal P_{\xi,\xi}X
=
-\Delta_{\xi}X-\text{Ric}[\xi](X).
\end{equation}
If \(\xi\) is Einstein with
\[
\text{Ric}[\xi]=-\frac{(n-1)}{n^{2}}\xi,
\]
then
\begin{equation}
\label{eq:P_background_hyperbolic_intro}
\mathcal P_{\xi,\xi}X
=
-\Delta_{\xi}X+\frac{(n-1)}{n^{2}}X.
\end{equation}
In particular, for every smooth vector field \(X\),
\begin{align}
\label{eq:background_coercive_intro}
\int_{M}\langle \mathcal P_{\xi,\xi}X,X\rangle_{\xi}\,d\mu_{\xi}
&=
\int_{M}\Bigl(|\nabla^{\xi}X|_{\xi}^{2}+\frac{(n-1)}{n^{2}}|X|_{\xi}^{2}\Bigr)\,d\mu_{\xi} \\
&\ge c_{0}\|X\|_{H^{1}(\xi)}^{2}, \nonumber
\end{align}
for some constant \(c_{0}>0\) depending only on \((M,\xi)\).  Thus
\(\mathcal P_{\xi,\xi}\) is injective, self-adjoint on \(L^{2}(TM,\xi)\), and, by standard
elliptic Fredholm theory, an isomorphism
\[
\mathcal P_{\xi,\xi}\colon H^{r+1}(TM,\xi)\longrightarrow H^{r-1}(TM,\xi)
\]
for every \(r\ge1\).  Moreover,
\begin{equation}
\label{eq:background_elliptic_estimate_intro}
\|X\|_{H^{r+1}(\xi)}
\le C_{r,\xi}\|\mathcal P_{\xi,\xi}X\|_{H^{r-1}(\xi)}.
\end{equation}

\noindent The small-data argument is then obtained by perturbing away from
\(\mathcal P_{\xi,\xi}\).  Fix \(s>\frac n2+1\).  If \(g\) satisfies
\begin{equation}
\label{eq:g_close_xi_intro}
\|g-\xi\|_{H^{s}(\xi)}\le \delta,
\end{equation}
with \(\delta\) sufficiently small, then Sobolev multiplication and the smooth dependence of
\(\Gamma[g]\), \(g^{-1}\), and \(\text{Ric}[g]\) on \(g\) imply the operator bound
\begin{equation}
\label{eq:operator_perturbation_intro}
\|(\mathcal P_{g,\xi}-\mathcal P_{\xi,\xi})X\|_{H^{r-1}(\xi)}
\le
C_{r,\xi}\,\|g-\xi\|_{H^{s}(\xi)}\,\|X\|_{H^{r+1}(\xi)},
\qquad 1\le r\le s.
\end{equation}
In particular, taking \(r=1\) and pairing against \(X\), one obtains
\begin{equation}
\label{eq:quadratic_perturbation_intro}
\biggl|
\int_{M}
\bigl\langle
(\mathcal P_{g,\xi}-\mathcal P_{\xi,\xi})X,X
\bigr\rangle_{\xi}\,d\mu_{\xi}
\biggr|
\le
C_{\xi}\,\|g-\xi\|_{H^{s}(\xi)}\,\|X\|_{H^{1}(\xi)}^{2}.
\end{equation}
Combining \eqref{eq:background_coercive_intro} and
\eqref{eq:quadratic_perturbation_intro}, we find that if
\[
C_{\xi}\delta\le \frac12 c_{0},
\]
then
\begin{equation}
\label{eq:P_g_xi_coercive_intro}
\int_{M}\langle \mathcal P_{g,\xi}X,X\rangle_{\xi}\,d\mu_{\xi}
\ge \frac12 c_{0}\|X\|_{H^{1}(\xi)}^{2}.
\end{equation}
Hence \(\mathcal P_{g,\xi}\) is injective. Since it is elliptic of index zero and depends
continuously on \(g\), it follows that
\[
\mathcal P_{g,\xi}\colon H^{r+1}(TM,\xi)\to H^{r-1}(TM,\xi)
\]
is an isomorphism for all \(1\le r\le s\), with quantitative estimate
\begin{equation}
\label{eq:shift_estimate_intro}
\|X\|_{H^{r+1}(\xi)}
\le
C_{r,\xi}\|\mathcal P_{g,\xi}X\|_{H^{r-1}(\xi)}.
\end{equation}
This is the analytic mechanism used in perturbative treatments such as
Fajman-Kr\"oncke \cite{fajman2020stable} and Mondal \cite{mondal2019attractors}: the shift equation is solvable because the evolving metric
remains inside a sufficiently small \(H^{s}\)-neighbourhood of a fixed negative Einstein
background.

\noindent The large-data setting of Theorem~\ref{main} lies outside this framework.  First, the
theorem does not assume that \(M\) carries any fixed negative Einstein metric at all; it
assumes only that \(M\) is of negative Yamabe type.  Second, even when such a metric
exists, the initial data constructed in the theorem are not required to satisfy any
smallness condition of the form \eqref{eq:g_close_xi_intro}.  Thus one has no a priori
estimate forcing
\[
\|g(T)-\xi\|_{H^{s}(\xi)}\ll1
\]
on the bootstrap interval, and therefore no quantitative perturbative reduction of
\(\mathcal P_{g,\xi}\) to the coercive background operator \(\mathcal P_{\xi,\xi}\).  In
particular, the first--order coefficient
\[
2\nabla^{j}X^{k}\bigl(\Gamma[g]^{i}_{jk}-\Gamma[\xi]^{i}_{jk}\bigr)
\]
and the zeroth--order Ricci term \(-\text{Ric}[g](X)\) can no longer be treated as controlled
perturbative errors relative to the background geometry.  For this reason, the spatial
harmonic gauge does not furnish a robust large--data gauge-fixing mechanism for the
problem studied here.

\noindent This is precisely why we work instead in CMC--transported spatial coordinates.  In that
gauge the shift vanishes identically, and the analysis is reduced to a hyperbolic evolution
system for \((\Sigma,\mathfrak T)\) coupled to a single elliptic equation for the lapse,
whose coefficients are handled directly in terms of the evolving geometry.
\medskip

\noindent A second point is the hierarchy of energy norms used to control the second fundamental form.  The top-order energy for \(\Sigma\) is uniformly bounded but, in general, does not decay.  By contrast, suitably renormalized lower-order norms of \(\Sigma\) do decay exponentially.  The proof exploits this discrepancy systematically: lower-order decay is first established and then fed back into the top-order energy estimates to control the nonlinear error terms that are borderline at the highest derivative level.  This hierarchy is one of the basic mechanisms by which the large-data bootstrap is closed.

\noindent Closely related to this is the structure of the nonlinear terms.  At the op order, the potentially dangerous interactions are accompanied by an additional factor \(e^{-T}\), and are therefore time-integrable.  This feature is not a classical null structure; rather, it is an integrable damping mechanism generated by the cosmological expansion.  The same phenomenon governs the lapse: although the normalized lapse defect \(\frac{N}{n}-1\) need not itself exhibit top-order decay, the terms in which its highest derivatives occur are always accompanied by an \(e^{-T}\)-weight, whereas the undamped appearances of \(\frac{N}{n}-1\) arise only at lower differential order.  This separation between top-order weighted terms and lower-order unweighted terms is crucial throughout the bootstrap argument; see Section~\ref{mainidea}.

\medskip

\noindent A further issue is that all Sobolev and elliptic norms are defined relative to the evolving metric \(g\).  Consequently, one must propagate quantitative control of the geometry of the slices simultaneously with the energy estimates.  The metric is governed by the transport equation
\begin{equation}
\label{eq:metric_transport}
\partial_T g_{ij}
=
-\frac{2\varphi}{\tau}N\Sigma_{ij}
-
2\Bigl(1-\frac{N}{n}\Bigr)g_{ij}.
\end{equation}
Integrating this identity along the CMC foliation yields control of the metric coefficients and uniform equivalence of the metrics \(g(T)\), provided \(\Sigma\) and \(\frac{N}{n}-1\) are already controlled in the relevant Sobolev norms.  What it does \emph{not} yield is a gain of one additional spatial derivative on \(g\).  For that reason, one cannot simply freeze the elliptic theory relative to a fixed background metric and appeal to standard coefficient-based regularity estimates at the top order.

\noindent The lapse equation illustrates this point.  In our normalization it takes the form
\begin{equation}
\label{eq:lapse_elliptic}
-\Delta_{g}\Bigl(\frac{N}{n}-1\Bigr)
+
\Bigl(|\Sigma|^{2}+\frac13\Bigr)\Bigl(\frac{N}{n}-1\Bigr)
=
-|\Sigma|^{2}.
\end{equation}
If one were to formulate all Sobolev norms relative to a fixed metric \(\xi_{0}\), then the usual elliptic theory would require coefficient control at a derivative level that is not directly available from the transport equation for \(g\).  The resolution is to commute \eqref{eq:lapse_elliptic} with covariant derivatives \(\nabla=\nabla[g]\) of the evolving metric itself.  In that formulation, the commutators involve curvature terms rather than uncontrolled higher derivatives of the metric coefficients.

\noindent This is precisely where the choice of unknowns becomes decisive.  Since the spatial dimension is three, the full Riemann tensor is algebraically determined by the Ricci tensor.  Writing
\[
\text{Ric}[g]
=
\mathfrak T+\frac13 R[g]\,g
\]
and using the Hamiltonian constraint
\begin{equation}
\label{eq:ham_constraint}
R[g]+\frac{2}{3}=|\Sigma|^{2},
\end{equation}
one expresses every curvature term appearing in the commuted lapse equation in terms of \(\mathfrak T\), \(\Sigma\), and lower-order contractions thereof.  Consequently, the curvature contributions in the elliptic commutators are controlled by the same Sobolev quantities that already enter the hyperbolic part of the argument.  This avoids derivative loss and yields the required estimates for the normalized lapse in the same regularity scale as the evolution variables.

\medskip

\noindent Another notable feature of the theorem is the generality of the allowed spatial topology.  The argument applies on every smooth, closed, connected, oriented \(3\)-manifold \(M\) with non-positive Yamabe invariant \(\sigma(M)\le 0\).  Combined with the explicit construction of the initial data class described earlier in the introduction, this places the result well beyond perturbative stability theory near a fixed hyperbolic background.  In particular, the theorem is not tied to a distinguished Einstein metric, nor to a center-manifold description of the long-time dynamics.

\medskip

\noindent Finally, the analytic framework developed here appears robust enough to extend beyond the vacuum problem.  The combination of expansion-driven integrable damping, lower-order decay, and elliptic estimates formulated relative to the evolving metric should be adaptable to Einstein-\(\Lambda\) systems with matter, including, for instance, Einstein-\(\Lambda\)-Maxwell and Einstein-\(\Lambda\)-Euler models.  Whether one can obtain comparably sharp large-data future-global results in those settings remains an interesting open problem.

\subsection{Main idea of the proof}
\label{mainidea}

We now explain the mechanism underlying the proof and isolate the precise role of the positive cosmological constant $\Lambda$.  The discussion in this subsection is intentionally schematic: all commuted equations, elliptic estimates, coercive energies, and bootstrap improvements are stated and proved later in the paper.  The point here is to identify the structural feature that makes large--data forward control possible in the Einstein--\(\Lambda\) setting and, at the same time, to clarify why the corresponding argument is unavailable in the vacuum case \(\Lambda=0\).

\medskip

\noindent Let \((\widehat M,\widehat g)\) be a globally hyperbolic spacetime foliated by compact constant--mean--curvature slices \(M_T\) of negative Yamabe type.  We illustrate this in the Andersson--Moncrief CMC--spatial harmonic gauge.  Denote by \(g\) the induced metric on \(M_T\), by \(N\) the lapse, by \(X\) the shift, and by \(\Sigma\) the trace--free second fundamental form.  In the vacuum case \(\Lambda=0\), the Einstein equations take the form
\begin{align}
\partial_T g_{ij}
 &= -2N\Sigma_{ij}
    -2\Bigl(1-\frac{N}{n}\Bigr)g_{ij}
    -(\mathscr L_X g)_{ij}, \label{eq:intro_metric_vac}\\
\partial_T \Sigma_{ij}
 &= -(n-1)\Sigma_{ij}
    -N\Bigl(\text{Ric}_{ij}-\frac{1}{3}R[g]g_{ij}\Bigr)
    +\nabla_i\nabla_j\Bigl(\frac{N}{n}-1\Bigr) \nonumber\\
 &\qquad
    +2N\Sigma_{ik}\Sigma^{k}{}_{j}
    -\frac1n\Bigl(\frac{N}{n}-1\Bigr)g_{ij}
    -(n-2)\Bigl(\frac{N}{n}-1\Bigr)\Sigma_{ij}
    -(\mathscr L_X\Sigma)_{ij}, \label{eq:intro_sigma_vac}
\end{align}
supplemented by the constraint equations
\begin{align}
R+\frac{n-1}{n}-|\Sigma|^2 &=0, \label{eq:intro_ham_vac}\\
\nabla_j\Sigma^j{}_i &=0. \label{eq:intro_mom_vac}
\end{align}
In the spatial harmonic gauge (together with the rescaled Hamiltonian constraint), the linearization of the map
\[
g\longmapsto \text{Ric}[g]-\frac{1}{3}R[g]g
\]
is elliptic, so that the gauge variables are determined by elliptic equations on each slice, whereas the pair \((g,\Sigma)\) evolves by a quasilinear hyperbolic system coupled to these elliptic constraints.

\medskip

\noindent The obstruction to large-data global control in \eqref{eq:intro_metric_vac}-\eqref{eq:intro_mom_vac} is already visible at the level of the \(\Sigma\)-equation.  The term \(-(n-1)\Sigma\) is linearly damping, but the remaining terms contain two genuinely nonperturbative mechanisms:
\begin{enumerate}
\item the curvature defect
\[
\mathfrak T_{ij}:=\text{Ric}_{ij}-\frac{1}{3}R[g]g_{ij},
\]
which measures the deviation from the constant negative sectional curvature geometry (homogeneous and isotropic) singled out by the gauge; and
\item the Riccati term \(2N\Sigma_i{}^k\Sigma_{kj}\), which is quadratic and of the same differential order as the damping term.
\end{enumerate}
Moreover, the lapse defect \(\frac{N}{n}-1\) is itself determined by an elliptic equation whose source is quadratic in the evolving geometric fields.  Thus, even at the level of the formal energy identities, the linearly decaying contribution \(-(n-1)\Sigma\) is coupled to nonlinear terms that are not accompanied by any time--integrable coefficient.  For this reason, every presently available forward stability argument in the \(\Lambda=0\) theory is fundamentally perturbative: one closes the estimates only after imposing smallness on the curvature defect and on the trace--free second fundamental form in scale--invariant Sobolev norms.  In particular, prescribing the data at a very large initial time \(T_0=a\gg1\) does not by itself create a new mechanism, since a bound of the form
\[
e^{\alpha T_0}\|(\Sigma,\mathfrak T)(T_0)\|\leq \mathcal I,
\qquad \alpha>0,
\]
is merely a reformulation of smallness of the unweighted fields at the initial slice.

\medskip

\noindent The situation changes decisively when \(\Lambda>0\).  After passing to the rescaled CMC variables used throughout the paper, every term in the evolution equations that is potentially dangerous from the point of view of long--time growth appears with an additional coefficient \(e^{-T}\).  More precisely, the commuted equations may be written schematically as
\begin{align}
\partial_T \Sigma
  &= -(n-1)\Sigma
     + e^{-T}\,\mathcal Q_1\bigl(g,N,\Sigma,\mathfrak T,\nabla N,\nabla^2N\bigr), \label{eq:intro_sigma_lambda_schematic}\\
\partial_T \mathfrak T
  &= e^{-T}\,\mathcal Q_2\bigl(g,N,\Sigma,\mathfrak T,\nabla\Sigma,\nabla^2\Sigma,\nabla N,\nabla^2N\bigr)
     + \Bigl(\frac{N}{n}-1\Bigr)g, \label{eq:intro_T_lambda_schematic}
\end{align}
where \(\mathcal Q_1\) and \(\mathcal Q_2\) denote universal linear combinations of tensorial contractions of the indicated quantities and their covariant derivatives.  In particular, the two terms that are most problematic in the vacuum case,
\[
\mathfrak T
\qquad\text{and}\qquad
\Sigma*\Sigma,
\]
now enter the evolution only through \(e^{-T}\mathfrak T\) and \(e^{-T}\Sigma*\Sigma\).  This is the basic structural fact on which the whole argument rests.

\medskip

\noindent We emphasize that this mechanism is not a null structure in the classical sense of Klainerman, since the gain does not arise from an algebraic cancellation in the quadratic form relative to the characteristic cone of the principal hyperbolic operator.  Rather, the positive cosmological constant produces a \emph{weak-null-type integrable damping mechanism}: the nonlinear interactions that would otherwise be borderline are multiplied by a universal coefficient \(e^{-T}\in L^1([T_0,\infty))\).  This integrability is exactly what makes large--data forward control possible once the initial CMC slice is placed sufficiently far in the expanding regime.

\medskip

\noindent To exploit this structure, we introduce three scale--invariant quantities.  The first is the high--order geometric energy
\[
\mathcal O(T)
 := \sum_{I\le 3}
 \Bigl(
   \|\nabla^I\Sigma\|_{L^2(M_T)}
   + \|\nabla^{I-1}\mathfrak T\|_{L^2(M_T)}
 \Bigr),
\]
which controls the differentiated curvature and deformation variables.  The second is the lower--order renormalized quantity
\[
\mathcal F(T)
 := \sum_{I\le 2}\|e^T\nabla^I\Sigma\|_{L^2(M_T)},
\]
equivalently
\[
E^{\mathrm{lower}}(T)
 := e^{-2T}\mathcal F(T)^2
 = \sum_{I\le 2}\|\nabla^I\Sigma\|_{L^2(M_T)}^2.
\]
The third is the pointwise control norm
\[
\mathcal N^\infty(T)
 := \|e^T\Sigma\|_{L^\infty(M_T)}
   + \Bigl\|e^{2T}\Bigl(\frac{N}{n}-1\Bigr)\Bigr\|_{L^\infty(M_T)}.
\]
The proof is organized around a bootstrap scheme for the triple
\[
\bigl(\mathcal O,\mathcal F,\mathcal N^\infty\bigr).
\]

\medskip

\noindent More precisely, on a bootstrap interval \([T_0,T^\ast)\) we assume
\[
\mathcal O(T)\le \Gamma,
\qquad
\mathcal F(T)\le \mathds Y,
\qquad
\mathcal N^\infty(T)\le \mathds L,
\qquad
T\in[T_0,T^\ast),
\]
with bootstrap constants chosen so that
\[
(\mathcal I^0)^2+\mathcal I^0+1<\min\{\Gamma,\mathds Y,\mathds L\},
\qquad
\Gamma+\mathds Y+\mathds L\le e^{T_0/5},
\]
where \(\mathcal I^0\) denotes the size of the initial data in the norms relevant to the main theorem.  The purpose of the argument is to show that, once \(T_0=a\) is taken sufficiently large such that \(\mathcal I^0e^{-a/10}<1\), one in fact has the stronger bounds
\[
\mathcal O(T)+\mathcal F(T)+\mathcal N^\infty(T)
 \lesssim \mathcal I^0+1
\qquad\text{for all }T\in[T_0,T^\ast),
\]
thereby improving the bootstrap assumptions and extending the solution globally to the future.

\medskip

\noindent The closure mechanism is hierarchical:
\[
\mathcal N^\infty
\;\Longrightarrow\;
\mathcal F
\;\Longrightarrow\;
\mathcal O.
\]
The first implication is obtained by combining the pointwise control of \(\Sigma\) and \(N\) with the lower--order evolution equation for \(\Sigma\).  The decisive term in the corresponding energy identity is the mixed contribution
\[
\sum_{I\le2} e^{-T}\int_{M_T} N\,\nabla^I\mathfrak T\,\nabla^I\Sigma\,\mu_g,
\]
which satisfies
\[
\biggl|
\sum_{I\le2} e^{-T}\int_{M_T} N\,\nabla^I\mathfrak T\,\nabla^I\Sigma\,\mu_g
\biggr|
\lesssim e^{-T}\mathcal O(T)^2.
\]
Accordingly, one obtains a differential inequality of the form
\[
\frac{d}{dT}E^{\mathrm{lower}}(T)
 \le -2(n-1)E^{\mathrm{lower}}(T)
      + C e^{-T}\mathcal O(T)^2.
\]
Since \(e^{-T}\) is integrable, this forcing term is genuinely lower order in time.  After integrating the inequality and using the bootstrap bound for \(\mathcal O\), one obtains first a uniform estimate for \(E^{\mathrm{lower}}\), and then, by a further iteration exploiting the decay already gained, the sharper bound
\[
\mathcal F(T)\lesssim 1+\mathcal I^0+\mathcal O(T)^2.
\]
The key point is that the curvature defect \(\mathfrak T\) does \emph{not} need to be small at the initial time; it only needs to be finite.  The factor \(e^{-T}\) is what converts large but finite geometric input into an integrable error.

\medskip

\noindent The second implication, from \(\mathcal F\) to \(\mathcal O\), enters at top order.  After commuting the equations and using the elliptic estimates for the lapse, one encounters cubic and quartic expressions in which one differentiated factor is paired with two lower order factors.  A representative contribution has the form
\[
\sum_{I\le2}\sum_{m=0}^I\sum_{J_1+J_2+J_3=m}
e^{-T}
\int_{M_T}
\nabla^{J_1}N\,
\nabla^{J_2+1}\Sigma\,
\nabla^{J_3+I-m}\Sigma\,
\nabla^{I+1}\Sigma\,
\mu_g.
\]
Using Sobolev, elliptic, and interpolation estimates, this is bounded by
\[
\lesssim e^{-2T}\mathcal F(T)\mathcal O(T)^2.
\]
Once the lower order estimate for \(\mathcal F\) is available, this becomes
\[
\lesssim e^{-2T}\mathcal O(T)^4
\lesssim e^{-2T}\Gamma^4.
\]
The coefficient \(e^{-2T}\) is now time--integrable with room to spare, and therefore
\[
\int_{T_0}^{T} e^{-2s}\Gamma^4\,ds
\lesssim e^{-2T_0}\Gamma^4.
\]
This is the point at which the ``relative smallness'' condition in the statement of the main theorem enters: for arbitrarily large but finite initial size \(\mathcal I^0\), one chooses \(T_0=a\) so large that the integrated top--order errors are perturbative.  In other words, the argument does not require absolute smallness of the data; it requires only that the largeness of the data be dominated by the expansion scale present at the initial slice.

\medskip

\noindent Finally, the metric and the gauge must be propagated without loss of geometric control.  The metric equation has the schematic form
\[
\partial_T g
 = e^{-T}N\Sigma
   -2\Bigl(1-\frac{N}{n}\Bigr)g,
\]
so once \(\mathcal N^\infty\) is controlled, the deformation of the metric is integrable in time.  This yields uniform equivalence of the evolving metrics, control of the volume form and isoperimetric constants, and hence uniform Sobolev inequalities on all future slices.  These geometric bounds are then fed back into the elliptic estimates for the lapse and the Sobolev estimates used in the energy argument, closing the system.

\medskip

\noindent We therefore obtain a self--consistent bootstrap mechanism in which lower order decay produces top order control, top order control feeds the elliptic theory, and the resulting gauge bounds preserve the geometric background needed for the Sobolev and energy estimates.  The entire argument hinges on the fact that, in the rescaled Einstein--\(\Lambda\) system, the nonlinear terms that are critical in the vacuum theory are accompanied by the integrable coefficient \(e^{-T}\).  This is the decisive large--data stabilizing effect of the positive cosmological constant and the basic reason the theorem proved here has no analogue in the known \(\Lambda=0\) CMC theory.








\begin{remark}
The analytical framework developed herein extends to the case in which the underlying manifold $M$ admits a positive Yamabe invariant, i.e., $\sigma(M)>0$, which includes in particular the setting with positive cosmological constant $\Lambda>0$. Under an additional spectral condition on the Laplace--Beltrami operator associated to the evolving metric, ensuring the uniqueness of the solution to the elliptic lapse equation, we establish convergence of the metric to one with constant positive scalar curvature.

\noindent This spectral condition, while technical, is natural in view of its role in guaranteeing elliptic solvability and uniqueness. In the absence of such a condition, uniqueness of the limiting metric generally fails. This phenomenon is closely related to the non-uniqueness observed in the convergence behavior of the Yamabe flow on manifolds with positive Yamabe type, as elucidated in the work of Brendle \cite{brendle}, where the flow is shown to converge to a constant scalar curvature metric, but the limit depends nontrivially on the initial data due to the lack of uniqueness in the conformal class.

\noindent In our context, the underlying cause of non-uniqueness is distinct and originates from the failure of uniqueness of solutions to the lapse equation, rather than the conformal degeneracy of the space of constant scalar curvature metrics. We emphasize that, even in the presence of convergence, the limiting metric is generally \emph{not} conformal to the initial metric. For completeness, and to illustrate the applicability of our method in this setting, we state the corresponding convergence result below; the proof proceeds with only minor modifications to the arguments presented in the negative Yamabe case. In addition, the construction of the data follows in an exact similar fashion.
\end{remark}
\noindent First, note that in case of $\sigma(M)>0$, the renormalized lapse equation takes the following form 
\begin{align}
\label{eq:positive}
-\Delta_{g}(1-\frac{N}{n})+(|\Sigma|^{2}-\frac{1}{3})(1-\frac{N}{n})=-|\Sigma|^{2},~\Delta_{g}:=g^{ij}\nabla_{i}\nabla_{j}.     
\end{align}
\noindent A subtle technical issue arises from the analysis of the elliptic equation \eqref{eq:positive} for the renormalized lapse $1 - \frac{N}{n}$. Namely, the associated operator
\[
-\Delta_{g} + n\left(|\Sigma|^{2} - \frac{1}{3}\right)
\]
may admit a nontrivial finite-dimensional kernel. To proceed, we impose the technical assumption that a solution to \eqref{eq:positive} exists. Note that such a solution need not be unique: even in the case of de Sitter spacetime, non-uniqueness of the lapse function is known for certain foliations.

\noindent Importantly, since the renormalized lapse function $1 - \frac{N}{n}$ appears as a source in the evolution equations for $(\Sigma, \mathfrak{T})$, this ambiguity propagates into the solution of the full reduced Einstein system. As a result, the limit geometry determined by the evolution is non-unique, reflecting the intrinsic degeneracy introduced by the kernel of the elliptic operator.


\begin{theorem}[Global well-posedness for $\Lambda > 0$ and $\sigma(M)>0$]
\label{main2}
Let $(\widehat{M}, \widehat{g})$ be a globally hyperbolic $(3+1)$-dimensional Lorentzian manifold, and let $(M,g)$ be a Cauchy hypersurface such that $M$ is a closed $3$-manifold of positive Yamabe type, i.e., 
\[
\sigma(M) > 0.
\] 
Fix $\mathcal{I}^0 > 0$. Then there exists a constant 
\[
a = a(\mathcal{I}^0) > 0
\]
sufficiently large such that 
\[
\mathcal{I}^0 e^{-\frac{a}{10}} < 1.
\]

\noindent Suppose that the initial data set 
\[
(g_0, \Sigma_0) \in C^\infty\big(\mathrm{Sym}^2(T^*M)\big) \times C^\infty\big(\mathrm{Sym}^2_0(T^*M)\big)
\]
for the reduced Einstein-$\Lambda$ system in constant mean curvature (CMC) transported spatial gauge satisfies the constraint equations and further verifies the bounds 
\begin{align}
\label{eq:metric_control}
C^{-1} \xi_0 \leq g_0 \leq C \xi_0,
\end{align}
and
\begin{align}
\label{eq:data_estimate}
\sum_{0 \leq I \leq 3} \|\nabla^I \Sigma_0\|_{L^2(M)} + \sum_{0 \leq I \leq 2} \|\nabla^I \mathfrak{T}[g_0]\|_{L^2(M)} + \sum_{0 \leq I \leq 2} \| e^{a} \nabla^I \Sigma_0 \|_{L^2(M)} \leq \mathcal{I}^0,
\end{align}
where $\xi_0$ is a fixed smooth background Riemannian metric on $M$ and $C > 0$ is a uniform constant.

\noindent Assume further that there exists a solution to the lapse equation corresponding to the positive Yamabe case (see Equation \eqref{eq:positive}).

\noindent Then the maximal developement 
\[
T \mapsto (g(T), \Sigma(T)) \in C^\infty\big(M; \mathrm{Sym}^2(T^*M)\big) \times C^\infty\big(M; \mathrm{Sym}^2_0(T^*M)\big),
\]
to the Einstein-$\Lambda$ system with initial data $(g_0, \Sigma_0)$ at CMC Newtonian time $T_0 = a$, exists globally for all
\[
T \in [T_0, \infty).
\]

\noindent Moreover, the solution exhibits the following asymptotic behavior as $T \to \infty$:
\begin{enumerate}[label=(\roman*)]
    \item The tensor $\Sigma(T)$ decays to zero in appropriate Sobolev norms;
    \item The metrics $g(T)$ converge smoothly to a limit metric $\widetilde{g}$ on $M$ with constant positive scalar curvature,
    \[
    R(\widetilde{g})(x) = \frac{n-1}{n}, \quad \forall x \in M.
    \]
\end{enumerate}

\noindent Finally, the limit metric $\widetilde{g}$ is non-unique precisely when the lapse equation admits non-unique solutions.

\end{theorem}



\begin{conjecture}
 Suppose Consider the following data-set for Einstein-$\Lambda$ vacuum system.    
\end{conjecture}

%\begin{remark}
 %Once again, the choice of $\gamma$ as stated in the theorem \ref{main2} is not optimal. The optimal value of $\gamma$ is $1$. This is discussed in section \ref{energy}   
%\end{remark}

\subsection{Consequences for Geometrization}

\noindent An important conceptual motivation underlying the present work concerns the potential role of the Einstein evolution equations in the geometric classification of closed $3$-manifolds. Specifically, we consider the relevance of the long-time behavior of solutions to the Einstein flow in connection with Thurston's geometrization conjecture. In this discussion, we assume throughout that the initial data satisfy the regularity assumptions stipulated in the main theorem.

\noindent A central open problem in mathematical general relativity is whether the Einstein equations on globally hyperbolic, spatially compact cosmological spacetimes can dynamically implement a form of geometrization of the underlying spatial manifold. In particular, we consider the case in which the Cauchy hypersurface \( M \) is a closed, oriented, connected $3$-manifold of negative Yamabe type, i.e., with \( \sigma(M) \leq 0 \).

\noindent The geometrization program in the setting of vacuum cosmological spacetimes was first proposed in the foundational works of Fischer--Moncrief \cite{fischer2000reduced,fischer2001reduced}, with further developments by Anderson \cite{anderson2001long}, who studied the asymptotic behavior of the vacuum Einstein flow and its possible connection to the Thurston decomposition. In the idealized scenario, one anticipates that the spatial metric induced on constant mean curvature (CMC) hypersurfaces asymptotically converges to a metric structure realizing a geometric decomposition of \( M \).

\noindent Following Anderson \cite{anderson2001long}, we recall the following definition of geometrization:

\begin{definition}[Geometrization]
Let \( M \) be a closed, oriented, connected $3$-manifold with \( \sigma(M) \leq 0 \). A \emph{weak geometrization} of \( M \) is a decomposition
\[
M = H \cup G,
\]
where:
\begin{itemize}
  \item \( H \) is a finite collection of embedded, complete, connected, finite-volume hyperbolic manifolds;
  \item \( G \) is a finite collection of embedded, connected graph manifolds;
  \item the union is along a finite collection of embedded tori \( \{T_i\} \), with \( \partial G = \partial H = \bigsqcup_i T_i \).
\end{itemize}
A \emph{strong geometrization} is a weak geometrization in which each gluing torus \( T_i \) is incompressible, i.e., the inclusion-induced homomorphism \( \pi_1(T_i) \to \pi_1(M) \) is injective.
\end{definition}

\noindent  Results due to Ringström \cite{ringstrom2008future} show that for a large class of initial data, including arbitrary spatial topologies, the Einstein-$\Lambda$ flow becomes essentially local in nature, and the large-scale topological structure of \( M \) becomes invisible to the dynamics. That is, the spatial topology is not imprinted on the asymptotic geometry of the evolving spacetime. We prove indeed this is the case for large initial data. 

\noindent A heuristic explanation for this phenomenon is furnished by comparison with a classical theorem of Cheeger and Gromov \cite{cheeger1,cheeger2}, which characterizes volume collapse of Riemannian manifolds with bounded curvature. Specifically, if a sequence \( \{(M_j, g_j)\} \) of compact Riemannian $3$-manifolds exhibits collapse with two-sided curvature bounds, then \( M_j \) admits an \( \mathcal{F} \)-structure of positive rank, and hence is a graph manifold for sufficiently large \( j \). Thus, any dynamical mechanism aiming to realize geometrization via Einstein evolution must, at minimum, produce volume collapse with bounded curvature in the graph component \( G \subset M \) for large times.

\noindent Our main theorem (Theorem~\ref{main}) establishes uniform bounds on the spacetime curvature under the Einstein–$\Lambda$ flow. However, we also prove that the volume of unit geodesic balls remains uniformly bounded from below for all future times, thereby precluding any form of collapse. More precisely, we obtain the following result:

\begin{corollary}[No-Collapse Theorem]
\label{nocollapse}
Let \( B(1) \subset M \) denote a geodesic ball of unit radius (scale $1$ in injectivity radius scale) with respect to the induced metric at CMC time \( T \in [T_0, \infty) \). Then the volume \( \operatorname{Vol}(B(1))_T \) remains uniformly bounded above by the initial volume \( \operatorname{Vol}(B(1))_{T_0} \) for all \( T \geq T_0 \).
\end{corollary}

\noindent The proof, given in Section~\ref{noncollapse}, follows directly from the estimates established in Theorem~\ref{main}. This result reflects the rapid expansion of the spacetime volume due to the presence of a positive cosmological constant, which dominates the large-scale dynamics in the asymptotic regime \( T \to \infty \). Consequently, the Einstein–$\Lambda$ flow does not differentiate between graph and hyperbolic components, and no meaningful geometric decomposition of \( M \) emerges in the limit.

\noindent The question of whether the Einstein equations with vanishing cosmological constant \( \Lambda = 0 \) can effectuate a geometric decomposition of \( M \) remains open and considerably more subtle. Unlike the Ricci flow, where Perelman's resolution of the geometrization conjecture \cite{perelman2002ricci,perelman2003ricci,perelman2003ricci2} crucially exploits the parabolic smoothing structure of the flow, the Einstein equations exhibit hyperbolic behavior, and their global-in-time analysis poses formidable technical challenges. Nonetheless, recent advances in the analysis of nonlinear wave equations in the large-data regime, such as \cite{C09}, offer promising avenues toward addressing these questions, and it is conceivable that similar techniques could be adapted to the cosmological setting.  

\noindent Ringström \cite{ringstrom2008future} provided a detailed local analysis of the Cauchy problem in general relativity with a positive cosmological constant and proposed a formal framework for quantifying the asymptotic inaccessibility of spatial topology by causal observers. His notion supports the view that long-time evolution under the Einstein-$\Lambda$ flow generically suppresses detectable topological information in the causal past of observers.

\begin{definition}
Let $(\widetilde{M}, \widehat{g})$ be a globally hyperbolic Lorentzian manifold with compact Cauchy hypersurfaces. We say that \emph{late-time observers are oblivious to topology} if there exists a Cauchy hypersurface $M \subset \widetilde{M}$ such that no future-directed inextendible causal curve $\gamma$ satisfies $M \subset J^{-}(\gamma)$. Conversely, we say that \emph{late-time observers are not oblivious to topology} if for every Cauchy hypersurface $M$, there exists a future-directed inextendible causal curve $\gamma$ with $M \subset J^{-}(\gamma)$.
\end{definition}

\noindent Based on this definition, Ringström proposed the following conjecture:

\begin{conjecture}[\cite{ringstrom2008future}]
Let $(\widetilde{M}, \widehat{g})$ be a future causally geodesically complete solution of the vacuum Einstein equations with positive cosmological constant $\Lambda > 0$, and compact Cauchy hypersurfaces. Then late-time observers in $(\widetilde{M}, \widehat{g})$ are oblivious to topology.
\end{conjecture}

\noindent We provide an affirmative to this conjecture in section \ref{mainproof}.  


%\subsection{Cosmological Implications}

%\noindent Although our analysis is grounded in the geometric analysis of the Einstein field equations with $\Lambda>0$, we briefly comment on an implication of our results for a longstanding open problem in mathematical cosmology. A central hypothesis in physical cosmology is the so-called \emph{cosmological principle}, which posits that, at sufficiently large spatial scales, the universe is approximately homogeneous and isotropic. This assumption underlies the derivation of the Friedmann–Lemaître–Robertson–Walker (FLRW) class of solutions, in which constant sectional curvature geometries model the spatial slices.

%\medskip

%\noindent From a rigorous mathematical perspective, a natural question is whether such large-scale homogeneity and isotropy can arise dynamically from generic initial data under the evolution dictated by the Einstein equations. In particular, one may ask whether spatial slices, even if initially far from being homogeneous or isotropic, can converge in an appropriate topology to a constant sectional curvature geometry—e.g., hyperbolic geometry—along the cosmological constant–driven CMC evolution. This idea was conjectured in \cite{MM}, where it was suggested that the long-time dynamics of the Einstein-$\Lambda$ flow may select the hyperbolic metric as a universal attractor.

%\medskip

%\noindent Our analysis, based on Theorem \ref{main}, provides a rigorous counterexample to this conjecture. Indeed, we work with the most general class of closed, connected, oriented three-manifolds of negative Yamabe type, encompassing all topological classes listed in Cases 1–3 of section 1. For such generic topologies, we demonstrate that no form of geometric convergence of the rescaled spatial metric to a hyperbolic (i.e., maximally symmetric) geometry occurs.

%\begin{corollary}[Failure of the Cosmological Principle as a Derived Consequence]
%Let $(\widetilde{M},\widehat{g})$ be a solution to the Einstein equations with positive cosmological constant $\Lambda>0$ in CMC-transported spatial gauge, constructed as in Theorem \ref{main}, and let $(M,g(T))$ denote the corresponding CMC foliation. Then, in $C^k$ or Sobolev, topology—the rescaled spatial metric $g(T)$ fails to converge, as $T \to \infty$, to a hyperbolic metric on $M$. In particular, the limiting spatial geometry is \emph{not} (even locally) homogeneous and isotropic.
%\end{corollary}

\noindent A proof of this statement is immediate from the structure of the limiting data described in Remark \ref{noconvergence}.
\medskip

\noindent While our analysis is confined to the vacuum Einstein equations with $\Lambda>0$, we conjecture that the failure of isotropization persists for a wide class of physically relevant matter models—including Einstein–$\Lambda$–Maxwell and Einstein–$\Lambda$–Euler systems. That is, we expect that generic initial data with nontrivial topology will not converge to spatial geometries exhibiting local homogeneity and isotropy, in the presence of matter fields, so long as the expansion is driven by a strictly positive cosmological constant.

 

\section*{Acknowledgement}
\noindent I thank Professor S-T Yau for many discussions regarding the scalar curvature geometry and Professor V. Moncrief for numerous discussions regarding Einstein flow. I thank the reviewer for the meticulous review, which substantially improved the quality of the article. This work was supported by the Center of Mathematical Sciences and Applications, Department of Mathematics at Harvard University and by BIMSA of YMSC at Tsinghua University.


\section{Preliminaries}

\noindent Throughout this work, we adopt the $(n+1)$-dimensional Arnowitt--Deser--Misner (ADM) formalism for the study of globally hyperbolic Lorentzian manifolds $(\widetilde{M},\widehat{g})$ with signature $(-,+,\ldots,+)$, where $n=3$. Specifically, we consider the canonical foliation 
\[
\widetilde{M} \cong \mathbb{R} \times M,
\]
where each leaf $\{t\} \times M$ is an orientable, compact $n$-dimensional Cauchy hypersurface diffeomorphic to $M$. The Lorentzian metric $\widehat{g}$ induces a Riemannian metric $g(t)$ on each slice $M(t) := \{t\}\times M$.

\noindent The ADM decomposition expresses the future-directed timelike unit normal vector field $\widehat{\mathbf{n}}$ orthogonal to the slices and the vector field $\partial_t$ generating the foliation as
\[
\partial_t = N \widehat{\mathbf{n}} + X,
\]
where $N: \widetilde{M} \to (0,\infty)$ is the lapse function, and $X \in \Gamma(T M)$ is the shift vector field, both of which possess the requisite regularity dictated by the function spaces to be specified later. The Lorentzian metric in the adapted coordinate system $\{x^\alpha\}_{\alpha=0}^n = \{t,x^i\}_{i=1}^n$ admits the canonical form
\begin{align}
\label{eq:spacetime_metric}
\widehat{g} = -N^2 dt \otimes dt + g_{ij} (dx^i + X^i dt) \otimes (dx^j + X^j dt),
\end{align}
where $g_{ij}(t,x)$ is the induced Riemannian metric on $M(t)$.

\noindent To quantify the extrinsic geometry of the embedding $M(t) \hookrightarrow \widetilde{M}$, we define the second fundamental form $k \in \Gamma(S^2 T^* M)$ by
\[
k_{ij} := -\langle \nabla_{\partial_i} \widehat{\mathbf{n}}, \partial_j \rangle = -\frac{1}{2N} \bigl(\partial_t g_{ij} - \mathscr{L}_X g_{ij}\bigr),
\]
where $\nabla$ denotes the Levi-Civita connection of $\widehat{g}$, and $\mathscr{L}_X$ is the Lie derivative along $X$. The mean extrinsic curvature, or trace of $k$, is given by
\[
\tau := \mathrm{tr}_g k = g^{ij} k_{ij},
\]
with $g^{ij}$ denoting the inverse metric on $M(t)$.

\noindent The vacuum Einstein field equations with cosmological constant $\Lambda \in \mathbb{R}$,
\begin{align}
\label{eq:einstein_equations}
\mathrm{Ric}(\widehat{g}) - \frac{1}{2} R(\widehat{g}) \widehat{g} + \Lambda \widehat{g} = 0,
\end{align}
admit an equivalent formulation in terms of the Cauchy data $(g,k)$ on each hypersurface $M(t)$ as the coupled evolution and constraint system:
\begin{align}
\label{eq:evol_g}
\partial_t g_{ij} &= -2N k_{ij} + \mathscr{L}_X g_{ij}, \\
\label{eq:evol_k}
\partial_t k_{ij} &= -\nabla_i \nabla_j N + N \left( R_{ij} + \tau k_{ij} - 2 k_i^m k_{mj} - \frac{2 \Lambda}{n-1} g_{ij} \right) + \mathscr{L}_X k_{ij}, \\
\label{eq:ham_constraint}
R - |k|_g^2 + \tau^2 &= 2\Lambda, \\
\label{eq:mom_constraint}
\nabla^i k_{ij} - \nabla_j \tau &= 0,
\end{align}
where $\nabla$ denotes the Levi-Civita connection of $g$, $R_{ij}$ and $R$ are the Ricci and scalar curvatures of $g$, respectively, and $|k|_g^2 = g^{im} g^{jn} k_{ij} k_{mn}$. Here and henceforth, indices are raised and lowered by $g$ and its inverse.

\noindent We denote by $\Gamma[g]$ the Christoffel symbols associated to $g$, and by $\mathrm{Riem}$ and $\mathrm{Ric}$ the corresponding Riemann and Ricci curvature tensors on $M$. The (negative-semidefinite) Laplace--Beltrami operator acting on functions $f \in C^\infty(M)$ is given by
\[
\Delta_g f := g^{ij} \nabla_i \nabla_j f,
\]
where the sign convention ensures that the spectrum of $-\Delta_g$ is nonnegative.

\noindent A central object in our analysis is the Lichnerowicz Laplacian
\[
\Delta_{L} : \Gamma(S^2 T^* M) \to \Gamma(S^2 T^* M),
\]
defined by
\[
(\Delta_{L} h)_{ij}
:=
g^{k\ell}\nabla_k\nabla_\ell h_{ij}
+2R_{ikj\ell}h^{k\ell}
-R_i{}^{k}h_{kj}
-R_j{}^{k}h_{ik}.,
\]
where the curvature operator acts via the Riemann curvature tensor. For metrics $g$ which are stable Einstein metrics, the operator $\mathcal{L}_g$ is elliptic and possesses strictly positive spectrum on the space of symmetric 2-tensors modulo gauge.

\noindent Finally, for nonnegative scalar functions $f,h : [0,\infty) \to \mathbb{R}_{\geq 0}$, we write $f(t) \lesssim h(t)$ to signify the existence of a constant $C > 0$, independent of $t$ but possibly dependent on the fixed background geometry and dimension $n$, such that
\[
f(t) \leq C h(t),
\]
and $f(t) \approx h(t)$ if there exist constants $0 < C_1 \leq C_2 < \infty$ such that
\[
C_1 h(t) \leq f(t) \leq C_2 h(t).
\]

\noindent Throughout, the spaces of smooth symmetric covariant 2-tensors and vector fields on $M$ are denoted by $S^0_2(M)$ and $\mathfrak{X}(M)$, respectively. The trace-free and transverse-traceless parts of a symmetric 2-tensor $A$ are denoted by $\widehat{A}$ and $A^{TT}$, respectively. We will also frequently denote the Cauchy slice at constant mean curvature time $T$ by $M(T)$. In addition, we frequently use the following transport identity, which is a consequence of the transport theorem \cite{marsden}
\begin{align}
 \frac{d}{dT}\int_{M(T)}f\mu_{g}=\int_{M(T)}(\partial_{T}f+\frac{1}{2}f\tr_{g}\partial_{T}g)\mu_{g}.   
\end{align}


\section{Gauge fixed Einstein's equations}
\label{mainsection}
\noindent We study the Cauchy problem in constant mean extrinsic curvature (CMC) gauge. To this end, we assume that the Lorentzian spacetime diffeomorphic to $M\times \mathbb{R}$ admits a constant mean curvature Cauchy slice. 

\noindent Now recall the vacuum Einstein's field equations with a cosmological constant $\Lambda$
\begin{align}
\label{eq:EFET}
 \text{Ric}(\widehat{g})_{\mu\nu}-\frac{1}{2}R(\widehat{g})+\Lambda\widehat{g}_{\mu\nu}=0   
\end{align}
expressed in $3+1$ formulation in a local chart $(t,x^{1},x^{2},x^{3})$ 
\begin{align}
\partial_{t}g_{ij}=-2Nk_{ij}+(\mathscr{L}_{X}g)_{ij},\\
\partial_{t}k_{ij}=-\nabla_{i}\nabla_{j}N+N(R_{ij}+\tr_{g}k k_{ij}-2k^{k}_{i}k_{jk}-\frac{2\Lambda}{n-1}g_{ij})+(\mathscr{L}_{X}k)_{ij},\\
2\Lambda=R(g)-|k|^{2}_{g}+(\tr_{g}k)^{2},\\
0=\nabla^{i}k_{ij}-\nabla_{j}\tr_{g}k.   
\end{align}
This evolution-constraint system leads to the following initial value problem
\begin{definition}
Initial data for \ref{eq:EFET} consist of an n dimensional manifold $M$, a Riemannian metric $g$, a covariant 2-tensor $k$, all assumed to be smooth and to satisfy 
\begin{align}
\label{eq:constraint1}
2\Lambda=R(g)-|k|^{2}_{g}+(\tr_{g}k)^{2},\\
\label{eq:constraint2}
0=\nabla^{i}k_{ij}-\nabla_{j}\tr_{g}k.       
\end{align}
Given a set of initial data, the corresponding initial value problem consists of constructing an $n+1$-dimensional Lorentzian manifold $(\widetilde{M},\widehat{g})$ satisfying the constraint equations \eqref{eq:constraint1} and \eqref{eq:constraint2}, together with an embedding $\iota:M\hookrightarrow \widetilde{M}$
such that $\iota(M)$ is a Cauchy hypersurface in $(\widetilde{M},\widehat{g})$. The embedding is required to satisfy $\iota^{*}\widehat{g}=g$ is the prescribed Riemannian metric on $M$, and, denoting by $\widehat{\mathbf{n}}$ the future-directed unit normal vector field to $\iota(M)$ and by $\kappa$ its second fundamental form in $(\widetilde{M},\widehat{g})$, one has $\iota^{*}\kappa=k$, where $k$ is the given symmetric 2-tensor on $M$. The spacetime $(\widetilde{M},\widehat{g})$ constructed in this manner is referred to as a globally hyperbolic development of the initial data.
\end{definition}

\noindent Now we define the CMC gauge. First, recall that  
\begin{definition}[Constant Mean Curvature (CMC) Gauge] Let \((\widetilde{M}, \widehat{g})\) be a globally hyperbolic Lorentzian manifold of dimension \( (3+1) \), and let \(\{M_t\}_{t \in I} \) denote a foliation of \(\widetilde{M}\) by spacelike Cauchy hypersurfaces defined as level sets of a smooth time function \(t: \mathcal{M} \to \mathbb{R}\). The foliation is said to be in \emph{constant mean curvature gauge} (CMC gauge) if, for each \(t \in I\), the hypersurface \(M_t := \{p \in \widetilde{M} \mid t(p) = t\}\) has constant mean curvature \(\tau(t)\), i.e., the trace of the second fundamental form \(k_t\) of \(M_t\) satisfies \[ \mathrm{tr}_{g_t} k_t = \tau(t) = t,\]
where $g_{t}$ denotes the induced Riemannian metric on $M_{t}$. In this gauge, the time coordinate function $t$ is chosen such that it coincides with the mean curvature of the hypersurface $M_{t}$.
\end{definition}
\noindent The definition automatically requires that the mean extrinsic curvature of level sets of the time be constant on the level sets i.e., $\nabla\tau=0$. More precisely the mean extrinsic curvature $\tau=\tr_{g}k$ of $M$ verifies the following 
\begin{align}
\label{eq:cmcg}
 \tau=\tr_{g}k=\text{monotonic function of $t$ alone}.  
\end{align}
In particular, we choose 
\begin{align}
    t=\tau.
\end{align}
The first analytical issue to address is the existence of a constant mean curvature (CMC) hypersurface $M$ within the spacetime manifold $M\times \mathbb{R}$. Unlike the setting of stability problems, where one typically begins with a known background solution possessing CMC slices, often of negative Einstein type, our situation lacks such an a priori reference geometry. In stability analyses, the presence of a negative Einstein background ensures the existence of CMC slices via a rigidity argument, whereby small perturbations of the data continue to admit CMC foliations. However, in the present context, no such background solution is assumed to be available. Rather, the goal is to determine whether such solutions exist at all, whether generic solutions asymptotically approach them, and whether any geometric or analytic obstructions prevent this convergence.

\noindent Consequently, we impose as a working hypothesis that the spacetime admits at least one CMC hypersurface. This assumption is motivated by the asymptotic structure of the models under consideration, which exhibit a big-bang-type singularity in the past. Specifically, we assume that past-directed inextendible causal geodesics terminate in a curvature singularity, in the sense of the $C^{\infty}$ formulation of the strong cosmic censorship conjecture in the vacuum cosmological setting. Under these conditions, the following result of Marsden and Tipler ensures the existence of a CMC hypersurface:

\begin{theorem-non}[\cite{marsden1980maximal}, Theorem 6]
Let $(\widetilde{M},\widehat{g})$ be a cosmological spacetime. If there exists a Cauchy hypersurface from which all orthogonal, future- or past-directed timelike geodesics terminate in a strong curvature singularity, then the singularity is crushing, and in particular, the spacetime admits a constant mean curvature hypersurface.
\end{theorem-non}
\noindent In the CMC gauge, the first notice a few basic important points. 
\begin{claim}
\label{important}
Let $(\widetilde{M},\widehat{g})=(M\times \mathbb{R},\widehat{g})$ be a $3+1$ dimensional globally hyperbolic spacetime admitting a foliation by constant mean curvature (CMC) hypersurfaces $M_{t}$, with each $M_{t}$ diffeomorphic to a compact manifold MM of negative Yamabe type. Then any initially expanding CMC solution does not admit a maximal hypersurface at any finite time. That is, the mean curvature $\tau(t)$ remains strictly negative for all $t$ in the domain of existence, and the evolution proceeds via continued expansion.
\end{claim}

\begin{proof}
The proof proceeds by a contradiction argument based on the Hamiltonian constraint and the structure of the space of initial data. We aim to establish the impossibility of the existence of a maximal slice (i.e., one for which $\tau=0$) in the future development of an initially expanding solution.

\noindent We begin by recalling that any symmetric $(0,2)$-tensor $k$ on $M$ can be orthogonally decomposed with respect to the standard York splitting (cf. \cite{koiso1979decomposition}):
\begin{align}
k = k^{\mathrm{TT}} + \frac{\tau}{3}g + \left( \mathscr{L}_Z g - \frac{2}{3} (\operatorname{div}_g Z) g \right),
\end{align}
where $k^{TT}$ denotes the transverse-traceless (TT) part of $k$, satisfying $tr_{g}k^{TT}$ and $\div_{g}k^{TT}=0$, and $Z$ is a smooth vector field on $M$.

\noindent Suppose $k$ satisfies the momentum constraint:
\begin{align}
\nabla^j k_{ij} = 0.
\end{align}
We show that the non-TT part involving the vector field Z must vanish. To that end, consider the $L^{2}$-inner product of the momentum constraint with the vector field Z, and apply Stokes’ theorem:
\begin{align}
0 &= \int_M \left(Z^i \nabla^j k_{ij} + Z^j \nabla^i k_{ji} \right) \mu_g \nonumber \
&= \int_M (\nabla^i Z^j + \nabla^j Z^i)\left( \nabla_i Z_j + \nabla_j Z_i - \frac{2}{3} (\operatorname{div}g Z) g{ij} \right) \mu_g \\\nonumber 
&= \int_M \left| \mathscr{L}_Z g - \frac{2}{3} (\operatorname{div}_g Z) g \right|^2  \mu_g.
\end{align}
Since the integrand is non-negative, this integral vanishes if and only if
\begin{align}
\mathscr{L}_Z g - \frac{2}{3} (\operatorname{div}_g Z) g = 0.
\end{align}
Thus, the extrinsic curvature simplifies to:
\begin{align}
k = k^{\mathrm{TT}} + \frac{\tau}{n}g.
\end{align}
We now invoke the Hamiltonian constraint on a CMC hypersurface $M_{t}$, which reads:
\begin{align}
R(g) - |k|^2 + \tau^2 &= 2\Lambda \
\Rightarrow R(g) - \left| k^{\mathrm{TT}} + \frac{\tau}{n}g \right|^2 + \tau^2 = 2\Lambda.
\end{align}
Expanding the norm
\begin{align}
|k^{\mathrm{TT}} + \frac{\tau}{n}g|^2 = |k^{\mathrm{TT}}|^2 + \frac{2\tau}{n} \operatorname{tr}_g k^{\mathrm{TT}} + \frac{\tau^2}{n^{2}} \operatorname{tr}_g g = |k^{\mathrm{TT}}|^2 + \frac{\tau^2}{n},
\end{align}
and using $\tr_{g}k^{TT}=0$ and $\tr_{g}g=n$, we obtain:
\begin{align}
R(g) - |k^{\mathrm{TT}}|^2 = 2\Lambda - \tau^2\left(1 - \frac{1}{3}\right).
\end{align}
For $n=3$, this simplifies to:
\begin{align}
R(g) - |k^{\mathrm{TT}}|^2 = 2\Lambda - \frac{2}{3}\tau^2.
\end{align}
Now, the hypothesis that $M$ is of negative Yamabe type implies that for all conformal representatives $g$ in the conformal class $[g]$, the scalar curvature $R(g)$ cannot be everywhere nonnegative. In particular, for all metrics $g$ conformally related to a constant negative scalar curvature metric, $R(g)<0$ somewhere. Thus, the left-hand side is strictly negative at some point, which enforces the inequality:
\begin{align}
\quad \tau^2 > 3\Lambda.
\end{align}
Since $\tau$ is constant on each CMC hypersurface, this implies that $\tau^{2}>3\Lambda$, hence, in particular, $\tau$ cannot vanish. Therefore, a maximal slice (where $\tau=0$) cannot form during the evolution.
Finally, since $\tau(t)$ is monotonically increasing in time by the choice of CMC gauge (cf. Equation~\eqref{eq:cmcg}), any initially expanding solution (i.e., with $\tau<0$) must remain strictly expanding throughout its domain of existence, with $\tau<0$ for all $t$. This concludes the proof.
\end{proof}
\begin{remark}
For convenience, one could simply set $\Lambda=n(n+1)$.
\end{remark}

\noindent A central function in the construction of a dimensionless dynamical formulation of the Einstein vacuum equations with cosmological constant $\Lambda$ is the use of a natural geometric scaling governed by the mean curvature. Consider the entity
\[
\varphi^2 := \tau^2 - 3\Lambda=\tau^{2}-3n(n+1) > 0,
\]
which is strictly positive and decays monotonically toward the future in a chosen temporal gauge (claim \ref{important}). The function $\varphi := \sqrt{\tau^2 - 3n(n+1)}$ serves as a natural scaling factor for the spatial geometry and lapse-shift data. To obtain a dimensionally consistent and dynamically well-posed evolution system, one must perform a rescaling of the Einstein field equations in terms of $\varphi$.

\noindent We adopt the dimensional convention of \cite{andersson2011einstein}, wherein the local spatial coordinates $\{x^i\}$ on the Cauchy hypersurface $M$ are taken to be dimensionless. Under this convention, the geometric and gauge quantities possess the following physical dimensions:
\begin{align} \label{eq:scaling}
[g_{ij}] = L^2, \quad [k_{ij}] = L, \quad [X^i] = L, \quad [N] = L^2, \quad [\tau] = L^{-1}, \quad [\varphi^2] = L^{-2},
\end{align}
where $L$ denotes the dimension of length. 

\noindent To facilitate the introduction of dimensionless variables, we adopt the notational convention that dimensional quantities are denoted with a tilde, while their dimensionless counterparts are written without decoration. The natural rescaling is then given by
\begin{align} \label{eq:scaling}
\widetilde{g}_{ij} = \frac{1}{\varphi^2} g_{ij}, \qquad \widetilde{N} = \frac{1}{\varphi^2} N, \qquad \widetilde{X}^i = \frac{1}{\varphi} X^i, \qquad \widetilde{k}_{ij}^{TT} = \frac{1}{\varphi} \Sigma_{ij},
\end{align}
where $\Sigma_{ij}$ denotes the trace-free part of the extrinsic curvature, and we define $\varphi = -\sqrt{\tau^2 - 3n(n+1)}$ so that $\varphi / \tau > 0$.

\noindent We now introduce a reparametrization of time via a new variable $T = T(\tau)$, defined implicitly by the vector field transformation
\begin{align} \label{eq:timechange}
\partial_T = -\frac{\varphi^2}{\tau} \partial_\tau.
\end{align}
Integrating this relation yields
\begin{align} \label{eq:ex1}
\varphi(T) = \varphi(T_0) e^{T_0} e^{-T},
\end{align}
for some fixed initial time $T_0$. By setting the normalization condition $\varphi(T_0) e^{T_0} = -C$ for a constant $C$ of dimension $L^{-1}$, we obtain explicit expressions for $\varphi$ and $\tau$ as functions of the new time coordinate:
\begin{align} \label{eq:ex2}
\varphi(T) = -Ce^{-T}, \qquad \tau(T) = -\sqrt{C^{2}e^{-2T} + 3n(n+1)}.
\end{align}
The new time variable $T$ ranges over the entire real line, i.e., $T \in (-\infty, \infty)$, and behaves analogously to Newtonian time in the rescaled formulation. 

\noindent Moreover, we record the asymptotic estimate
\begin{align} \label{eq:asymptotic}
\frac{\varphi(T)}{\tau(T)} = \frac{Ce^{-T}}{\sqrt{C^{2}e^{-2T} + 3n(n+1)}} \leq \frac{Ce^{-T}}{\sqrt{3n(n+1)}} \sim n^{-\frac{1}{2}}(n+1)^{-\frac{1}{2}}e^{-T},
\end{align}
demonstrating the exponential decay of the ratio $\varphi/\tau$ in the future time direction.




\noindent The rescaled Einstein's equations in the CMC gauge can be cast into the following system
\begin{align}
\label{eq:cmc1}
 \frac{\partial g_{ij}}{\partial T}=-\frac{2\varphi}{\tau}N\Sigma_{ij}-2(1-\frac{N}{n})g_{ij}-\frac{\varphi}{\tau}(\mathscr{L}_{X}g)_{ij}\\
 \frac{\partial \Sigma_{ij}}{\partial T}=-(n-1)\Sigma_{ij}-\frac{\varphi}{\tau}N(\text{Ric}_{ij}+\frac{n-1}{n^{2}}g_{ij})+\frac{\varphi}{\tau}\nabla_{i}\nabla_{j}(\frac{N}{n}-1)+\frac{2\varphi}{\tau}N\Sigma_{ik}\Sigma^{k}_{j}\\
 -\frac{\varphi}{n\tau}(\frac{N}{n}-1)g_{ij}-(n-2)(\frac{N}{n}-1)\Sigma_{ij}-\frac{\varphi}{\tau}(\mathscr{L}_{X}\Sigma)_{ij}\\
 \label{eq:HC11}
 R+\frac{n-1}{n}-|\Sigma|^{2}=0,\\
 \label{eq:cmc2}
 \nabla_{j}\Sigma^{j}_{i}=0.
\end{align}
For the purpose of our study, this system is not suitable. Instead of treating this system as a coupled weakly wave system (which would inevitably require invoking spatial harmonic gauge to turn the system into a strong hyperbolic system), we will treat the equation for the metric components $g_{ij}$ as transport equations. We define the following new entity that we shall study 
\begin{align}
    \mathfrak{T}_{ij}:=\text{Ric}_{ij}-\frac{1}{3}R(g)g_{ij}.
\end{align}
We obtain the following manifestly wave system for $(\Sigma,\mathfrak{T})$.
\begin{remark}
 We note that Choquet-Bruhat and York \cite{YCBY} considered a similar coupled system for $\Sigma$ and $\text{Ric}$.
\end{remark}
\begin{lemma}[Coupled Wave System for \((\Sigma, \mathfrak{T})\)]
\label{evolutionsystem}
Let \( (g, \Sigma) \) satisfy the reduced Einstein evolution equations \eqref{eq:cmc1}--\eqref{eq:cmc2} in constant mean curvature (CMC) and transported spatial gauge. Then the pair \( (\Sigma, \mathfrak{T}) \), consisting respectively of the transverse-traceless part of the second fundamental form and the obstruction tensor $\mathfrak{T}$, satisfies the following system of coupled second-order quasilinear hyperbolic equations, modulo a spatial diffeomorphism \( \Psi \) generated by a smooth shift vector field \( X \).

\noindent The evolution equations for \( \Sigma_{ij} \) and \( \mathfrak{T}_{ij} \) read:
\begin{align}
\partial_T \Sigma_{ij} &= -\frac{\varphi}{\tau} (\mathscr{L}_X \Sigma)_{ij} 
- (n-1) \Sigma_{ij}
- \frac{\varphi}{\tau} N \mathfrak{T}_{ij} 
+ \frac{\varphi}{\tau} \nabla_i \nabla_j \left( \frac{N}{n} - 1 \right) 
+ \frac{2\varphi}{\tau} N \Sigma_{ik} \Sigma^k{}_j \notag\\
&\quad - \frac{\varphi}{n\tau} \left( \frac{N}{n} - 1 \right) g_{ij} 
- (n-2) \left( \frac{N}{n} - 1 \right) \Sigma_{ij}, \label{eq:sigma_evol}
\end{align}
\begin{align}
\partial_T \mathfrak{T}_{ij} &= -\frac{\varphi}{\tau} (\mathscr{L}_X \mathfrak{T})_{ij} 
- N \frac{\varphi}{\tau} \left( \Delta_g \Sigma_{ij} - R^m{}_{jli} \Sigma^l{}_m - R^m{}_{ilj} \Sigma^l{}_m \right) \notag\\
&\quad + n \frac{\varphi}{\tau} \nabla^l \nabla_i \left( \frac{N}{n} - 1 \right) \Sigma_{jl} 
+ n \frac{\varphi}{\tau} \nabla^l \left( \frac{N}{n} - 1 \right) \nabla_i \Sigma_{jl} \notag\\
&\quad + n \frac{\varphi}{\tau} \nabla^l \nabla_j \left( \frac{N}{n} - 1 \right) \Sigma_{il}
+ n \frac{\varphi}{\tau} \nabla^l \left( \frac{N}{n} - 1 \right) \nabla_j \Sigma_{il} \notag\\
&\quad - n \frac{\varphi}{\tau} \Delta_g \left( \frac{N}{n} - 1 \right) \Sigma_{ij}
- 2n \frac{\varphi}{\tau} \nabla^l \left( \frac{N}{n} - 1 \right) \nabla_l \Sigma_{ij} \notag\\
&\quad - \Delta_g \left( \frac{N}{n} - 1 \right) g_{ij}
+ N \frac{\varphi}{\tau} \left( \mathfrak{T}_{ki} \Sigma^k{}_j + \mathfrak{T}_{kj} \Sigma^k{}_i \right)
+ \frac{2(n-1)}{n^2} \left( \frac{N}{n} - 1 \right) g_{ij}. \label{eq:T_evol}
\end{align}
\end{lemma}


\begin{proof}
 The proof follows directly from Einstein's equations and the variation formula for Ricci and scalar curvature. 
 
 %Recall the definition $\mathfrak{T}:=\text{Ric}[g]+\frac{n-1}{n^{2}}g$. Now recall the variation formula for the Ricci curvature 
 %\begin{align}
 % \frac{\partial R_{ij}}{\partial T}=\frac{1}{2}\left(\nabla^{l}\nabla_{i}\partial_{T}g_{jl}+\nabla^{l}\nabla_{j}\partial_{T}g_{il}-\nabla^{l}\nabla_{l}\partial_{T}g_{ij}\right)-\frac{1}{2}\nabla_{i}\nabla_{j}\tr_{g}\partial_{T}g_{ij}.     
 %\end{align}
 %Compute each term individually using the evolution equation for the metric $g_{ij}$. 
 %\begin{align}
 % \nabla^{l}\nabla_{i}\partial_{T}g_{jl}&=&-\frac{2\varphi}{\tau}\left((\nabla^{l}\nabla_{i}N)\Sigma_{jl}+\nabla_{i}N\nabla^{l}\Sigma_{jl}+\nabla^{l}N\nabla_{i}\Sigma_{jl}+N\nabla^{l}\nabla_{i}\Sigma_{jl}\right)\\\nonumber 
 %&&+2\nabla_{i}\nabla_{j}(\frac{N}{n}-1)-\frac{\varphi}{\tau}\nabla^{l}\nabla_{i}\mathscr{L}_{X}g_{lj}   
 %\end{align}
 %Similarly $i\leftrightarrow j$, one obtains 
%\begin{align}
% \nabla^{l}\nabla_{j}\partial_{T}g_{il}&=&-\frac{2\varphi}{\tau}\left((\nabla^{l}\nabla_{j}N)\Sigma_{il}+\nabla_{j}N\nabla^{l}\Sigma_{il}+\nabla^{l}N\nabla_{j}\Sigma_{il}+N\nabla^{l}\nabla_{j}\Sigma_{il}\right)\\\nonumber 
% &&+2\nabla_{i}\nabla_{j}(\frac{N}{n}-1)-\frac{\varphi}{\tau}\nabla^{l}\nabla_{j}\mathscr{L}_{X}g_{il}.   
%\end{align}
%Now notice that there are two terms $N\nabla^{l}\nabla_{i}\Sigma_{jl}$ and $N\nabla^{l}\nabla_{j}\Sigma_{il}$ involving the second derivative of $\Sigma$ that can further be reduced using the momentum constraint $\nabla^{j}\Sigma_{ij}=0$ i.e., by Ricci identity
%\begin{align}
% \nabla_{l}\nabla_{i}\Sigma^{l}_{j}=\nabla_{i}\nabla_{l}\Sigma^{l}_{j}+R^{l}~_{kli}\Sigma^{k}_{j}-R^{m}~_{jli}\Sigma^{l}_{m}=R_{ki}\Sigma^{k}_{j}-R^{m}~_{jli}\Sigma^{l}_{m}.
%\end{align}
%This observation yields the ellipticity of the map $\Sigma\mapsto \partial_{T}\text{Ric}$ thereby making the evolution system $(\partial_{T}\Sigma,\partial_{T}\mathfrak{T})$ manifestly hyperbolic.   
%The last relation that we need is the following identity regarding the Lie derivative and the Fr\'echet derivative 
%\begin{align}
 %\mathscr{L}_{X}\text{Ric}=D\text{Ric}\cdot\mathscr{L}_{X}g.   
%\end{align}
%Therefore the terms involving shift vector field $X$ combine to form the Lie derivative term $\mathscr{L}_{X}\mathfrak{T}$. The collection of every term completes the proof. 
\end{proof}

\noindent There are advantages of working with the pair $(\Sigma,\mathfrak{T})$ instead of $(g,\Sigma)$. Firstly, the equations are manifestly hyperbolic up to a spatial diffeomorphism and one does not need to deal with the Gribov ambiguities associated with the spatial harmonic gauge often used. In particular, the spatial harmonic gauge is not suited for large data problems in the current context without substantial technical machinery.      

\noindent In the previous section, we have chosen the CMC gauge as the time gauge. We need to choose a spatial gauge to analyze the evolution system \ref{evolutionsystem}. We choose CMC transported spatial gauge (previously \cite{fajman2022cosmic,rodnianski2018regime} used this gauge to study the dynamical stability of big-bang singularity formation) 
\subsection{Transported Spatial Coordinates and the Re-Scaled Einstein System}
\noindent In this section, we introduce the notion of a \emph{transportedspatial gauge}, which serves as a coordinate choice adapted to the constant mean curvature (CMC) foliation. 

\begin{definition}[Transported Spatial Gauge]
\label{spatial}
Let $(x^1_0, x^2_0, x^3_0)$ be local coordinates on an open neighborhood $\Omega \subset M$ of the initial Cauchy hypersurface $\{T' = T_0\} \subset \widetilde{M} := \mathbb{R} \times M$. We say that a coordinate system $(T', x^1, x^2, x^3)$ on $[T_0, T] \times \Omega$ defines a \emph{transported spatial gauge} if the spatial coordinate functions are Lie transported along the future-directed unit normal vector field $\widehat{\mathbf{n}}$ to the CMC hypersurfaces, i.e.,
\begin{align}
\label{eq:transport}
\mathscr{L}_{\widehat{\mathbf{n}}}(x^i) = 0, \qquad i = 1,2,3,
\end{align}
with initial condition $(x^1, x^2, x^3)|_{T'=T_0} = (x^1_0, x^2_0, x^3_0)$.
\end{definition}

\noindent
From~\eqref{eq:transport}, and the fact that $\widehat{\mathbf{n}} = \frac{1}{3}(\partial_{T'} - X^i \partial_i)$ in general coordinates, we deduce:
\begin{align}
\label{eq:shift-vanishing}
\mathscr{L}_{\widehat{\mathbf{n}}}(x^{i})=\widehat{\mathbf{n}}(x^i) = \frac{1}{3}(\partial_{T'} - X^j \partial_j)(x^i) = -\frac{1}{3} X^i \Rightarrow X^i = 0.
\end{align}
Thus, the shift vector field vanishes identically in these coordinates.

\vspace{1em}
\noindent
In this gauge, the rescaled Einstein evolution equations take a particularly tractable form. Let $\Sigma_{ij}$ denote the trace-free part of the second fundamental form (the shear), and let $\mathfrak{T}_{ij}$ denote the trace-free part of the rescaled Ricci curvature. The evolution equations then read:

\begin{align}
\label{eq:curvature1}
\partial_T \Sigma_{ij} &= -(n-1)\Sigma_{ij} - \frac{\varphi}{\tau} N \mathfrak{T}_{ij} + \frac{\varphi}{\tau} \nabla_i \nabla_j \left( \frac{N}{n} - 1 \right) + \frac{2\varphi}{\tau} N \Sigma_{ik} \Sigma^k{}_j \\
&\quad - \frac{\varphi}{n\tau} \left( \frac{N}{n} - 1 \right) g_{ij} - (n-2) \left( \frac{N}{n} - 1 \right) \Sigma_{ij}\nonumber, 
\end{align}

\begin{align}
 \label{eq:curvature2}
\partial_T \mathfrak{T}_{ij} &= -N\frac{\varphi}{\tau} \left( \Delta_g \Sigma_{ij} - R^m{}_{jli} \Sigma^l{}_m - R^m{}_{ilj} \Sigma^l{}_m \right) + n \frac{\varphi}{\tau} \nabla^l \nabla_i \left( \frac{N}{n} - 1 \right) \Sigma_{jl} \\
&\quad + n \frac{\varphi}{\tau} \nabla^l \left( \frac{N}{n} - 1 \right) \nabla_i \Sigma_{jl} + n \frac{\varphi}{\tau} \nabla^l \nabla_j \left( \frac{N}{n} - 1 \right) \Sigma_{il} \nonumber \\
&\quad + n \frac{\varphi}{\tau} \nabla^l \left( \frac{N}{n} - 1 \right) \nabla_j \Sigma_{il} - n \frac{\varphi}{\tau} \Delta_g \left( \frac{N}{n} - 1 \right) \Sigma_{ij} \nonumber \\
&\quad - 2n \frac{\varphi}{\tau} \nabla^l \left( \frac{N}{n} - 1 \right) \nabla_l \Sigma_{ij} - \Delta_g \left( \frac{N}{n} - 1 \right) g_{ij} \nonumber \\
&\quad + N \frac{\varphi}{\tau} \left( \mathfrak{T}_{ki} \Sigma^k{}_j + \mathfrak{T}_{kj} \Sigma^k{}_i \right) + \frac{2(n-1)}{n^2} \left( \frac{N}{n} - 1 \right) g_{ij}\nonumber, 
\end{align}

\noindent
These equations are supplemented by the elliptic lapse equation:
\begin{align}
\label{eq:lapse3}
- \Delta_g \left( \frac{N}{n} - 1 \right) + \left( |\Sigma|^2 + \frac{1}{3} \right) \left( \frac{N}{n} - 1 \right) = - |\Sigma|^2.
\end{align}







\subsection{Integration and Norms}

\noindent Let \( (M, g) \) be a smooth, closed, oriented Riemannian $3$-manifold. Fix a smooth partition of unity \( \{ \xi_U \}_{U \in \mathcal{U}} \) subordinate to a finite atlas \( \{ (U, \varphi_U) \}_{U \in \mathcal{U}} \) of coordinate charts \( \varphi_U: U \to \mathbb{R}^3 \). For a measurable function \( f: M \to \mathbb{R} \), the integral of \( f \) with respect to the Riemannian volume form \( \mu_g \) is defined by

\begin{align} 
\label{eq:integration}
\int_M f \, \mu_g := \sum_{U \in \mathcal{U}} \int_{\varphi_U(U)} f \circ \varphi_U^{-1}(x) \, \xi_U \circ \varphi_U^{-1}(x) \sqrt{\det g_{ij}(x)} \, dx^1 dx^2 dx^3,
\end{align}

\noindent where \( g_{ij} \) denotes the components of the metric in the chart \( \varphi_U \), and \( \det g_{ij} \) is the determinant of the metric matrix.

\noindent Given a smooth tensor field \( \psi \in \Gamma\left( (T^{\otimes L} M) \otimes (T^* M)^{\otimes K} \right) \), where \( K + L \geq 1 \), we define its pointwise norm using the Riemannian metric \( g \) by

\begin{align} \label{eq:pointwise-norm}
\langle \psi(x), \psi(x) \rangle_g := g^{i_1 j_1} \cdots g^{i_L j_L} g_{m_1 n_1} \cdots g_{m_K n_K} \,
\psi^{m_1 \dots m_K}_{i_1 \dots i_L}(x) \, \psi^{n_1 \dots n_K}_{j_1 \dots j_L}(x).
\end{align}

\noindent This inner product is well-defined and independent of the choice of coordinates. The \( L^p \) norm of \( \psi \) over \( (M, g) \) is then defined by

\begin{align} \label{eq:Lp-norm}
\| \psi \|_{L^p(M)}^p := \int_M \left( \langle \psi, \psi \rangle_g \right)^{p/2} \, \mu_g, \qquad 1 \leq p < \infty,
\end{align}

\noindent and the essential supremum (or \( L^\infty \)) norm is given by

\begin{align} \label{eq:Linf-norm}
\| \psi \|_{L^\infty(M)} := \sup_{x \in M} \left( \langle \psi(x), \psi(x) \rangle_g \right)^{1/2}.
\end{align}

\noindent We now define the energy norms used in the analysis of the Einstein$-\Lambda$ flow. Let \( \Sigma \) denote the trace-free part of the second fundamental form and \( \mathfrak{T} \) the associated energy-momentum correction tensor. For integers \( I \in \{0, 1, 2, 3\} \), define:

\begin{align} \label{eq:OI}
\mathcal{O}_I := 
\begin{cases}
\| \nabla^I \Sigma \|_{L^2(M)} + \| \nabla^{I-1} \mathfrak{T} \|_{L^2(M)}, & \text{if } I \geq 1, \\
\| \Sigma \|_{L^2(M)}, & \text{if } I = 0.
\end{cases}
\end{align}

\noindent We also introduce exponentially weighted norms in time to account for the expanding geometry in cosmological spacetimes. Let \( T \) denote the CMC time parameter, and let \( 1\geq \gamma > 0 \) be a scaling parameter (which will ultimately be chosen to be \( \gamma = 1 \)). For \( I \leq 2 \), define the weighted Sobolev norms

\begin{align} \label{eq:FI}
\mathcal{F}_I := \| e^{\gamma T} \nabla^I \Sigma \|_{L^2(M)}.
\end{align}

\noindent Define also the pointwise weighted norm:
\begin{align} \label{eq:infty-norm}
\mathcal{N}^\infty := \| e^{\gamma T} \Sigma \|_{L^\infty(M)} + \left\| e^{2\gamma T} \left( \frac{N}{n} - 1 \right) \right\|_{L^\infty(M)},
\end{align}
where \( N \) denotes the spacetime lapse function.

\noindent The total energy norms used throughout this work are defined by

\begin{align} \label{eq:O}
\mathcal{O} := \sum_{I=0}^3 \mathcal{O}_I, \qquad \mathcal{F} := \sum_{I=0}^2 \mathcal{F}_I.
\end{align}
As will be demonstrated in the analysis below, the optimal exponential weight corresponds to \( \gamma = 1 \), which balances the natural scaling of the lapse and the shear in expanding CMC spacetimes.

\section{Construction of the Initial data}
\label{data}
\noindent The heuristic of the data characterization is provided in the introduction. The goal of this section is to explicitly construct such data that verify the Einstein-$\Lambda$ constraint equations. In addition, we specifically show that such data does not fall under the category treated by the previous mathematically rigorous studies (e.g., \cite{fajman2020stable}) in this exact context. The most important point to note here is that we are in the negative Yamabe context and apply the transverse-traceless perturbation and conformal technique. First, we construct Riemannian metrics on $M$ that verify the condition 
\[
\sum_{I = 0}^{2}  \| \nabla^I \mathfrak{T}[g_1] \|_{L^2(M)}  \leq \mathcal{I}^0=O(e^{a/10}),~a\gg 1
\]
The second step is to solve the momentum constraint in the CMC gauge i.e., construct an appropriate $g_{1}$ $TT-$ tensor verifying the estimate 
\[
\|\Sigma\|_{H^{2}}\lesssim e^{-a}\mathcal I^{0},
\qquad
\|\Sigma\|_{H^{3}}\lesssim \mathcal I^{0},
\]
In the third step, we solve the Hamiltonian constraint using the conformal method and prove that the estimates proved for $\mathfrak{T}$ and $\Sigma$ are modified by a negligible amount in the conformal transformation process. The main proposition that we prove is the existence of an open set of initial data $(g_{0},k_{0})$ that verifies the estimates 
\[
\sum_{I = 0}^{3} \| \nabla^I \Sigma_0 \|_{L^2(M)} + \sum_{I = 0}^{2} \left( \| \nabla^I \mathfrak{T}[g_0] \|_{L^2(M)} + \| e^a \nabla^I \Sigma_0 \|_{L^2(M)} \right) \leq \mathcal{I}^0
\]

\noindent First, we recall the well-known lemma relating the trace-free Ricci curvature of $\gamma$ and the perturbed metric $g=\gamma+\epsilon h$ for $\epsilon\ll1$.

\begin{lemma}
\label{nonlinearexpansion}
Let $(M^{n},\gamma)$, $n\geq 2$, be a smooth closed connected Riemannian manifold, and let
$h\in C^\infty(\text{Sym}^2T^*M)$ be a smooth symmetric $(0,2)$-tensor. For $\epsilon\in\mathbb R^{+}$
with $\epsilon$ sufficiently small, define
\[
g_\epsilon:=\gamma+\epsilon h.
\]
Then $g_\epsilon$ is a smooth Riemannian metric, and the following assertions hold. First, the inverse metric expression reads
\begin{equation}\label{eq:inverse-expansion}
(g_{\epsilon})^{ij}
=
\gamma^{ij}
-
\epsilon h^{ij}
+
\epsilon^{2}O_{\gamma}(h*h).
\end{equation}
\medskip
\noindent Now let $\nabla=\nabla^\gamma$ denote the Levi--Civita connection of $\gamma$, and let
\[
A_\epsilon{}^{k}{}_{ij}
:=
\Gamma(g_\epsilon)^{k}{}_{ij}-\Gamma(\gamma)^{k}{}_{ij}.
\]
Then
\begin{equation}
\label{eq:A-exact}
A_\epsilon{}^{k}{}_{ij}
=
\frac12 (g_\epsilon)^{k\ell}
\big(
\nabla_i(g_\epsilon)_{j\ell}
+\nabla_j(g_\epsilon)_{i\ell}
-\nabla_\ell(g_\epsilon)_{ij}
\big)
=
\frac{\epsilon}{2}(g_\epsilon)^{k\ell}
\big(
\nabla_i h_{j\ell}
+\nabla_j h_{i\ell}
-\nabla_\ell h_{ij}
\big).
\end{equation}
Moreover,
\begin{equation}
\label{eq:Ric-exact}
\text{Ric}(g_\epsilon)_{ij}
=
\text{Ric}(\gamma)_{ij}
+\nabla_k A_\epsilon{}^{k}{}_{ij}
-\nabla_j A_\epsilon{}^{k}{}_{ik}
+A_\epsilon{}^{k}{}_{ij}A_\epsilon{}^{\ell}{}_{k\ell}
-A_\epsilon{}^{k}{}_{i\ell}A_\epsilon{}^{\ell}{}_{jk},
\end{equation}
and hence
\begin{equation}
\label{eq:R-exact}
R[g_\epsilon]
=
(g_\epsilon)^{ij}
\Big(
\text{Ric}(\gamma)_{ij}
+\nabla_k A_\epsilon{}^{k}{}_{ij}
-\nabla_j A_\epsilon{}^{k}{}_{ik}
+A_\epsilon{}^{k}{}_{ij}A_\epsilon{}^{\ell}{}_{k\ell}
-A_\epsilon{}^{k}{}_{i\ell}A_\epsilon{}^{\ell}{}_{jk}
\Big).
\end{equation}
\noindent The scalar curvature admits the expansion
\begin{equation}
\label{eq:R-linearization-expansion}
R[g_\epsilon]
=
R[\gamma]+\epsilon\,DR_\gamma[h]+\epsilon^2 Q^R_\gamma(h;\epsilon),
\end{equation}
where
\begin{equation}
\label{eq:DR}
DR_\gamma[h]
=
\nabla^i\nabla^j h_{ij}
-\Delta_\gamma(\tr_\gamma h)
-\langle \text{Ric}(\gamma),h\rangle_\gamma.
\end{equation}
Here $\Delta_\gamma=\gamma^{ij}\nabla_i\nabla_j$ is the rough Laplacian on functions. The remainder $Q^R_\gamma(h;\epsilon)$ is smooth in $(x,\epsilon)$ and in schematic notation,
\begin{equation}
\label{eq:QR-schematic}
Q^R_\gamma(h;\epsilon)
=
h*\nabla^2 h+\nabla h * \nabla h+\text{Rm}(\gamma)*h*h.
\end{equation}
\noindent Now if
\[
\mathfrak T[g]
:=\text{Ric}(g)-\frac1n R[g]\,g.
\]
Then
\begin{equation}
\label{eq:T-expansion}
\mathfrak T[g_\epsilon]
=
\mathfrak T[\gamma]
+\epsilon\,D\mathfrak T_\gamma[h]
+\epsilon^2 Q^{\mathfrak T}_\gamma(h;\epsilon),
\end{equation}
where
\begin{equation}
\label{eq:DT-general}
D\mathfrak T_\gamma[h]
=
D\text{Ric}_\gamma[h]
-\frac1n\,DR_\gamma[h]\,\gamma
-\frac1n\,R[\gamma]\,h,
\end{equation}
and
\begin{align}
\label{eq:DRic-general}
D\text{Ric}_\gamma[h]_{ij}
&=
\frac12\Big(
-\Delta_\gamma h_{ij}
-\nabla_i\nabla_j(\tr_\gamma h)
+\nabla_i(\delta_\gamma h)_j
+\nabla_j(\delta_\gamma h)_i
\Big)
-\text{Rm}(\gamma)_{ikj\ell}h^{k\ell}
+\frac12\Big(\text{Ric}(\gamma)_{ik}h^{k}{}_{j}+\text{Ric}(\gamma)_{jk}h^{k}{}_{i}\Big).
\end{align}
Here
\[
(\delta_\gamma h)_j:=\nabla^i h_{ij}.
\]
The quadratic remainder $Q^{\mathfrak T}_\gamma(h;\epsilon)$ is smooth in $(x,\epsilon)$ and, 
in particular,
\begin{equation}
\label{eq:QT-schematic}
Q^{\mathfrak T}_\gamma(h;\epsilon)
=
h*\nabla^2 h+\nabla h * \nabla h+\text{Rm}(\gamma)*h*h.
\end{equation}

\medskip
\noindent Assume in addition that
\[
\tr_\gamma h=0,
\qquad
\delta_\gamma h=0.
\]
Then
\begin{equation}
\label{eq:DR-TT}
DR_\gamma[h]=-\langle \text{Ric}(\gamma),h\rangle_\gamma
=-\langle \mathfrak T[\gamma],h\rangle_\gamma,
\end{equation}
and
\begin{equation}
\label{eq:DRic-TT}
D\text{Ric}_\gamma[h]
=
\frac12\,\Delta_{L,\gamma}h,
\end{equation}
where the Lichnerowicz Laplacian on symmetric $(0,2)$-tensors is defined by
\begin{equation}
\label{eq:Lich-def}
(\Delta_{L,\gamma}h)_{ij}
:=
-\Delta_\gamma h_{ij}
-2\,\text{Rm}(\gamma)_{ikj\ell}h^{k\ell}
+\text{Ric}(\gamma)_{ik}h^{k}{}_{j}
+\text{Ric}(\gamma)_{jk}h^{k}{}_{i}.
\end{equation}
Consequently,
\begin{equation}
\label{eq:DT-TT}
D\mathfrak T_\gamma[h]_{ij}
=
\frac12(\Delta_{L,\gamma}h)_{ij}
+\frac1n\langle \mathfrak T[\gamma],h\rangle_\gamma\,\gamma_{ij}
-\frac1n R[\gamma]\,h_{ij}.
\end{equation}
\end{lemma}



\begin{proof}
For $\epsilon$ sufficiently small, the bilinear form $g_\epsilon=\gamma+\epsilon h$ remains positive
definite, hence defines a smooth Riemannian metric. Since $M$ is compact and $h$ is fixed, all
smallness thresholds and constants below depend only on finitely many norms of $(\gamma,h)$.

\noindent We write $\nabla=\nabla^\gamma$. Since $\nabla\gamma=0$ and $(g_\epsilon)_{ij}=\gamma_{ij}+\epsilon h_{ij}$,
the difference tensor between the Levi--Civita connections of $g_\epsilon$ and $\gamma$ is given by the
standard formula
\[
A_\epsilon{}^{k}{}_{ij}
=
\Gamma(g_\epsilon)^{k}{}_{ij}-\Gamma(\gamma)^{k}{}_{ij}
=
\frac12 (g_\epsilon)^{k\ell}
\big(
\nabla_i(g_\epsilon)_{j\ell}
+\nabla_j(g_\epsilon)_{i\ell}
-\nabla_\ell(g_\epsilon)_{ij}
\big),
\]
which immediately yields \eqref{eq:A-exact}. The curvature transformation law for two torsion-free
connections implies
\[
\text{Ric}(g_\epsilon)_{ij}
=
\text{Ric}(\gamma)_{ij}
+\nabla_k A_\epsilon{}^k{}_{ij}
-\nabla_j A_\epsilon{}^k{}_{ik}
+A_\epsilon{}^k{}_{ij}A_\epsilon{}^\ell{}_{k\ell}
-A_\epsilon{}^k{}_{i\ell}A_\epsilon{}^\ell{}_{jk},
\]
hence also \eqref{eq:R-exact} after contraction with $(g_\epsilon)^{ij}$.
\noindent Next, since
\[
g_\epsilon=\gamma\circ(\text{Id}+\epsilon \gamma^{-1}h),
\]
the inverse metric is obtained by Neumann expansion:
\[
g_\epsilon^{-1}
=
(\text{Id}+\epsilon \gamma^{-1}h)^{-1}\circ \gamma^{-1}
=
\big(\text{Id}-\epsilon \gamma^{-1}h+\epsilon^2(\gamma^{-1}h)^2+\epsilon^3\mathcal E_\epsilon\big)\circ \gamma^{-1},
\]
which gives \eqref{eq:inverse-expansion}. The $C^m$ bounds on $\mathcal E_\epsilon$ follow from
smooth dependence of matrix inversion on the coefficients of the metric and compactness of $M$.

\noindent We now derive the linearization of scalar curvature. Differentiating \eqref{eq:A-exact} at $\epsilon=0$ gives
\[
\dot A^{k}{}_{ij}
:=
\left.\frac{d}{d\epsilon}\right|_{\epsilon=0}A_\epsilon{}^{k}{}_{ij}
=
\frac12\gamma^{k\ell}
\big(
\nabla_i h_{j\ell}
+\nabla_j h_{i\ell}
-\nabla_\ell h_{ij}
\big).
\]
Differentiating \eqref{eq:R-exact} at $\epsilon=0$, using $(g_\epsilon)^{ij}=\gamma^{ij}+O(\epsilon)$ and
$A_\epsilon=O(\epsilon)$, we find
\[
DR_\gamma[h]
=
-h^{ij}\text{Ric}(\gamma)_{ij}
+\gamma^{ij}\big(\nabla_k\dot A^{k}{}_{ij}-\nabla_j\dot A^{k}{}_{ik}\big).
\]
A direct computation shows
\[
\gamma^{ij}\nabla_k\dot A^{k}{}_{ij}
=
\nabla^i\nabla^j h_{ij}-\frac12\Delta_\gamma(\tr_\gamma h),
\]
while
\[
\gamma^{ij}\nabla_j\dot A^{k}{}_{ik}
=
\frac12\Delta_\gamma(\tr_\gamma h).
\]
Substituting these identities yields \eqref{eq:DR}.

\noindent The quadratic expansion \eqref{eq:R-linearization-expansion} follows by substituting
\eqref{eq:inverse-expansion} and \eqref{eq:A-exact} into \eqref{eq:R-exact}. Indeed,
\[
A_\epsilon=\epsilon\,\nabla h+\epsilon^2 h*\nabla h,
\qquad
\nabla A_\epsilon=\epsilon\,\nabla^2 h+\epsilon^2(\nabla h*\nabla h+h*\nabla^2 h),
\]
and the exact formula \eqref{eq:R-exact} then yields
\[
R[g_\epsilon]-R[\gamma]-\epsilon DR_\gamma[h]
=
\epsilon^2\Big(h*\nabla^2 h+\nabla h * \nabla h+\text{Rm}(\gamma)*h*h\Big),
\]
schematically.

\noindent We turn to the trace-free Ricci tensor
\[
\mathfrak T[g]=\text{Ric}(g)-\frac1nR[g]\,g.
\]
Differentiating at $\gamma$ gives
\[
D\mathfrak T_\gamma[h]
=
D\text{Ric}_\gamma[h]-\frac1n DR_\gamma[h]\gamma-\frac1nR[\gamma]h,
\]
which is \eqref{eq:DT-general}. Thus, it remains only to record the standard formula for
$D\text{Ric}_\gamma[h]$.

\noindent Differentiating \eqref{eq:Ric-exact} at $\epsilon=0$ eliminates the quadratic terms in $A_\epsilon$ and yields
\[
D\text{Ric}_\gamma[h]_{ij}
=
\nabla_k\dot A^k{}_{ij}-\nabla_j\dot A^k{}_{ik}.
\]
Substituting the expression for $\dot A$ and commuting covariant derivatives in the third-order terms
gives
\[
D\text{Ric}_\gamma[h]_{ij}
=
\frac12\Big(
-\Delta_\gamma h_{ij}
-\nabla_i\nabla_j(\tr_\gamma h)
+\nabla_i(\delta_\gamma h)_j
+\nabla_j(\delta_\gamma h)_i
\Big)
-\text{Rm}(\gamma)_{ikj\ell}h^{k\ell}
+\frac12\Big(\text{Ric}(\gamma)_{ik}h^k{}_j+\text{Ric}(\gamma)_{jk}h^k{}_i\Big).
\]
This proves the variation of the Ricci curvature. 
The expansion \eqref{eq:T-expansion} and the schematic structure \eqref{eq:QT-schematic}
follow immediately from \eqref{eq:DT-general}, \eqref{eq:R-linearization-expansion},
\eqref{eq:DRic-general}, and the already established bounds for the scalar-curvature remainder.

\noindent Assume now that $h$ is transverse-traceless relative to $\gamma$, i.e.
\[
\tr_\gamma h=0,\qquad \delta_\gamma h=0.
\]
Then \eqref{eq:DR} reduces to
\[
DR_\gamma[h]=-\langle \text{Ric}(\gamma),h\rangle_\gamma.
\]
Since
\[
\text{Ric}(\gamma)=\mathfrak T[\gamma]+\frac1nR[\gamma]\gamma
\]
and $\langle \gamma,h\rangle_\gamma=\tr_\gamma h=0$, $DR_{\gamma}[h]$ is also equal to
\[
-\langle \mathfrak T[\gamma],h\rangle_\gamma,
\]
proving \eqref{eq:DR-TT}.

\noindent Under the same TT assumptions, \eqref{eq:DRic-general} simplifies to
\[
D\text{Ric}_\gamma[h]_{ij}
=
\frac12\Big(
-\Delta_\gamma h_{ij}
-2\text{Rm}(\gamma)_{ikj\ell}h^{k\ell}
+\text{Ric}(\gamma)_{ik}h^k{}_j+\text{Ric}(\gamma)_{jk}h^k{}_i
\Big)
=
\frac12(\Delta_{L,\gamma}h)_{ij},
\]
which is \eqref{eq:DRic-TT}. Substituting this and \eqref{eq:DR-TT} into \eqref{eq:DT-general}
yields
\[
D\mathfrak T_\gamma[h]_{ij}
=
\frac12(\Delta_{L,\gamma}h)_{ij}
+\frac1n\langle \text{Ric}(\gamma),h\rangle_\gamma\,\gamma_{ij}
-\frac1nR[\gamma]h_{ij}.
\]
Since $h$ is trace-free,
\[
\langle \text{Ric}(\gamma),h\rangle_\gamma
=
\langle \mathfrak T[\gamma],h\rangle_\gamma,
\]
and \eqref{eq:DT-TT} follows. This completes the proof of the lemma.
\end{proof}
\noindent Now we prove the following basic lemma regarding the smooth dependence of the transverse-traceless projection on the metric. This will be important in the next proposition, where we use high frequency $TT-$ eigentensors (of appropriately constructed second order elliptic operator)
\begin{lemma}[TT projection depends smoothly on the metric]
\label{lem:TT-projection}
Fix $s\ge4$ and a background metric $\bar g$. There exists a neighborhood
$\mathcal U\subset \text{Met}^{H^{s}}(M)$ of $\bar g$ and a bounded linear operator
\[
\Pi^{TT}_g: H^{s}(M;S^{2}T^{*}M)\to H^{s}(M;S^{2}T^{*}M),
\qquad g\in\mathcal U,
\]
such that:
\begin{enumerate}
\item $\Pi^{TT}_g$ is a projection: $(\Pi^{TT}_g)^2=\Pi^{TT}_g$.
\item $\Pi^{TT}_g(h)$ is $g$-TT for every $h$.
\item The map $g\mapsto \Pi^{TT}_g$ is $C^\infty$ from $\mathcal U$ into
$\mathcal L(H^{s},H^{s})$.
\item There exists $C=C(\mathcal U)$ such that
\begin{equation}
\label{eq:TTproj-bound}
\|\Pi^{TT}_g(h)\|_{H^{s}} \le C\,\|h\|_{H^{s}}
\quad\text{and}\quad
\|(\Pi^{TT}_g-\Pi^{TT}_{\bar g})(h)\|_{H^{s}} \le C\,\|g-\bar g\|_{H^{s}}\,\|h\|_{H^{s}}.
\end{equation}
\end{enumerate}
\end{lemma}

\begin{proof}
Consider the elliptic operator on $1$--forms
\[
\mathcal L_g := \delta_g\circ \mathcal D_g : H^{s-1}(T^*M)\to H^{s-3}(T^*M),
\]
where $(\mathcal D_g X)_{ij}:=\frac12(\nabla_i X_j+\nabla_j X_i)-\frac13(\div_g X)g_{ij}$
is the conformal Killing operator. For $g$ with no nontrivial conformal Killing fields
(or after restricting to the $L^2$--orthogonal complement of $\ker \mathcal L_g$),
$\mathcal L_g$ is invertible and depends smoothly on $g$; invertibility holds for $g$
in an $H^{s}$--neighborhood of $\bar g$ by stability of elliptic isomorphisms.
Define
\[
\Pi^{TT}_g(h) := h - \mathcal D_g\big(\mathcal L_g^{-1}(\delta_g h)\big)
- \frac13(\tr_g h)\,g.
\]
Then $\Pi^{TT}_g(h)$ is $g$--TT, it is a projection, and the bounds follow from elliptic estimates.
(If $\ker\mathcal L_{\bar g}\neq\{0\}$, impose the standard gauge condition
$\mathcal L_g^{-1}$ on the orthogonal complement; the same estimates hold.)
\end{proof}



\noindent In this last subsection, we formulate the geometric mechanism underlying the construction of the CMC initial data as stated in the main theoerm \ref{main}. The construction consists of three logically separate steps.

\noindent First, starting from a background metric of constant negative scalar curvature, we introduce a high-frequency transverse-traceless perturbation. The linearized trace-free Ricci tensor in such directions is governed by the Lichnerowicz Laplacian and is therefore of second-order size in the oscillation parameter, whereas the scalar curvature variation is only first-order and, in transverse-traceless gauge, is governed by contraction against the background trace-free Ricci tensor. This produces a regime in which the trace-free Ricci tensor is large in \(H^{2}\), while the normalized scalar curvature defect remains small in \(H^{2}\).

\noindent Second, on this manifold $(M,g)$, our task is to now find a $TT$ tensor $\Sigma$ that verifies the estimates 
\begin{align}
||\Sigma||_{H^{2}}=O(e^{-9a/10}),~||\Sigma||_{H^{3}}=O(e^{a/10})
\end{align}
so that we can use such $\Sigma$ as the free data for our constraint system that is to be solved via conformal method in the next step. In particular in the next (and third) step we perform the conformal transformation 
\[
g_{0}=\varphi^{\frac{4}{n-2}}g,
\qquad
\Sigma_{0}=\varphi^{-2}\Sigma.
\]
where we intend to use $\Sigma$ as the data for the momentum constraint 
\begin{align}
\label{eq:TT1}
\div_{g_{0}}\Sigma_{0}=0.
\end{align}
By covariance the $TT$ property of $\Sigma$ with respect to $g$ i.e., $\tr_{g}\Sigma=0$ and $\div_{g}\Sigma=0$ implies $\tr_{g_{0}}\Sigma_{0}=0$ and the momentum constraint (\ref{eq:TT1}). Therefore it is crucial to explicitly construct such $\Sigma$ on $(M,g)$. 

\noindent Third, we correct the scalar curvature defect by solving the Hamiltonian constraint through a conformal deformation. The resulting Lichnerowicz equation is solved perturbatively around the constant solution \(1\), and the resulting conformal factor is shown to be so close to \(1\) in $H^{4}$ Sobolev norm that the trace-free Ricci tensor changes only by a lower-order amount in \(H^{2}\).




\noindent We now record the quantitative form of the three-step construction described heuristically in the introduction. Throughout this subsection, \(M\) denotes a closed connected smooth \(3\)-manifold of negative Yamabe type. We fix once and for all a smooth metric \(\gamma\) belonging to a conformal class $[\widehat{\gamma}]$ satisfying
\[
R[\gamma]=-\frac23.
\]
All covariant derivatives, contractions, volume forms, and Sobolev norms are taken with respect to \(\gamma\), unless explicitly indicated otherwise.



\begin{proposition}
\label{prop:YorkTTperturb}
Let $(M^{3},\widehat{\gamma})$ be a smooth closed manifold of negative Yamabe type. Choose a smooth metric $\gamma\in[\widehat{\gamma}]$ such that
\[
R[\gamma]\equiv -\frac23.
\]
Assume, for simplicity, that $\gamma$ admits no nontrivial conformal Killing fields. Define the conformal Killing operator on $1$-forms by
\[
(\mathcal D_\gamma W)_{ij}
:=
\nabla_i W_j+\nabla_j W_i-\frac23(\nabla^kW_k)\gamma_{ij},
\]
and let $S^2_0T^*M$ denote the bundle of $\gamma$-trace-free symmetric $2$-tensors. Since $\delta_\gamma\mathcal D_\gamma$ is then invertible, the York projector
\[
\Pi_{TT}
:=
I-\mathcal D_\gamma(\delta_\gamma\mathcal D_\gamma)^{-1}\delta_\gamma
\]
is a well-defined bounded operator on $H^s(S^2_0T^*M)$ for every integer $s\ge0$, with range
\[
TT_\gamma
:=
\{h\in C^\infty(S^{2}T^{*}M):\tr_\gamma h=0,\ \delta_\gamma h=0\}.
\]
Consider the positive self-adjoint elliptic pseudodifferential operator of order two
\[
A_\gamma
:=
\Pi_{TT}(I+\nabla_\gamma^{*}\nabla_\gamma)\Pi_{TT}\big|_{L^2(TT_\gamma)}
\]
acting on $L^2(TT_\gamma)$.

\noindent Then there exist constants
\[
\varepsilon_{*}>0,\qquad C_m\ge1\ \ (m\ge0),\qquad c_m>0\ \ (m\ge0),
\]
depending only on $(M,\gamma)$, together with an $L^2(\gamma)$-orthonormal sequence of smooth tensors
\[
h_\nu\in TT_\gamma,\qquad \nu\in\mathbb N,
\]
and a sequence of positive numbers
\[
\Lambda_\nu\to+\infty,
\]
such that
\[
A_\gamma h_\nu=\Lambda_\nu h_\nu,
\]
and, for every integer $m\ge0$,
\[
\|h_\nu\|_{H^m(\gamma)}
\le C_m(1+\Lambda_\nu)^{m/2}
\]
for all $\nu$, while
\[
\|h_\nu\|_{H^m(\gamma)}
\ge c_m\Lambda_\nu^{m/2}
\]
for all sufficiently large $\nu$.

\noindent Moreover, if
\[
g_{\varepsilon,\nu}:=\gamma+\varepsilon h_\nu,
\qquad
0<\varepsilon\le \varepsilon_{*}\Lambda_\nu^{-3/2},
\]
then $g_{\varepsilon,\nu}$ is a smooth Riemannian metric and
\begin{align}
\label{eq:Tquantitative}
\|\mathfrak T[g_{\varepsilon,\nu}]\|_{H^2(\gamma)}
&\ge
c_{4}\,\varepsilon \Lambda_\nu^{2}
-
C\,\varepsilon^{2}\Lambda_\nu^{3}
-
C,\\
\label{eq:Rquantitative}
\left\|R[g_{\varepsilon,\nu}]+\frac23\right\|_{H^2(\gamma)}
&\le
C\,\varepsilon \Lambda_\nu
+
C\,\varepsilon^{2}\Lambda_\nu^{3},
\end{align}
where
\[
\mathfrak T[g]:=\text{Ric}[g]-\frac13R[g]\,g.
\]
\end{proposition}

\begin{proof} The basic idea is to perturb the metric along transverse-traceless ($TT$) direction. In addition, the frequency content of such perturbations is restricted to be very large. Such perturbations can modify the trace-free Ricci curvature $\mathfrak{T}[g]:=\text{Ric}[g]-\frac{1}{3}R[g]g$ in high enough Sobolev norm while keeping the scalar curvature almost invariant. 
The main idea is to construct a second order pseudo-differential operator (in the sense of H\"ordmander) with smooth $L^{2}$ normalized $TT$ family of eigentensors.  We divide the argument into several steps. 
Let
\[
L_\gamma:=\delta_\gamma\mathcal D_\gamma
\]
be the vector Laplacian associated with $\gamma$. Since $\gamma$ has no nontrivial conformal Killing fields, the kernel of $L_\gamma$ is trivial. Indeed, by the standard adjoint relation between $\delta_\gamma$ and $\mathcal D_\gamma$, one has by an exact similar calculation as in the proof of claim (\ref{important})
\[
\langle L_\gamma W,W\rangle_{L^2(\gamma)}
=
\frac12\|\mathcal D_\gamma W\|_{L^2(\gamma)}^{2},
\]
up to the harmless sign dictated by the convention for $\delta_\gamma$; in particular,
\[
L_\gamma W=0
\quad\Longrightarrow\quad
\mathcal D_\gamma W=0.
\]
Thus the assumption on conformal Killing fields implies $\ker L_\gamma=\{0\}$. Since $L_\gamma$ is a second-order strongly elliptic self-adjoint operator on the closed manifold $M$, elliptic Fredholm theory yields that
\[
L_\gamma:H^{s+2}(T^*M)\to H^s(T^*M)
\]
is an isomorphism for every integer $s\ge0$.

\noindent It follows that the operator
\[
\Pi_{TT}
=
I-\mathcal D_\gamma L_\gamma^{-1}\delta_\gamma
\]
is well defined and bounded on $H^s(S^2_0T^*M)$ for every $s\ge0$. We now verify that its range is precisely the TT space. Let $k\in H^s(S^2_0T^*M)$ and define
\[
W:=L_\gamma^{-1}\delta_\gamma k,\qquad k^{TT}:=k-\mathcal D_\gamma W.
\]
Since $\mathcal D_\gamma W$ is trace-free by construction, $k^{TT}$ is trace-free as well. Moreover,
\[
\delta_\gamma k^{TT}
=
\delta_\gamma k-\delta_\gamma\mathcal D_\gamma W
=
\delta_\gamma k-L_\gamma W
=
0.
\]
Hence $k^{TT}\in TT_\gamma$. Conversely, if $h\in TT_\gamma$, then $\delta_\gamma h=0$, so
\[
\Pi_{TT}h
=
h-\mathcal D_\gamma L_\gamma^{-1}\delta_\gamma h
=
h.
\]
Therefore $\Pi_{TT}$ is a projection with range $TT_\gamma$.

\noindent In particular, for each $s\ge0$ one has the topological direct sum decomposition
\[
H^s(S^2_0T^*M)
=
H^s(TT_\gamma)\oplus \mathcal D_\gamma H^{s+1}(T^*M).
\]
This is the usual York splitting \cite{york}. Since the principal symbol of $\Pi_{TT}$ is the orthogonal projection onto the algebraic subspace
\[
\{q\in S^2_0T_x^*M:\ \xi^i q_{ij}=0\},
\]
whose rank equals $2$ in dimension $3$, the range of $\Pi_{TT}$ is infinite-dimensional. Thus $TT_\gamma$ is an infinite-dimensional closed subspace of $L^2(S^2_0T^*M)$.

\smallskip
\noindent
Next we consider the high-frequency TT basis that we are going to use to perturb the metric.
We now consider
\[
A_\gamma
=
\Pi_{TT}(I+\nabla_\gamma^{*}\nabla_\gamma)\Pi_{TT}\big|_{L^2(TT_\gamma)}.
\]
Because $\Pi_{TT}$ is a self-adjoint pseudodifferential operator of order $0$ and $I+\nabla_\gamma^*\nabla_\gamma$ is a positive self-adjoint elliptic differential operator of order $2$, the operator $A_\gamma$ is a positive self-adjoint elliptic pseudodifferential operator of order $2$ on the Hilbert space $L^2(TT_\gamma)$. Its principal symbol on the TT symbol bundle is
\[
\sigma_2(A_\gamma)(x,\xi)
=
|\xi|_\gamma^2\,\text{Id} .
\]
In particular, $A_\gamma$ is elliptic and has a compact resolvent. Standard spectral theory for positive self-adjoint elliptic operators on closed manifolds therefore, yields an $L^2(\gamma)$-orthonormal basis of smooth eigentensors
\[
\{h_\nu\}_{\nu\ge1}\subset TT_\gamma
\]
with corresponding eigenvalues
\[
0<\Lambda_1\le\Lambda_2\le\cdots,\qquad \Lambda_\nu\to+\infty,
\]
such that
\[
A_\gamma h_\nu=\Lambda_\nu h_\nu.
\]
 We next record the Sobolev bounds for this eigenbasis. Since $A_\gamma$ is elliptic of order $2$, the graph norm of $(I+A_\gamma)^{m/2}$ is equivalent to the $H^m$ norm on $TT_\gamma$; that is, for each integer $m\ge0$ there exist constants $c_m',C_m'>0$ such that
\[
c_m'\|u\|_{H^m(\gamma)}
\le
\|(I+A_\gamma)^{m/2}u\|_{L^2(\gamma)}
\le
C_m'\|u\|_{H^m(\gamma)}
\]
for all $u\in H^m(TT_\gamma)$. Applying this to $u=h_\nu$ and using
\[
(I+A_\gamma)^{m/2}h_\nu=(1+\Lambda_\nu)^{m/2}h_\nu,
\qquad
\|h_\nu\|_{L^2(\gamma)}=1,
\]
we obtain
\[
c_m'(1+\Lambda_\nu)^{m/2}
\le
\|h_\nu\|_{H^m(\gamma)}
\le
C_m'(1+\Lambda_\nu)^{m/2}.
\]
After adjusting the constants, this implies
\[
\|h_\nu\|_{H^m(\gamma)}
\le C_m(1+\Lambda_\nu)^{m/2}
\]
for all $\nu$, and
\[
\|h_\nu\|_{H^m(\gamma)}
\ge c_m\Lambda_\nu^{m/2}
\]
for all sufficiently large $\nu$.
\noindent In particular,
\[
\|h_\nu\|_{H^2(\gamma)}\le C_2\Lambda_\nu,
\qquad
\|h_\nu\|_{H^4(\gamma)}\ge c_4\Lambda_\nu^2
\]
for all sufficiently large $\nu$.

\smallskip
\noindent
Now we show that the $TT$ perturbation that we apply to the metric preseves its Riemannian character. 
Since $\dim M=3$, the Sobolev embedding $H^{2}(M)\hookrightarrow C^0(M)$ gives
\[
\|h_\nu\|_{C^0(\gamma)}
\le
C\|h_\nu\|_{H^2(\gamma)}
\le
C\Lambda_\nu.
\]
Hence, if
\[
0<\varepsilon\le \varepsilon_*\Lambda_\nu^{-3/2},
\]
then
\[
\varepsilon\|h_\nu\|_{C^0(\gamma)}
\le
C\varepsilon_*\Lambda_\nu^{-1/2}.
\]
By choosing $\varepsilon_*>0$ sufficiently small, and then decreasing it once more to absorb the finitely many low-frequency modes, we may ensure that
\[
\varepsilon\|h_\nu\|_{C^0(\gamma)}\le \frac12
\]
for every $\nu$ under the above restriction on $\varepsilon$. Now let $X\in T_xM$. Then
\[
g_{\varepsilon,\nu}(X,X)
=
\gamma(X,X)+\varepsilon h_\nu(X,X)
\ge
\Bigl(1-\varepsilon\|h_\nu\|_{C^0(\gamma)}\Bigr)\gamma(X,X)
\ge
\frac12\,\gamma(X,X).
\]
Therefore, $g_{\varepsilon,\nu}$ is positive definite, hence a smooth Riemannian metric.

\noindent Now using the previous lemma \ref{nonlinearexpansion}, we obtain the necessary estimates for the trace-free Ricci curvature and the scalar curvature of the perturbed metric.
We write
\[
\mathfrak T[g]=\text{Ric}[g]-\frac13R[g]\,g.
\]
The map
\[
g\longmapsto \mathfrak T[g]
\]
is a smooth quasilinear differential operator of order two. Accordingly, for $k$ small in $H^4$, one has the Taylor expansion
\[
\mathfrak T[\gamma+k]
=
\mathfrak T[\gamma]
+
D\mathfrak T_\gamma[k]
+
\mathcal R_{\mathfrak T,\gamma}(k),
\]
where the remainder is quadratic in $k$ and its derivatives up to order two. More precisely, because curvature depends smoothly on the inverse metric and on the first and second derivatives of the metric, and because $H^2(M)$ is a Banach algebra in dimension $3$, there exists a neighborhood of the origin in $H^4(S^2T^*M)$ and a constant $C$ such that
\begin{equation}
\label{eq:RTremainder}
\|\mathcal R_{\mathfrak T,\gamma}(k)\|_{H^2(\gamma)}
\le
C\|k\|_{H^4(\gamma)}\|k\|_{H^2(\gamma)}
\end{equation}
for all $k$ in that neighborhood.

\noindent We now restrict to TT directions. Set
\[
L_\gamma:=D\mathfrak T_\gamma\big|_{TT_\gamma}.
\]
Since the principal second-order part of the linearized Ricci tensor is $-\frac12\nabla^*\nabla h$, and since the additional contributions coming from the scalar-curvature term are of lower order on TT tensors, the principal symbol of $L_\gamma$ is
\[
\sigma_2(L_\gamma)(x,\xi)
=
-\frac12|\xi|_\gamma^2\,\text{Id}
\]
on the TT symbol space. Thus $L_\gamma$ is elliptic on $TT_\gamma$.

\noindent Therefore elliptic regularity gives, for $u\in H^4(TT_\gamma)$,
\[
\|u\|_{H^4(\gamma)}
\le
C\bigl(\|L_\gamma u\|_{H^2(\gamma)}+\|u\|_{L^2(\gamma)}\bigr).
\]
Applying this estimate to $u=h_\nu$ and using $\|h_\nu\|_{L^2}=1$, we obtain
\[
\|L_\gamma h_\nu\|_{H^2(\gamma)}
\ge
C^{-1}\|h_\nu\|_{H^4(\gamma)}-C
\ge
c\,\Lambda_\nu^2-C.
\]
After decreasing $c$ if necessary, we may rewrite this as
\begin{equation}
\label{eq:Lgamma-lower}
\|L_\gamma h_\nu\|_{H^2(\gamma)}
\ge
c_4\Lambda_\nu^2-C
\end{equation}
for all $\nu$.

\noindent Now substitute $k=\varepsilon h_\nu$ into the Taylor formula. Since $h_\nu\in TT_\gamma$,
\[
\mathfrak T[g_{\varepsilon,\nu}]
=
\mathfrak T[\gamma]
+
\varepsilon L_\gamma h_\nu
+
\mathcal R_{\mathfrak T,\gamma}(\varepsilon h_\nu).
\]
Using \eqref{eq:RTremainder}, together with the Sobolev bounds for $h_\nu$, we find
\[
\|\mathcal R_{\mathfrak T,\gamma}(\varepsilon h_\nu)\|_{H^2(\gamma)}
\le
C\varepsilon^2
\|h_\nu\|_{H^4(\gamma)}
\|h_\nu\|_{H^2(\gamma)}
\le
C\varepsilon^2\Lambda_\nu^3.
\]
Hence, by the reverse triangle inequality,
\begin{align*}
\|\mathfrak T[g_{\varepsilon,\nu}]\|_{H^2(\gamma)}
&\ge
\varepsilon\|L_\gamma h_\nu\|_{H^2(\gamma)}
-
\|\mathfrak T[\gamma]\|_{H^2(\gamma)}
-
\|\mathcal R_{\mathfrak T,\gamma}(\varepsilon h_\nu)\|_{H^2(\gamma)}\\
&\ge
\varepsilon(c_4\Lambda_\nu^2-C)
-
C
-
C\varepsilon^2\Lambda_\nu^3.
\end{align*}
Since $\varepsilon\le1$ after decreasing $\varepsilon_*$ once more if necessary, the term $\varepsilon C$ may be absorbed into the final background constant, and we conclude that
\[
\|\mathfrak T[g_{\varepsilon,\nu}]\|_{H^2(\gamma)}
\ge
c_4\,\varepsilon\Lambda_\nu^2
-
C\varepsilon^2\Lambda_\nu^3
-
C.
\]
This proves \eqref{eq:Tquantitative}.

\smallskip
\noindent
Now we focus on the scalar curvature expansion.
The scalar curvature map is likewise a smooth quasilinear differential operator of order two, so one has
\[
R[\gamma+k]
=
R[\gamma]
+
DR_\gamma[k]
+
\mathcal R_{R,\gamma}(k),
\]
with
\begin{equation}
\label{eq:Rremainder}
\|\mathcal R_{R,\gamma}(k)\|_{H^2(\gamma)}
\le
C\|k\|_{H^4(\gamma)}\|k\|_{H^2(\gamma)}.
\end{equation}
Since $R[\gamma]\equiv-\frac23$, it remains to estimate the linear term. The standard formula for the linearization of scalar curvature is
\[
DR_\gamma[h]
=
-\Delta_\gamma(\tr_\gamma h)+\delta_\gamma\delta_\gamma h-\langle \text{Ric}[\gamma],h\rangle_\gamma.
\]
For $h=h_\nu\in TT_\gamma$, the first two terms vanish, and therefore
\[
DR_\gamma[h_\nu]
=
-\langle \text{Ric}[\gamma],h_\nu\rangle_\gamma.
\]
This is the main mechanism that cancels the principal part in the expansion of the scalar curvature and therefore does not cost frequency.
Because $\text{Ric}[\gamma]$ is a fixed smooth tensor, multiplication by $\text{Ric}[\gamma]$ is a bounded operator on $H^2$, and hence
\[
\|DR_\gamma[h_\nu]\|_{H^2(\gamma)}
\le
C\|h_\nu\|_{H^2(\gamma)}
\le
C\Lambda_\nu.
\]

\noindent Substituting $k=\varepsilon h_\nu$ into the expansion for $R$, and using \eqref{eq:Rremainder}, we obtain
\begin{align*}
\left\|R[g_{\varepsilon,\nu}]+\frac23\right\|_{H^2(\gamma)}
&\le
\varepsilon\|DR_\gamma[h_\nu]\|_{H^2(\gamma)}
+
\|\mathcal R_{R,\gamma}(\varepsilon h_\nu)\|_{H^2(\gamma)}\\
&\le
C\varepsilon\Lambda_\nu
+
C\varepsilon^2
\|h_\nu\|_{H^4(\gamma)}
\|h_\nu\|_{H^2(\gamma)}\\
&\le
C\varepsilon\Lambda_\nu
+
C\varepsilon^2\Lambda_\nu^3.
\end{align*}
This is exactly \eqref{eq:Rquantitative}. The proof is complete.
\end{proof}

\begin{remark}
Note that instead of a general negative Yamabe background, if it is Einstein, then one can simply use the spectrum of the Lichnerowicz Laplacian $\Delta_{L}$ since that preserves the transverse-traceless (TT) subspace of $S^{2}T^{*}M$. On a general closed manifold, this is not true and one ought use the TT projector machinery used in the proposition     
\end{remark}


\noindent Now in the corollary, we choose the smallness parameter $\epsilon$ appearing in the previous proposition \ref{prop:YorkTTperturb} and the Lichnerowicz eigenvalue $\lambda_{\nu}$ in terms of the largeness parameter $a$ in our study. This provides the estimates necessary in this context.
\begin{corollary}
\label{cor:separation-scales}
There exist constants \(a_{0}\geq 1\), \(c>0\), and \(C\geq 1\), depending only on \((M,\gamma)\), such that for every \(a\geq a_{0}\) there exist an index \(\nu(a)\) and a parameter \(\varepsilon(a)>0\) for which the metric
\[
g_{1,a}:=\gamma+\varepsilon(a) h_{\nu(a)}
\]
satisfies
\begin{align}
\label{eq:Tlargecor}
\|\mathfrak T[g_{1,a}]\|_{H^{2}(\gamma)}&=O(e^{a/10}),\\
\label{eq:Rsmallcor}
\left\|R[g_{1,a}]+\frac23\right\|_{H^{2}(\gamma)}&= O(e^{-77a/20}).
\end{align}
\end{corollary}

\begin{proof}
First, assume that $||\mathfrak{T}[\gamma]||_{H^{2}}\leq C$, where $C\in [0,O(e^{\frac{a}{10}-\alpha})]$ for $\alpha>0$ since if $||\mathfrak{T}[\gamma]||_{H^{2}}=O(e^{a/10})$, then there is nothing to prove. In particular, if $M$ admits an Einstein metric, then $C=0$ and in that case $\mathfrak{T}$ vanishes identically. Therefore, the data is large is in the sense that its large deviation from the Einstein background (if it exists on $M$). 
Fix
\[
\varepsilon(a):=e^{-8a}.
\]
Since \(\lambda_{\nu}\to +\infty\), for every sufficiently large \(a\) one may choose \(\nu(a)\) such that
\[
e^{\frac{81a}{20}}\leq \lambda_{\nu(a)}\leq 2e^{\frac{81a}{20}}.
\]
Then
\[
\varepsilon(a)\lambda_{\nu(a)}^{3/2}
\leq
Ce^{-8a}e^{\frac{243a}{40}}
=
Ce^{-77a/40},
\]
and so for \(a\) sufficiently large,
\[
\varepsilon(a)\leq \varepsilon_{\ast}\lambda_{\nu(a)}^{-3/2},
\]
and proposition ~\ref{prop:YorkTTperturb} applies.
\noindent Next we note, 
\[
\varepsilon(a)\lambda_{\nu(a)}^{2}
\geq
c\,e^{-8a}e^{\frac{81a}{10}}
=
c\,e^{\frac{a}{10}},
\]
while
\[
\varepsilon(a)\lambda_{\nu(a)}
\leq
Ce^{-8a}e^{\frac{81a}{20}}
=
Ce^{-79a/20},
\]
and
\[
\varepsilon(a)^{2}\lambda_{\nu(a)}^{3}
\leq
Ce^{-16a}e^{\frac{243a}{20}}
=
Ce^{-\frac{77a}{20}}.
\]
Substituting these bounds into the estimates obtained in the previous proposition \ref{prop:YorkTTperturb} i.e., into \eqref{eq:Tquantitative} and \eqref{eq:Rquantitative}, we obtain
\[
\|\mathfrak T[g_{1,a}]\|_{H^{2}(\gamma)}
\geq
c\,e^{\frac{a}{10}}
-
Ce^{-\frac{77a}{20}}
-
C,
\]
and
\[
\left\|R[g_{1,a}]+\frac23\right\|_{H^{2}(\gamma)}
\leq
Ce^{-\frac{79a}{20}}
+
Ce^{-\frac{77a}{20}}.
\]
For \(a\) sufficiently large, the first estimate implies
\[
\|\mathfrak T[g_{1,a}]\|_{H^{2}(\gamma)}\geq c\,e^{a/10},
\]
and the second implies
\[
\left\|R[g_{1,a}]+\frac23\right\|_{H^{2}(\gamma)}\leq C\,e^{-\frac{77a}{20}}<ce^{-2a}.
\]
This proves the corollary.
\end{proof}
\noindent Now we want to construct the final physical metric $g_{0}$ and the transverse-traceless second fundamental form $\Sigma_{0}$ that verify the constraint equations 
\begin{align}
R[g_{0,a}]+\frac23&=|\Sigma_{0,a}|_{g_{0,a}}^{2},\\
\div_{g_{0,a}}\Sigma_{0,a}&=0.
\end{align}
To this end, we perform the conformal transformation 
\begin{align}
 g_{0}=\varphi^{\frac{4}{n-2}}g,
\qquad
\Sigma_{0}=\varphi^{-2}\Sigma,   
\end{align}
where $g$ is chosen to be $g_{1,a}$ and $\Sigma$ is the free data which is transverse-traceless with respect to $g$ (as is standard in the conformal technique) that verifies 
\begin{align}
\label{eq:estimate}
||\Sigma||_{H^{2}}=O(e^{-9a/10}),~||\Sigma||_{H^{3}}=O(e^{a/10}).
\end{align}
Now question may arise: how to construct such a $TT$ tensor on $(M,g)$ that verifies the estimate (\ref{eq:estimate}). We do this now in the following proposition.
\begin{proposition}
\label{prop:TT_construction_separated_scales}
Let $M$ be a smooth closed manifold of dimension $n\geq 3$, and let $g$ be a smooth Riemannian metric on $M$. In particular, the conclusion applies when $M$ is of negative Yamabe type.

\noindent Then there exist constants
\[
a_{0}\geq 1,
\qquad
0<c\leq C<\infty,
\]
depending only on $(M,g)$, and for each $a\geq a_{0}$ a smooth symmetric $2$-tensor $\Sigma_{a}$ on $M$ such that
\[
\tr_{g}\Sigma_{a}=0,
\qquad
\div_{g}\Sigma_{a}=0,
\]
and
\[
c\,e^{-9a/10}\leq \|\Sigma_{a}\|_{H^{2}(M,g)}\leq C\,e^{-9a/10},
\qquad
c\,e^{a/10}\leq \|\Sigma_{a}\|_{H^{3}(M,g)}\leq C\,e^{a/10}.
\]
In particular,
\[
\|\Sigma_{a}\|_{H^{2}(M,g)}=O(e^{-9a/10}),
\qquad
\|\Sigma_{a}\|_{H^{3}(M,g)}=O(e^{a/10}),
\]
and the family $\Sigma_{a}$ is transverse-traceless with respect to $g$.
\end{proposition}

\begin{proof}
We divide the construction into three steps.

\medskip
\noindent  First we construct a high-frequency trace-free seed with one-derivative smaller divergence.
Choose a coordinate chart $U\subset M$ and smooth coordinates
\[
x=(x^{1},\dots,x^{n}):U\to B(0,2)\subset \mathbb R^{n}.
\]
Let
\[
\psi:=x^{1}\in C^{\infty}(U).
\]
Shrinking $U$ if necessary, we may assume that $d\psi$ is nowhere vanishing on $U$. Define the smooth unit vector field
\[
E_{1}:=\frac{\nabla \psi}{|\nabla \psi|_{g}}
\]
on $U$, and extend $E_{1}$ to a smooth local $g$-orthonormal frame
\[
E_{1},E_{2},E_{3},\dots,E_{n}
\]
on $U$. Let
\[
\theta^{1},\theta^{2},\theta^{3},\dots,\theta^{n}
\]
denote the dual coframe. Define the smooth symmetric $2$-tensor
\[
q:=\theta^{2}\otimes \theta^{2}-\theta^{3}\otimes \theta^{3}
\]
on $U$. Then note 
\[
\tr_{g} q = 0
\]
and, since $q(E_{1},\cdot)=0$, also
\[
q(\nabla \psi,\cdot)=0
\qquad\text{on }U.
\]

\noindent Fix a cutoff $\chi\in C_{c}^{\infty}(U)$ such that $0\leq \chi\leq 1$ and $\chi\equiv 1$ on some nonempty open set $V\Subset U$. For each parameter $\mu\geq 1$, define the modulation
\[
S_{\mu}:=\chi \cos(\mu\psi)\, q,
\]
extended by zero outside $U$. Then $S_{\mu}$ is a smooth symmetric $2$-tensor on $M$ and, because $q$ is trace-free,
\[
\tr_{g} S_{\mu}=0.
\]

\noindent We first record the $H^{m}$-size of $S_{\mu}$ for $m=2,3$.

\smallskip
\noindent Since $\chi$ and $q$ are fixed smooth tensors supported in $U$, repeated differentiation of $\cos(\mu\psi)$ yields
\[
\nabla^{j}\big(\cos(\mu\psi)\big)
=
\sum_{\ell=1}^{j} \mu^{\ell}\,A_{j,\ell},
\]
where each $A_{j,\ell}$ is a smooth tensor depending only on $g$, $\psi$, and finitely many derivatives thereof, multiplied by either $\sin(\mu\psi)$ or $\cos(\mu\psi)$. By Leibniz' rule, for each integer $j\geq 0$ there exists a constant $C_{j}$, depending only on $(M,g)$ and the auxiliary choices $\psi,\chi,q$, such that
\[
\|\nabla^{j} S_{\mu}\|_{L^{2}(M,g)}
\leq C_{j}\,\mu^{j}.
\]
Consequently, for $m=2,3$,
\begin{equation}
\label{eq:seed_upper}
\|S_{\mu}\|_{H^{m}(M,g)}\leq C\,\mu^{m}.
\end{equation}

\smallskip
\noindent Now consider the following. On the open set $V$ we have $\chi\equiv 1$, and hence
\[
S_{\mu}=\cos(\mu\psi)\,q
\qquad\text{on }V.
\]
Applying the covariant derivative repeatedly in the $E_{1}$-direction gives
\[
\nabla_{E_{1}}^{m} S_{\mu}
=
\mu^{m}\,b_{m}(\mu\psi)\,\big(E_{1}\psi\big)^{m}q
+
\mu^{m-1}R_{m,\mu},
\qquad m=2,3,
\]
where $b_{m}$ is either $\sin$ or $\cos$, and where the remainder $R_{m,\mu}$ is a smooth tensor field supported in $V$ satisfying
\[
\|R_{m,\mu}\|_{L^{2}(V,g)}\leq C.
\]
Since
\[
E_{1}\psi = |\nabla \psi|_{g}>0
\quad\text{on }V,
\]
and $q$ is nowhere zero on $V$, there exists $c_{0}>0$ such that
\[
\left\| b_{m}(\mu\psi)\,\big(E_{1}\psi\big)^{m}q\right\|_{L^{2}(V,g)}\geq c_{0}
\]
uniformly for all $\mu\geq 1$. It follows that
\[
\|\nabla^{m}S_{\mu}\|_{L^{2}(M,g)}
\geq
\|\nabla_{E_{1}}^{m}S_{\mu}\|_{L^{2}(V,g)}
\geq c_{1}\mu^{m}-C\mu^{m-1}.
\]
Hence, after increasing $\mu$ if necessary,
\[
\|S_{\mu}\|_{H^{m}(M,g)}\geq c\,\mu^{m},
\qquad m=2,3.
\]
Together with \eqref{eq:seed_upper}, this yields
\begin{equation}
\label{eq:seed_two_sided}
c\,\mu^{m}\leq \|S_{\mu}\|_{H^{m}(M,g)}\leq C\,\mu^{m},
\qquad m=2,3,
\end{equation}
for all sufficiently large $\mu$.


\noindent We next estimate the divergence of $S_{\mu}$. We use the convention
\[
(\div_{g} h)_{j}:=\nabla^{i}h_{ij}.
\]
A direct computation gives
\[
(\div_{g}S_{\mu})_{j}
=
\nabla^{i}\chi \,\cos(\mu\psi)\, q_{ij}
-\mu \chi \sin(\mu\psi)\,\nabla^{i}\psi\, q_{ij}
+\chi \cos(\mu\psi)\,\nabla^{i}q_{ij}.
\]
The middle term vanishes identically because $q(\nabla \psi,\cdot)=0$. Therefore
\[
\div_{g}S_{\mu}
=
(\nabla \chi)*q*\cos(\mu\psi)
+
\chi\,(\nabla q)*\cos(\mu\psi),
\]
where $*$ denotes a universal contraction. In particular, no term of size $\mu$ appears in the amplitude of $\div_{g}S_{\mu}$. This is precisely the reason of choosing $q$ such that $q(\nabla\psi,\cdot)=0$. Differentiating this identity and arguing as above, we obtain for $m=2,3$,
\begin{equation}
\label{eq:div_seed}
\|\div_{g}S_{\mu}\|_{H^{m-1}(M,g)}\leq C\,\mu^{m-1}.
\end{equation}

\noindent Now we explicitly construct a $TT$ tensor through York-projection \cite{york}.
Let
\[
(\mathbb L_{g}W)_{ij}
:=
\nabla_{i}W_{j}
+
\nabla_{j}W_{i}
-
\frac{2}{n}(\div_{g}W)\,g_{ij}
\]
and denote the conformal Killing operator acting on one-forms $W$. By construction,
\[
\tr_{g}(\mathbb L_{g}W)=0
\]
for every one-form $W$. Define the vector Laplacian
\[
\Delta_{\mathbb L}W:=-\div_{g}(\mathbb L_{g}W).
\]
This is a second-order self-adjoint elliptic operator on one-forms. Moreover, by a similar computation as the previous case
\[
\int_{M}\langle \Delta_{\mathbb L}W,W\rangle\,d\mu_{g}
=
\frac12\int_{M} |\mathbb L_{g}W|_{g}^{2}\,d\mu_{g},
\]
and hence
\[
\ker(\Delta_{\mathbb L})=\ker(\mathbb L_{g}),
\]
the finite-dimensional space of conformal Killing one-forms.



\noindent We claim that $\div_{g}S_{\mu}$ is $L^{2}$-orthogonal to $\ker(\Delta_{\mathbb L})$. Indeed, let $X\in \ker(\Delta_{\mathbb L})=\ker(\mathbb L_{g})$. Since $S_{\mu}$ is trace-free, the formal adjoint relation gives
\[
\int_{M}\langle \div_{g}S_{\mu},X\rangle\,d\mu_{g}
=
-\frac12\int_{M}\langle S_{\mu},\mathbb L_{g}X\rangle\,d\mu_{g}
=0.
\]
Therefore $\div_{g}S_{\mu}$ belongs to the range of $\Delta_{\mathbb L}$, and there exists a unique one-form $W_{\mu}$ satisfying
\[
W_{\mu}\perp \ker(\Delta_{\mathbb L}) \quad\text{in }L^{2}(M,g),
\]
and
\begin{equation}
\label{eq:York_eqn}
\Delta_{\mathbb L}W_{\mu}=\div_{g}S_{\mu}.
\end{equation}
Define
\[
\tau_{\mu}:=S_{\mu}+\mathbb L_{g}W_{\mu}.
\]
Since both $S_{\mu}$ and $\mathbb L_{g}W_{\mu}$ are trace-free, we have
\[
\tr_{g}\tau_{\mu}=0.
\]
Also, by \eqref{eq:York_eqn},
\[
\div_{g}\tau_{\mu}
=
\div_{g}S_{\mu}+\div_{g}(\mathbb L_{g}W_{\mu})
=
\div_{g}S_{\mu}-\Delta_{\mathbb L}W_{\mu}
=0.
\]
Thus $\tau_{\mu}$ is transverse-traceless.

\smallskip
\noindent We now estimate the correction term. Since $\Delta_{\mathbb L}$ is elliptic and invertible on the $L^{2}$-orthogonal complement of its kernel, standard elliptic regularity yields, for $m=2,3$,
\begin{equation}
\label{eq:elliptic_W}
\|W_{\mu}\|_{H^{m+1}(M,g)}
\leq
C\,\|\div_{g}S_{\mu}\|_{H^{m-1}(M,g)}.
\end{equation}
Combining \eqref{eq:elliptic_W} with \eqref{eq:div_seed}, we obtain
\[
\|W_{\mu}\|_{H^{m+1}(M,g)}\leq C\,\mu^{m-1},
\qquad m=2,3.
\]
Since $\mathbb L_{g}$ is first order,
\begin{equation}
\label{eq:LW_est}
\|\mathbb L_{g}W_{\mu}\|_{H^{m}(M,g)}
\leq
C\,\|W_{\mu}\|_{H^{m+1}(M,g)}
\leq C\,\mu^{m-1},
\qquad m=2,3.
\end{equation}
\noindent Now we obtain the two-sided bounds for the TT tensor by the choice of an appropriate scale. By definition,
\[
\tau_{\mu}=S_{\mu}+\mathbb L_{g}W_{\mu}.
\]
From \eqref{eq:seed_two_sided} and \eqref{eq:LW_est}, for $m=2,3$ we have
\[
\|\tau_{\mu}\|_{H^{m}}
\leq
\|S_{\mu}\|_{H^{m}}+\|\mathbb L_{g}W_{\mu}\|_{H^{m}}
\leq C\mu^{m},
\]
and also
\[
\|\tau_{\mu}\|_{H^{m}}
\geq
\|S_{\mu}\|_{H^{m}}-\|\mathbb L_{g}W_{\mu}\|_{H^{m}}
\geq
c\mu^{m}-C\mu^{m-1}.
\]
Therefore, for all sufficiently large $\mu$,
\begin{equation}
\label{eq:tau_two_sided}
c\,\mu^{m}\leq \|\tau_{\mu}\|_{H^{m}(M,g)}\leq C\,\mu^{m},
\qquad m=2,3.
\end{equation}

\noindent Now let
\[
\mu_{a}:=e^{a},
\qquad
A_{a}:=e^{-29a/10},
\]
and define
\[
\Sigma_{a}:=A_{a}\,\tau_{\mu_{a}}
=
e^{-29a/10}\tau_{e^{a}}.
\]
Since $\tau_{e^{a}}$ is transverse-traceless, so is $\Sigma_{a}$. Using \eqref{eq:tau_two_sided} with $\mu=e^{a}$, we find
\[
\|\Sigma_{a}\|_{H^{2}(M,g)}
=
e^{-29a/10}\|\tau_{e^{a}}\|_{H^{2}(M,g)}
\sim
e^{-29a/10}e^{2a}
=
e^{-9a/10},
\]
and
\[
\|\Sigma_{a}\|_{H^{3}(M,g)}
=
e^{-29a/10}\|\tau_{e^{a}}\|_{H^{3}(M,g)}
\sim
e^{-29a/10}e^{3a}
=
e^{a/10}.
\]
More precisely, there exist $a_{0}\geq 1$ and constants $0<c\leq C<\infty$ such that for all $a\geq a_{0}$,
\[
c\,e^{-9a/10}\leq \|\Sigma_{a}\|_{H^{2}(M,g)}\leq C\,e^{-9a/10},
\qquad
c\,e^{a/10}\leq \|\Sigma_{a}\|_{H^{3}(M,g)}\leq C\,e^{a/10}.
\]
This completes the proof.
\end{proof}




\begin{proposition}[Existence of a unique positive solution to the Lichnerowicz equation]
\label{prop:conformal-correction}
Fix a bounded-geometry class \(\mathcal G\) of smooth Riemannian metrics on the closed
\(3\)-manifold \(M\). Then there exist constants
\[
\delta_{0}>0,\qquad C\ge 1,
\]
depending only on \(\mathcal G\), with the following property.

\noindent Let \(g\in \mathcal G\), and let \(\Sigma\) be a smooth symmetric \(2\)-tensor satisfying
\[
\tr_{g}\Sigma=0,\qquad \div_{g}\Sigma=0,
\]
together with
\[
\left\|R[g]+\frac23\right\|_{H^{2}(g)}\le \delta_{0},
\qquad
\|\Sigma\|_{H^{2}(g)}\le \delta_{0}.
\]
Define
\[
\mathcal E(g,\Sigma):=R[g]+\frac23-|\Sigma|_{g}^{2}.
\]
Then there exists a unique positive solution
\[
\varphi\in H^{4}(M)
\]
of the Lichnerowicz equation
\begin{equation}
\label{eq:Lichnerowicz-exact-corrected}
-8\Delta_{g}\varphi + R[g]\varphi + \frac23 \varphi^{5} - |\Sigma|_{g}^{2}\varphi^{-7}=0.
\end{equation}
Moreover, \(\varphi\) is smooth, and if one defines
\[
g_{0}:=\varphi^{4}g,\qquad \Sigma_{0}:=\varphi^{-2}\Sigma,
\]
then
\begin{align}
R[g_{0}]+\frac23&=|\Sigma_{0}|_{g_{0}}^{2},\\
\div_{g_{0}}\Sigma_{0}&=0,\\
\tr_{g_{0}}\Sigma_{0}&=0.
\end{align}
In addition,
\begin{align}
\label{eq:phi-close-corrected}
\|\varphi-1\|_{H^{4}(g)}
&\leq
C\|\mathcal E(g,\Sigma)\|_{H^{2}(g)},\\
\label{eq:T-change-corrected}
\|\mathfrak T[g_{0}]-\mathfrak T[g]\|_{H^{2}(g)}
&\leq
C\|\mathcal E(g,\Sigma)\|_{H^{2}(g)}.
\end{align}
\end{proposition}

\begin{proof}
Set
\[
S:=|\Sigma|_{g}^{2},
\qquad
\rho:=R[g]+\frac23.
\]
Then define the following
\[
\mathcal E(g,\Sigma)=\rho-S
\]
which is small in $H^{2}$ norm. 
Since \(M\) is \(3\)-dimensional and \(g\in\mathcal G\), the Sobolev embeddings and
multiplication estimates are uniform in \(g\). In particular,
\[
\|S\|_{H^{2}(g)}\le C\|\Sigma\|_{H^{2}(g)}^{2},
\qquad
\|S\|_{C^{0}(g)}\le C\|\Sigma\|_{H^{2}(g)}^{2},
\qquad
\|\rho\|_{C^{0}(g)}\le C\|\rho\|_{H^{2}(g)}.
\]
Hence, for \(\delta_{0}\) sufficiently small,
\begin{equation}
\label{eq:defect-small-proof}
\|\mathcal E(g,\Sigma)\|_{H^{2}(g)}
\le C\delta_{0},
\qquad
\|\rho\|_{C^{0}(g)}+\|S\|_{C^{0}(g)}\le C\delta_{0}.
\end{equation}

\noindent We write
\[
\varphi=1+u
\]
and define
\[
\mathcal F_{g,\Sigma}(\varphi)
:=
-8\Delta_{g}\varphi + R[g]\varphi + \frac23\varphi^{5}-S\varphi^{-7}.
\]
Then \eqref{eq:Lichnerowicz-exact-corrected} is equivalent to
\[
\mathcal F_{g,\Sigma}(1+u)=0.
\]
Moreover,
\[
\mathcal F_{g,\Sigma}(1)=\rho-S=\mathcal E(g,\Sigma).
\]

\noindent First, we understand the linearization of $\mathcal F_{g,\Sigma}$. The linearization at \(u=0\) is
\[
L_{g,\Sigma}u
:=
D\mathcal F_{g,\Sigma}|_{\varphi=1}[u]
=
-8\Delta_{g}u+\Bigl(R[g]+\frac{10}{3}+7S\Bigr)u.
\]
Since
\[
R[g]+\frac{10}{3}+7S
=
\frac83+\rho+7S,
\]
the pointwise bound \eqref{eq:defect-small-proof} implies, after shrinking \(\delta_{0}\),
that
\[
R[g]+\frac{10}{3}+7S\ge 1
\qquad\text{on }M.
\]
Therefore
\[
\int_{M} \langle L_{g,\Sigma}u,u\rangle\,d\mu_{g}
=
8\|\nabla u\|_{L^{2}(g)}^{2}
+
\int_{M}\Bigl(R[g]+\frac{10}{3}+7S\Bigr)u^{2}\,d\mu_{g}
\ge c\|u\|_{H^{1}(g)}^{2}.
\]
Hence \(L_{g,\Sigma}\) has trivial kernel. Since it is a strongly elliptic second-order
operator on a closed manifold, it is Fredholm of index zero, and thus
\[
L_{g,\Sigma}\colon H^{4}(M)\to H^{2}(M)
\]
is an isomorphism. By uniform elliptic regularity on the bounded-geometry class
\(\mathcal G\), there is a constant \(K\) such that
\begin{equation}
\label{eq:elliptic-L-corrected}
\|u\|_{H^{4}(g)}
\le
K\|L_{g,\Sigma}u\|_{H^{2}(g)}.
\end{equation}

\noindent Next we isolate the nonlinear remainder. For \(u\) sufficiently small in \(C^{0}\),
\[
(1+u)^{5}=1+5u+Q_{5}(u),
\qquad
(1+u)^{-7}=1-7u+Q_{-7}(u),
\]
where \(Q_{5}(u)\) and \(Q_{-7}(u)\) vanish quadratically at \(u=0\). Thus
\[
\mathcal F_{g,\Sigma}(1+u)
=
\mathcal E(g,\Sigma)+L_{g,\Sigma}u+\mathcal N_{g,\Sigma}(u),
\]
with
\[
\mathcal N_{g,\Sigma}(u)
=
\frac23 Q_{5}(u)-S\,Q_{-7}(u).
\]

\noindent Let \(C_{\mathrm{emb}}\) denote the uniform norm of the embedding
\(H^{4}(M)\hookrightarrow C^{0}(M)\) on \(\mathcal G\), and set
\[
\rho_{*}:=\frac{1}{4C_{\mathrm{emb}}}.
\]
If \(\|u\|_{H^{4}(g)}\le \rho_{*}\), then \(\|u\|_{C^{0}(g)}\le \frac14\), so
\[
\frac34\le 1+u\le \frac54.
\]
On this set, uniform Moser estimates imply
\begin{align}
\label{eq:Nquad-corrected}
\|\mathcal N_{g,\Sigma}(u)\|_{H^{2}(g)}
&\le
C\|u\|_{H^{4}(g)}^{2},\\
\label{eq:Nlip-corrected}
\|\mathcal N_{g,\Sigma}(u)-\mathcal N_{g,\Sigma}(v)\|_{H^{2}(g)}
&\le
C\bigl(\|u\|_{H^{4}(g)}+\|v\|_{H^{4}(g)}\bigr)\|u-v\|_{H^{4}(g)}.
\end{align}

\noindent Define
\[
\Phi(u):=
-L_{g,\Sigma}^{-1}\bigl(\mathcal E(g,\Sigma)+\mathcal N_{g,\Sigma}(u)\bigr).
\]
Let
\[
r:=2K\|\mathcal E(g,\Sigma)\|_{H^{2}(g)}.
\]
If \(\delta_{0}\) is sufficiently small, then by \eqref{eq:defect-small-proof},
\[
r\le \rho_{*},
\qquad
CKr\le \frac12,
\qquad
CKr^{2}\le \frac r2.
\]
Now let \(\|u\|_{H^{4}(g)}\le r\). Using \eqref{eq:elliptic-L-corrected} and
\eqref{eq:Nquad-corrected},
\[
\|\Phi(u)\|_{H^{4}(g)}
\le
K\|\mathcal E(g,\Sigma)\|_{H^{2}(g)}
+
K\|\mathcal N_{g,\Sigma}(u)\|_{H^{2}(g)}
\le
\frac r2+CKr^{2}
\le r.
\]
Hence \(\Phi\) maps the closed \(H^{4}\)-ball \(B_{r}\) into itself. Likewise, for
\(u,v\in B_{r}\), \eqref{eq:Nlip-corrected} gives
\[
\|\Phi(u)-\Phi(v)\|_{H^{4}(g)}
\le
K\|\mathcal N_{g,\Sigma}(u)-\mathcal N_{g,\Sigma}(v)\|_{H^{2}(g)}
\le
2CKr\,\|u-v\|_{H^{4}(g)}
\le
\frac12\|u-v\|_{H^{4}(g)}.
\]
Thus \(\Phi\) is a contraction on \(B_{r}\). By Banach's fixed-point theorem, there exists
a unique \(u\in B_{r}\) such that \(\Phi(u)=u\). Setting \(\varphi:=1+u\), we obtain a
solution of \eqref{eq:Lichnerowicz-exact-corrected} with
\[
\|\varphi-1\|_{H^{4}(g)}=\|u\|_{H^{4}(g)}
\le
C\|\mathcal E(g,\Sigma)\|_{H^{2}(g)},
\]
which proves \eqref{eq:phi-close-corrected}. Since \(\|u\|_{C^{0}(g)}\le \frac14\), we also
have \(\varphi\ge \frac34>0\).

\noindent We now prove uniqueness among \emph{all} positive solutions. Let \(\varphi\) be any
positive solution of \eqref{eq:Lichnerowicz-exact-corrected}. Let
\[
m:=\min_{M}\varphi,
\qquad
M:=\max_{M}\varphi.
\]
At a point of maximum,
\[
0
=
-8\Delta_{g}\varphi + R[g]\varphi + \frac23\varphi^{5}-S\varphi^{-7}
\ge
R[g]M+\frac23 M^{5}-S M^{-7},
\]
since $\Delta_{g} \varphi<0$ at that point.
Hence
\[
\frac23 M(M^{4}-1)\le |\rho|\,M + S M^{-7}.
\]
Similarly, at a point of minimum,
\[
0
=
-8\Delta_{g}\varphi + R[g]\varphi + \frac23\varphi^{5}-S\varphi^{-7}
\le
R[g]m+\frac23 m^{5}-S m^{-7},
\]
since $\Delta_{g} \varphi>0$ at that point.
Hence
\[
\frac23 m(1-m^{4})\le |\rho|\,m + S m^{-7}.
\]
Using \eqref{eq:defect-small-proof}, one now chooses \(\delta_{0}\) so small that these two
inequalities force
\[
\frac34\le m\le M\le \frac54.
\]
Thus every positive solution takes values in \([\frac34,\frac54]\).

\noindent  Now let \(\varphi_{1},\varphi_{2}\) be two positive solutions, and set
\[
w:=\varphi_{1}-\varphi_{2}.
\]
Subtracting the two equations and using the mean value theorem pointwise yields
\[
-8\Delta_{g}w + a(x)w=0,
\]
where
\[
a(x)
=
R[g]
+
\frac{10}{3}\theta_{1}(x)^{4}
+
7S(x)\theta_{2}(x)^{-8},
\]
for some \(\theta_{1}(x),\theta_{2}(x)\in[\frac34,\frac54]\). Hence
\[
a(x)\ge -\|\rho\|_{C^{0}(g)}-\frac23+\frac{10}{3}\Bigl(\frac34\Bigr)^{4}.
\]
After possibly shrinking \(\delta_{0}\) once more, the right-hand side is bounded below by
a positive constant \(c_{*}>0\). Therefore
\[
-8\Delta_{g}w+c_{*}w\le 0
\qquad\text{and}\qquad
-8\Delta_{g}(-w)+c_{*}(-w)\le 0,
\]
since the order of $\varphi_{1}$ and $\varphi_{2}$ here does not matter.
So the maximum principle implies \(w\equiv0\). Thus the positive solution is unique.

\noindent Define
\[
g_{0}:=\varphi^{4}g,\qquad \Sigma_{0}:=\varphi^{-2}\Sigma.
\]
The standard conformal covariance identities in dimension \(3\) give
\[
R[g_{0}]
=
\varphi^{-5}\bigl(-8\Delta_{g}\varphi+R[g]\varphi\bigr),
\qquad
|\Sigma_{0}|_{g_{0}}^{2}
=
\varphi^{-12}|\Sigma|_{g}^{2}.
\]
Hence \eqref{eq:Lichnerowicz-exact-corrected} is equivalent to
\[
R[g_{0}]+\frac23=|\Sigma_{0}|_{g_{0}}^{2}.
\]
Since \(\Sigma\) is \(g\)-trace-free and \(g\)-divergence-free, the standard conformal
transformation law for TT tensors implies
\[
\tr_{g_{0}}\Sigma_{0}=0,
\qquad
\div_{g_{0}}\Sigma_{0}=0.
\]

\noindent It remains to prove \eqref{eq:T-change-corrected}. In dimension \(3\),
\[
\mathfrak T[\varphi^{4}g]
=
\mathfrak T[g]
-
2\varphi^{-1}\Bigl(\nabla^{2}\varphi-\frac13(\Delta_{g}\varphi)g\Bigr)
+
6\varphi^{-2}\Bigl(d\varphi\otimes d\varphi-\frac13|\nabla\varphi|_{g}^{2}g\Bigr).
\]
Therefore
\begin{align*}
\mathfrak T[g_{0}]-\mathfrak T[g]
&=
-2\varphi^{-1}\Bigl(\nabla^{2}\varphi-\frac13(\Delta_{g}\varphi)g\Bigr)\\
&\qquad
+
6\varphi^{-2}\Bigl(d\varphi\otimes d\varphi-\frac13|\nabla\varphi|_{g}^{2}g\Bigr).
\end{align*}
Since \(\varphi\in[\frac34,\frac54]\) pointwise and
\(\|\varphi-1\|_{H^{4}(g)}\le C\|\mathcal E(g,\Sigma)\|_{H^{2}(g)}\), uniform composition
estimates on the bounded-geometry class imply
\[
\|\varphi^{-1}\|_{H^{4}(g)}+\|\varphi^{-2}\|_{H^{4}(g)}\le C.
\]
Using that \(H^{2}(M)\) is an algebra in dimension \(3\),
\[
\left\|
\varphi^{-1}\Bigl(\nabla^{2}\varphi-\frac13(\Delta_{g}\varphi)g\Bigr)
\right\|_{H^{2}(g)}
\le
C\|\varphi-1\|_{H^{4}(g)},
\]
and
\[
\left\|
\varphi^{-2}\Bigl(d\varphi\otimes d\varphi-\frac13|\nabla\varphi|_{g}^{2}g\Bigr)
\right\|_{H^{2}(g)}
\le
C\|\varphi-1\|_{H^{4}(g)}^{2}.
\]
Hence
\[
\|\mathfrak T[g_{0}]-\mathfrak T[g]\|_{H^{2}(g)}
\le
C\|\varphi-1\|_{H^{4}(g)}
+
C\|\varphi-1\|_{H^{4}(g)}^{2}.
\]
For \(\delta_{0}\) sufficiently small, the quadratic term is absorbed into the linear one,
and \eqref{eq:phi-close-corrected} yields
\[
\|\mathfrak T[g_{0}]-\mathfrak T[g]\|_{H^{2}(g)}
\le
C\|\mathcal E(g,\Sigma)\|_{H^{2}(g)}.
\]
This proves \eqref{eq:T-change-corrected}.

\noindent Finally, since \(g\) and \(\Sigma\) are smooth and \(\varphi\) is bounded above and below
away from zero, standard elliptic bootstrapping applied to
\eqref{eq:Lichnerowicz-exact-corrected} shows that \(\varphi\in C^{\infty}(M)\).
\end{proof}
\begin{corollary}[Persistence of a large \(H^{2}\)-norm of the trace-free Ricci tensor]
\label{cor:persistence}
Under the hypotheses of Proposition~\ref{prop:conformal-correction}, let \(A>0\). Assume in addition that
\[
\|\mathfrak T[g]\|_{H^{2}(g)}\ge A,
\qquad
\|\mathcal E(g,\Sigma)\|_{H^{2}(g)}\le \frac{A}{2C},
\]
where \(C\) is the constant in \eqref{eq:T-change-corrected}. Then
\[
\|\mathfrak T[g_{0}]\|_{H^{2}(g)}\ge \frac{A}{2}.
\]
\end{corollary}

\begin{proof}
By \eqref{eq:T-change-corrected},
\[
\|\mathfrak T[g_{0}]-\mathfrak T[g]\|_{H^{2}(g)}
\le
C\|\mathcal E(g,\Sigma)\|_{H^{2}(g)}
\le
\frac{A}{2}.
\]
Therefore, by the reverse triangle inequality,
\[
\|\mathfrak T[g_{0}]\|_{H^{2}(g)}
\ge
\|\mathfrak T[g]\|_{H^{2}(g)}
-
\|\mathfrak T[g_{0}]-\mathfrak T[g]\|_{H^{2}(g)}
\ge
A-\frac{A}{2}
=
\frac{A}{2}.
\]
This proves the claim.
\end{proof}
\noindent The following corollary provides the exact initial data that we can address in this study.
\begin{corollary}[Exact CMC initial data]
\label{cor:final-CMC}
There exist constants \(a_{0}\ge 1\) and \(0<c\le C\), depending only on \((M,\gamma)\), with the following property. For every \(a\ge a_{0}\), there exists a smooth initial data set \((g_{0,a},\Sigma_{0,a})\) on \(M\) such that
\begin{align}
R[g_{0,a}]+\frac23&=|\Sigma_{0,a}|_{g_{0,a}}^{2},\\
\div_{g_{0,a}}\Sigma_{0,a}&=0,\\
\tr_{g_{0,a}}\Sigma_{0,a}&=0.
\end{align}
Moreover,
\begin{equation}
\label{eq:g0a-gamma-equivalence}
C^{-1}\gamma\le g_{0,a}\le C\gamma,
\end{equation}
and, with all Sobolev norms and covariant derivatives taken with respect to \(\gamma\),
\begin{equation}
\label{eq:final-CMC-upper}
\sum_{I=0}^{3}\|\nabla_\gamma^{I}\Sigma_{0,a}\|_{L^{2}(\gamma)}
+\sum_{I=0}^{2}\Bigl(
\|\nabla_\gamma^{I}\mathfrak T[g_{0,a}]\|_{L^{2}(\gamma)}
+
\|e^{a}\nabla_\gamma^{I}\Sigma_{0,a}\|_{L^{2}(\gamma)}
\Bigr)
\le C e^{a/10},
\end{equation}
and moreover
\begin{equation}
\label{eq:final-CMC-lower}
\sum_{I=0}^{3}\|\nabla_\gamma^{I}\Sigma_{0,a}\|_{L^{2}(\gamma)}\lesssim e^{a/10},\sum_{I=0}^{2}\Bigl(
\|\nabla_\gamma^{I}\mathfrak T[g_{0,a}]\|_{L^{2}(\gamma)}
\Bigr)
\lesssim e^{a/10}, 
\|e^{a}\nabla_\gamma^{I}\Sigma_{0,a}\|_{L^{2}(\gamma)}\lesssim e^{a/10},
\end{equation}
for constants $\leq C$.
\end{corollary}



\begin{proof}
We combine the three ingredients established earlier in the construction of the initial data: the seed metric with controlled scalar-curvature defect and large \(H^{2}\)-size of the renormalized trace-free Ricci tensor, the construction of a transverse-traceless tensor of the required size, and the perturbative conformal correction which solves the Hamiltonian constraint exactly.

\medskip

\noindent
By the seed-metric construction proved earlier, there exist constants \(a_{1}\ge 1\) and \(C\ge 1\), depending only on \((M,\gamma)\), such that for every \(a\ge a_{1}\) there exists a smooth metric \(\widehat g_{a}\) on \(M\) satisfying
\begin{equation}
\label{eq:seed-metric-cor}
C^{-1}\gamma\le \widehat g_{a}\le C\gamma,
\end{equation}
together with
\begin{equation}
\label{eq:seed-scalar-cor}
\left\|R[\widehat g_{a}]+\frac23\right\|_{H^{2}(\gamma)}
\le C e^{-9a/5},
\qquad
\left\|R[\widehat g_{a}]+\frac23\right\|_{H^{3}(\gamma)}
\le C e^{a/5},
\end{equation}
and
\begin{equation}
\label{eq:seed-T-cor}
\|\mathfrak T[\widehat g_{a}]\|_{H^{2}(\gamma)}
\le C e^{a/10}.
\end{equation}

\medskip

\noindent By the transverse-traceless construction proved earlier, after possibly increasing \(a_{1}\), there exists a smooth symmetric \(2\)-tensor \(\widehat\Sigma_{a}\) such that
\begin{equation}
\label{eq:seed-TT-cor}
\div_{\widehat g_{a}}\widehat\Sigma_{a}=0,
\qquad
\tr_{\widehat g_{a}}\widehat\Sigma_{a}=0,
\end{equation}
and
\begin{equation}
\label{eq:seed-sigma-cor}
\|\widehat\Sigma_{a}\|_{H^{2}(\gamma)}\le C e^{-9a/10},
\qquad
\|\widehat\Sigma_{a}\|_{H^{3}(\gamma)}\le C e^{a/10}.
\end{equation}

\noindent Since \(M\) is a closed \(3\)-manifold, \(H^{2}(M)\) and \(H^{3}(M)\) are Banach algebras. Using also the uniform equivalence \eqref{eq:seed-metric-cor}, we obtain
\begin{equation}
\label{eq:sigma-square-H2-cor}
\bigl\||\widehat\Sigma_{a}|_{\widehat g_{a}}^{2}\bigr\|_{H^{2}(\gamma)}
\le C \|\widehat\Sigma_{a}\|_{H^{2}(\gamma)}^{2}
\le C e^{-9a/5},
\end{equation}
and similarly
\begin{equation}
\label{eq:sigma-square-H3-cor}
\bigl\||\widehat\Sigma_{a}|_{\widehat g_{a}}^{2}\bigr\|_{H^{3}(\gamma)}
\le C \|\widehat\Sigma_{a}\|_{H^{3}(\gamma)}^{2}
\le C e^{a/5}.
\end{equation}

\noindent Therefore the Hamiltonian defect of the seed pair \((\widehat g_{a},\widehat\Sigma_{a})\) obeys
\begin{equation}
\label{eq:seed-defect-cor}
\left\|R[\widehat g_{a}]+\frac23-|\widehat\Sigma_{a}|_{\widehat g_{a}}^{2}\right\|_{H^{2}(\gamma)}
\le C e^{-9a/5},
\end{equation}
and likewise
\begin{equation}
\label{eq:seed-defect-H3-cor}
\left\|R[\widehat g_{a}]+\frac23-|\widehat\Sigma_{a}|_{\widehat g_{a}}^{2}\right\|_{H^{3}(\gamma)}
\le C e^{a/5}.
\end{equation}

\medskip

\noindent
We now apply the perturbative conformal correction theorem to the seed pair \((\widehat g_{a},\widehat\Sigma_{a})\). It yields a smooth positive function \(\varphi_{a}\) such that, upon setting
\begin{equation}
\label{eq:conformal-cor}
g_{0,a}:=\varphi_{a}^{4}\widehat g_{a},
\qquad
\Sigma_{0,a}:=\varphi_{a}^{-2}\widehat\Sigma_{a},
\end{equation}
the pair \((g_{0,a},\Sigma_{0,a})\) satisfies
\begin{align}
R[g_{0,a}]+\frac23&=|\Sigma_{0,a}|_{g_{0,a}}^{2},\\
\div_{g_{0,a}}\Sigma_{0,a}&=0,\\
\tr_{g_{0,a}}\Sigma_{0,a}&=0.
\end{align}
Moreover, the conformal theorem yields the perturbative estimate
\begin{equation}
\label{eq:phi-small-cor}
\|\varphi_{a}-1\|_{H^{4}(\gamma)}\le C e^{-9a/5}.
\end{equation}

\noindent Since \(H^{4}(M)\hookrightarrow C^{2}(M)\) on a closed \(3\)-manifold, \eqref{eq:phi-small-cor} implies
\[
\|\varphi_{a}-1\|_{L^{\infty}(M)}\le C e^{-9a/5}.
\]
\noindent Hence, after increasing \(a_{0}\ge a_{1}\) if necessary, we may assume
\[
\frac12\le \varphi_{a}\le 2
\qquad\text{on }M.
\]
Combining this with \eqref{eq:seed-metric-cor}, we obtain
\[
C^{-1}\gamma\le g_{0,a}\le C\gamma,
\]
which proves \eqref{eq:g0a-gamma-equivalence}.

\medskip

\noindent
From \eqref{eq:conformal-cor},
\[
\Sigma_{0,a}-\widehat\Sigma_{a}=(\varphi_{a}^{-2}-1)\widehat\Sigma_{a}.
\]
By standard composition estimates in Sobolev spaces and \eqref{eq:phi-small-cor},
\begin{equation}
\label{eq:phi-inv-small-cor}
\|\varphi_{a}^{-2}-1\|_{H^{4}(\gamma)}
\le C\|\varphi_{a}-1\|_{H^{4}(\gamma)}
\le C e^{-9a/5}.
\end{equation}
Using the algebra property of \(H^{2}\) and \(H^{3}\), together with \eqref{eq:seed-sigma-cor}, we infer
\begin{align}
\|\Sigma_{0,a}\|_{H^{2}(\gamma)}
&\le \|\widehat\Sigma_{a}\|_{H^{2}(\gamma)}
   +\|(\varphi_{a}^{-2}-1)\widehat\Sigma_{a}\|_{H^{2}(\gamma)} \nonumber\\
&\le C e^{-9a/10}+C e^{-9a/5}e^{-9a/10}
\le C e^{-9a/10},
\label{eq:sigma-H2-cor}
\end{align}
and likewise
\begin{align}
\|\Sigma_{0,a}\|_{H^{3}(\gamma)}
&\le \|\widehat\Sigma_{a}\|_{H^{3}(\gamma)}
   +\|(\varphi_{a}^{-2}-1)\widehat\Sigma_{a}\|_{H^{3}(\gamma)} \nonumber\\
&\le C e^{a/10}+C e^{-9a/5}e^{a/10}
\le C e^{a/10}.
\label{eq:sigma-H3-cor}
\end{align}
This proves the first and last bounds in \eqref{eq:final-CMC-lower}.

\medskip

\noindent
From \eqref{eq:conformal-cor},
\[
g_{0,a}-\widehat g_{a}=(\varphi_{a}^{4}-1)\widehat g_{a}.
\]
Again by the composition estimates and \eqref{eq:phi-small-cor},
\[
\|\varphi_{a}^{4}-1\|_{H^{4}(\gamma)}\le C e^{-9a/5},
\]
and therefore
\begin{equation}
\label{eq:g-difference-cor}
\|g_{0,a}-\widehat g_{a}\|_{H^{4}(\gamma)}\le C e^{-9a/5}.
\end{equation}

\noindent Now the map
\[
g\mapsto \mathfrak T[g]
\]
is a smooth nonlinear differential operator of order \(2\). On any bounded subset of the \(H^{4}\)-neighborhood of \(\gamma\), it is locally Lipschitz from \(H^{4}\) to \(H^{2}\). Hence \eqref{eq:g-difference-cor} implies
\begin{equation}
\label{eq:T-difference-cor}
\|\mathfrak T[g_{0,a}]-\mathfrak T[\widehat g_{a}]\|_{H^{2}(\gamma)}
\le C \|g_{0,a}-\widehat g_{a}\|_{H^{4}(\gamma)}
\le C e^{-9a/5}.
\end{equation}
Combining \eqref{eq:T-difference-cor} with \eqref{eq:seed-T-cor}, we conclude
\begin{equation}
\label{eq:T-final-cor}
\|\mathfrak T[g_{0,a}]\|_{H^{2}(\gamma)}
\le C e^{a/10}.
\end{equation}
This is exactly the second bound asserted in \eqref{eq:final-CMC-lower}.

\medskip

\noindent
By definition of Sobolev norms with respect to \(\gamma\),
\[
\sum_{I=0}^{3}\|\nabla_{\gamma}^{I}\Sigma_{0,a}\|_{L^{2}(\gamma)}
\le C \|\Sigma_{0,a}\|_{H^{3}(\gamma)}
\le C e^{a/10}
\]
by \eqref{eq:sigma-H3-cor}. Also,
\[
\sum_{I=0}^{2}\|e^{a}\nabla_{\gamma}^{I}\Sigma_{0,a}\|_{L^{2}(\gamma)}
\le C e^{a}\|\Sigma_{0,a}\|_{H^{2}(\gamma)}
\le C e^{a}e^{-9a/10}
= C e^{a/10}
\]
by \eqref{eq:sigma-H2-cor}. Finally,
\[
\sum_{I=0}^{2}\|\nabla_{\gamma}^{I}\mathfrak T[g_{0,a}]\|_{L^{2}(\gamma)}
\le C \|\mathfrak T[g_{0,a}]\|_{H^{2}(\gamma)}
\le C e^{a/10}
\]
by \eqref{eq:T-final-cor}. This proves \eqref{eq:final-CMC-upper}.

\noindent Therefore, for every \(a\ge a_{0}\), there exists a smooth exact CMC initial data set \((g_{0,a},\Sigma_{0,a})\) satisfying all the stated properties. This completes the proof.
\end{proof}


\begin{remark}
 Notice the vital point: by choosing the perturbations to be transverse-traceless (or $TT$), we eliminate the principal leading order term in the scalar curvature deviation. This is the precise mechanism that yields almost unmodified scalar curvature in $H^{2}$ norm and hence by Sobolev embedding, in the point-wise sense as well. On the other hand, the principal leading order term in the deviation of the trace-free Ricci curvature $\mathfrak{T}$ is the main contributor to its large   $H^{2}$ norm.
\end{remark}



\section{Hyperbolic Estimates}

\subsection{Setting up the bootstrap argument}
\label{bootstrapargument}
\noindent In order to establish global-in-time control of the solutions of the Einstein vacuum equations with cosmological constant $\Lambda>0$—in the rescaled variables introduced earlier—we employ a bootstrap argument. Specifically, we aim to derive uniform bounds on three scale-invariant norms:
\[
\mathcal{O}, \mathcal{F}, \mathcal{N}^\infty,
\]
which respectively control curvature energies, renormalized entropy for the transverse-traceless ($TT$) second fundamental form, and point-wise norm of gauge variable and $TT$ tensor $\Sigma$. These norms are defined in the rescaled setting associated with the dynamical time variable $T \in [T_0, T_\infty)$.

\noindent We assume the initial data at time $T = T_0$ satisfies the bound
\begin{align} \label{eq:initial-bound}
\mathcal{O}(T_0) + \mathcal{F}(T_0) + \mathcal{N}^\infty(T_0) \lesssim \mathcal{I}^0,
\end{align}
where $\mathcal{I}^0 \gg 1$ denotes the (large) initial data size. Our main objective is to prove that for all $T \in [T_0, T_\infty]$, the combined norm satisfies the uniform a priori estimate
\begin{align} \label{eq:goal-bound}
\mathcal{O}(T) + \mathcal{F}(T) + \mathcal{N}^\infty(T) \lesssim \mathcal{I}^0 + 1.
\end{align}
This suffices to propagate the solution globally in time via standard local well-posedness and continuation arguments (see Section~\ref{local}).

\noindent To achieve \eqref{eq:goal-bound}, we proceed via a bootstrap scheme. We assume that on a given time interval $[T_0, T] \subset [T_0, T_\infty)$, the following bounds hold:
\begin{align} \label{eq:strap2}
\mathcal{O}(T) \leq \Gamma, \qquad \mathcal{F}(T) \leq \mathds{Y}, \qquad \mathcal{N}^\infty(T) \leq \mathds{L},
\end{align}
where the constants $\Gamma, \mathds{Y}, \mathds{L}$ satisfy
\begin{align} \label{eq:bootstrap-size}
(\mathcal{I}^0)^2 + \mathcal{I}^0 + 1 < \min\{\Gamma, \mathds{Y}, \mathds{L}\},
\end{align}
and are additionally constrained by
\begin{align} \label{eq:bootstrap-growth}
\Gamma + \mathds{Y} + \mathds{L} \leq e^{\frac{T_0}{10}}.
\end{align}

\noindent Define the set of admissible bootstrap times:
\[
\mathcal{S} := \left\{ T \in [T_0, T_\infty] \mid \text{the bounds \eqref{eq:strap2} hold on } [T_0, T] \right\}.
\]
Our aim is to prove that $\mathcal{S} = [T_0, T_\infty]$ by showing that $\mathcal{S}$ is both open and closed.

\paragraph{Step 1: Non-emptiness} By local well-posedness theory for the Einstein--$\Lambda$ system in the chosen gauge, there exists $\varepsilon > 0$ such that a solution exists on $[T_0, T_0 + \varepsilon]$ and satisfies
\begin{align}
\mathcal{O}(T) + \mathcal{F}(T) + \mathcal{N}^\infty(T) \lesssim 2 \mathcal{I}^0, \qquad \forall T \in [T_0, T_0 + \varepsilon],
\end{align}
thus implying $\mathcal{S} \neq \{\emptyset\}$.

\paragraph{Step 2: Closedness} The closedness of $\mathcal{S}$ follows from standard continuation results and the fact that the norms $\mathcal{O}, \mathcal{F}, \mathcal{N}^\infty$ are continuous in time. More precisely, the uniform bounds implied by the bootstrap assumptions extend continuously to the supremum of $\mathcal{S}$, which must then also belong to $\mathcal{S}$.

\paragraph{Step 3: Openness} The core of the argument lies in establishing a hierarchy of estimates:
\[
\mathcal{N}^\infty \longrightarrow \mathcal{F} \longrightarrow \mathcal{O},
\]
that enable one to strictly improve the assumptions in \eqref{eq:strap2}. These estimates are derived as follows:

\begin{itemize}
  \item The gauge quantities in $\mathcal{N}^\infty$ are controlled via elliptic regularity theory and Sobolev inequalities on the spatial slices $\Sigma_T$, leading to the estimate
  \begin{align} \label{eq:estN}
  \mathcal{N}^\infty(T) \lesssim \mathcal{O}(T) + \mathcal{F}(T)^2 + \mathcal{F}(T).
  \end{align}
  
  \item The norm $\mathcal{F}$, which captures the trace-free part of the second fundamental form, is controlled by integrating the transport equation for $\Sigma$ forward in $T$, yielding
  \begin{align} \label{eq:estF}
  \mathcal{F}(T) \lesssim \mathcal{O}(T) + \mathcal{I}^0 + 1.
  \end{align}
  
  \item Finally, $\mathcal{O}$ is estimated via energy methods, applied to the hyperbolic system, giving the uniform bound
  \begin{align} \label{eq:estO}
  \mathcal{O}(T) \lesssim \mathcal{I}^0 + 1.
  \end{align}
\end{itemize}

\noindent Combining \eqref{eq:estN}--\eqref{eq:estO}, and choosing the bootstrap constants sufficiently large relative to $\mathcal{I}^0$ yet satisfying \eqref{eq:bootstrap-growth}, we obtain strict improvements of the assumptions in \eqref{eq:strap2}. Therefore, by the local well-posedness theory, the solution to the Einsten-$\Lambda$ equation can be extended a bit towards a larger value of $T$, implying the openness of $\mathcal{S}$.

\noindent Since $\mathcal{S} \subset [T_0, T_\infty]$ is non-empty, open, and closed, and $[T_0, T_\infty]$ is connected, it follows that
\[
\mathcal{S} = [T_0, T_\infty].
\]
Thus, the solution exists globally in time, and the norms $\mathcal{O}, \mathcal{F}, \mathcal{N}^\infty$ obey uniform bounds independent of $T_\infty$. Sending $T_\infty \to \infty$ yields global existence and uniform control for all future times.
   

\subsection{Estimates for the metric components}
\label{metriccontrol}
\begin{proposition}[Uniform Control of the Metric Volume Density]
\label{metric1}
Let $(\widetilde{M}, \widehat{g}(T))$ be a smooth globally hyperbolic spacetime foliated by constant mean curvature (CMC) slices, and suppose that the metric \( g(T) \) evolves according to the Einstein vacuum equations with cosmological constant in the CMC-spatially transported gauge. Let the initial data at \( T = T_0 \) satisfy the assumptions of Theorem~\ref{main}, and assume the bootstrap condition
\[
\left\| \frac{N}{n} - 1 \right\|_{L^\infty(M)} \leq \mathds{L} e^{-2\gamma T_{0}}, \qquad \text{for all } T \in [T_0, T],
\]
where \( \mathds{L}, \gamma > 0 \) are constants. Then the Riemannian volume density \( \mu_g := \sqrt{\det g_{ij}} \) associated with the induced metric \( g_{ij}(T) \) satisfies the uniform bounds
\[
C_1 \leq \det(g_{ij}(T)) \leq C_2, \qquad \text{for all } T \in [T_0, T],
\]
for some positive constants \( C_1, C_2 > 0 \) depending only on the initial data at \( T_0 \).
\end{proposition}

\begin{proof}
We begin with the evolution equation for the spatial metric \( g_{ij} \) in CMC-spatially transported gauge:
\[
\partial_T g_{ij} = -\frac{2\varphi}{\tau} N \Sigma_{ij} - 2\left(1 - \frac{N}{n}\right) g_{ij}.
\]
Taking the trace with respect to \( g^{ij} \), and using the fact that \( \operatorname{tr}_g \Sigma = 0 \), we obtain:
\[
\operatorname{tr}_g(\partial_T g) = -2n\left(1 - \frac{N}{n}\right).
\]
Let \( \mu_g := \sqrt{\det g_{ij}} \) denote the Riemannian volume density. The evolution of \( \mu_g \) is governed by:
\[
\partial_T \mu_g = \frac{1}{2} \mu_g \operatorname{tr}_g(\partial_T g) = -n \mu_g \left(1 - \frac{N}{n}\right),
\]
which implies
\[
\mu_g(T) = \mu_g(T_0) \exp\left(n \int_{T_0}^{T} \left(\frac{N}{n} - 1 \right)(s) \, ds \right).
\]
Applying the bootstrap assumption gives:
\[
\left| \log\left(\frac{\mu_g(T)}{\mu_g(T_0)}\right) \right| \leq n \int_{T_0}^{T} \left| \frac{N}{n} - 1 \right|(s) ds \leq \frac{n \mathds{L}}{2\gamma} e^{-2\gamma T_0}.
\]
Exponentiating, we obtain the bounds:
\[
\mu_g(T_0) e^{-\frac{n\mathds{L}}{2\gamma} e^{-2\gamma T_0}} \leq \mu_g(T) \leq \mu_g(T_0) e^{\frac{n\mathds{L}}{2\gamma} e^{-2\gamma T_0}},
\]
or 
\[
|\mu_{g}(T)-\mu_{g}(T_{0})|\leq \mu_{g}(T_{0})|e^{\frac{n\mathds{L}}{2\gamma}e^{-2\gamma T_{0}}}-1|\lesssim \mathds{L}e^{-2\gamma T_{0}}. 
\]
which imply, in particular, the uniform bounds on \( \det g_{ij} \). The constants \( C_1, C_2 \) depend only on the initial data.
\end{proof}

\begin{corollary}[Exponential Stability of Volume Function]
\label{cor:volume-stability}
Under the same assumptions as in Proposition~\ref{metric1}, the total volume functional satisfies the bound:
\[
\left| \operatorname{Vol}(g(T)) - \operatorname{Vol}(g(T_0)) \right| \lesssim \mathds{L} e^{-2\gamma T_0}, \qquad \text{for all } T \in [T_0, T_\infty],
\]
where \( \operatorname{Vol}(g(T)) = \int_M \mu_g(T) \, dx \) denotes the total volume with respect to \( g(T) \).
\end{corollary}

\begin{proof}
From the previous proposition and the evolution equation for \( \mu_g \), we compute:
\[
\frac{d}{dT} \operatorname{Vol}(g(T)) = \int_M \partial_T \mu_g \, dx = n \int_M \mu_g\left(\frac{N}{n} - 1\right) dx.
\]
Integrating in time and using the uniform bound \( \left\| \frac{N}{n} - 1 \right\|_{L^\infty(M)} \leq \mathds{L} e^{-2\gamma T} \), we deduce:
\[
\left| \operatorname{Vol}(g(T)) - \operatorname{Vol}(g(T_0)) \right| \lesssim \operatorname{Vol}(g(T_0))  \int_{T_0}^{T} \mathds{L} e^{-2\gamma t} dt \lesssim \mathds{L} e^{-2\gamma T_0}.
\]
This completes the proof.
\end{proof}
\begin{proposition}
\label{metric2}
Let $(x,T)\in M\times [T_{0},T_{\infty}]$ denote coordinates on a globally hyperbolic spacetime foliated by constant mean curvature (CMC) hypersurfaces, and let $g(T,x)$ denote the induced Riemannian metric on the CMC slices in CMC-transported coordinates. Denote by $\alpha(T)$ and $\beta(T)$ the maximal and minimal eigenvalues, respectively, of the symmetric bilinear form $g(T,x)$ with respect to the initial metric $g(T_0,x)$. Then under the bootstrap assumption \eqref{eq:strap2}, we have the estimate
\[
|\alpha(T) - 1| + |\beta(T) - 1| \lesssim \mathds{L} e^{-2\gamma T_0},
\]
where $\mathds{L}$ is a constant depending on the norms in the bootstrap assumptions and $\gamma > 0$ is a fixed decay rate.
\end{proposition}

\begin{proof}
Fix a point $x \in M$, and consider the symmetric bilinear form $g(T,x)$ on $T_x M$. Define the maximal and minimal eigenvalues of $g(T,x)$ with respect to $g(T_0,x)$ at the point $x$ as
\begin{align}
\label{eq:eigdef}
\alpha(x,T) := \sup_{0 \neq Y \in T_x M} \frac{g(T,x)(Y,Y)}{g(T_0,x)(Y,Y)}, \quad
\beta(x,T) := \inf_{0 \neq Y \in T_x M} \frac{g(T,x)(Y,Y)}{g(T_0,x)(Y,Y)}.
\end{align}
We aim to estimate $|\alpha(x,T)-1|$ and $|\beta(x,T)-1|$ uniformly in $x$ and $T \in [T_0,T_\infty]$.
Recall the evolution equation for the metric in CMC-transported coordinates:
\begin{align}
\label{eq:metric_evolution}
\partial_T g_{ij} = -\frac{2 \varphi}{\tau} N \Sigma_{ij} - 2\left(1 - \frac{N}{n}\right) g_{ij},
\end{align}
where $N$ is the lapse function, $\Sigma_{ij}$ is the transverse-traceless part of the second fundamental form, $\tau$ is the mean curvature, and $\varphi$ is a smooth weight function. For a fixed tangent vector $Y \in T_x M$, we obtain from \eqref{eq:metric_evolution}
\begin{align}
\label{eq:scalar_evolution}
\partial_T g(Y,Y) = -\frac{2\varphi}{\tau} N \Sigma(Y,Y) - 2\left(1 - \frac{N}{n} \right) g(Y,Y).
\end{align}
To estimate the right-hand side of \eqref{eq:scalar_evolution}, we begin with the bound
\begin{align}
\label{eq:sigma_pointwise}
|\Sigma(Y,Y)| \leq \sqrt{g(T_0,x)^{ik} g(T_0,x)^{jl} \Sigma_{ij} \Sigma_{kl}} \cdot g(T_0,x)(Y,Y),
\end{align}
since $\Sigma$ is a symmetric tensor and $g(T_0,x)$ is a fixed reference metric. Thus, using the bootstrap assumption \eqref{eq:strap2}, we have
\[
|\Sigma(Y,Y)| \lesssim \|\Sigma\|_{L^\infty(M)} \cdot g(T_0,x)(Y,Y).
\]
Now, using \eqref{eq:scalar_evolution}, we integrate from $T_0$ to $T$:
\begin{align}
|g(T,x)(Y,Y) - g(T_0,x)(Y,Y)| &\leq \int_{T_0}^T \left| \partial_s g(s,x)(Y,Y) \right| ds \\
&\lesssim \int_{T_0}^T \left( \left|\frac{\varphi(s)}{\tau(s)}\right| \cdot |N \Sigma(Y,Y)| + \left|1 - \frac{N}{n}\right| g(Y,Y) \right) ds.
\end{align}
From the bootstrap assumption, we have the decay bounds:
\[
\left| \frac{\varphi(s)}{\tau(s)} \right| \lesssim e^{-s}, \quad \|\Sigma\|_{L^\infty(M)} \lesssim e^{-\gamma s}, \quad \left\| 1 - \frac{N}{n} \right\|_{L^\infty(M)} \lesssim e^{-2\gamma s}.
\]
Therefore,
\begin{align}
|g(T,x)(Y,Y) - g(T_0,x)(Y,Y)| &\lesssim \int_{T_0}^T e^{-s} e^{-\gamma s} g(T_0,x)(Y,Y) ds + \int_{T_0}^T e^{-\gamma s} g(T_0,x)(Y,Y) ds \\
&\lesssim \left( \int_{T_0}^\infty e^{-(1+\gamma) s} ds + \int_{T_0}^\infty e^{-2\gamma s} ds \right) g(T_0,x)(Y,Y).
\end{align}
This yields
\[
|g(T,x)(Y,Y) - g(T_0,x)(Y,Y)| \lesssim e^{-2\gamma T_0} g(T_0,x)(Y,Y).
\]
To improve this to an estimate of relative deviation, we write
\[
\left| \frac{g(T,x)(Y,Y)}{g(T_0,x)(Y,Y)} - 1 \right| \lesssim e^{-2\gamma T_0}.
\]
Taking the supremum and infimum over unit vectors $Y$ with respect to $g(T_0,x)$ as in \eqref{eq:eigdef}, we obtain
\[
|\alpha(x,T) - 1| + |\beta(x,T) - 1| \lesssim \mathds{L} e^{-2\gamma T_0}.
\]
Since the estimates are uniform in $x$, the result follows.
\end{proof}

\begin{remark}
Propositions~\ref{metric1} and~\ref{metric2} establish uniform two-sided bounds on the evolving spatial metric $g(T,x)$ in terms of the initial metric $g_0 := g(T_0,x)$, valid on the entire bootstrap interval $T \in [T_0, T_{\infty}]$, with $T_0 \gg 1$ sufficiently large. More precisely, there exists a numerical constant $C \geq 1$, such that
\begin{equation}
    C^{-1} g_0 \leq g(T,x) \leq C g_0,
\end{equation}
in the sense of positive-definite symmetric bilinear forms. In particular, the evolving metric remains uniformly equivalent to the initial metric along the foliation. This uniform equivalence enables the definition of time-dependent Sobolev norms $\|\cdot\|_{H^s(g(T))}$ for tensor fields on the spatial slices $M_{T}$, and ensures that all Sobolev inequalities, elliptic estimates, derived from the geometry of $g_0$ continue to hold with constants depending only on $C$. This equivalence plays a critical role in establishing quantitative a priori estimates.
\end{remark}
\noindent We prove the estimates on the Sobolev constants in the following section.

\subsection{Controlling the Sobolev Constants}
\label{sobolev2}
\noindent One of the fundamental aspects of this problem is to control the appropriate Sobolev norms of the solution variables ($\Sigma,\mathfrak{T},N)$ for large times. In addition, one needs to make use of the Sobolev inequalities for the tensor fields on a manifold. Recall that the Sobolev norms are defined with respect to the dynamical metric $g$. This metric verifies the transport equation \ref{eq:cmc1}. Therefore, one needs to make sure that the metric $g$ does not degenerate during the evolution. Using the estimates derived on the metric components, we define the isoperimetric constant for the Cauchy slice
\begin{align}
 \text{I}(M):=\sup_{U\subset M,\partial U\in C^{1}}\frac{\min\{\text{vol}(U),\text{vol}(U^{c})\}}{[\text{Area}(\partial U)]^{\frac{3}{2}}}.   
\end{align}
The following proposition yields an upper bound for $I(M)$ 
\begin{proposition}\label{prop:isoperimetric-control}
Let $(\mathcal{M},g)$ be a globally hyperbolic spacetime endowed with a
smooth time function $T\colon \mathcal{M}\to\mathbb{R}$ whose level sets 
\[
M_{T}:=\{T=\mathrm{const}\}
\]
form a foliation by compact $2$-dimensional Riemannian manifolds 
(with induced metric still denoted $g(T)$).  
Assume the initial data on $M_{T_{0}}$ satisfy the hypotheses of Section~2
and that the bootstrap assumption \emph{(2.10)} holds along the entire slab 
$T\in[T_{0},T_{\infty})$.
Then for every $T\in [T_{0},T_{\infty})$ the isoperimetric constant satisfies
\begin{equation}\label{eq:iso-bound}
     I(M_{T}) \;\leq\; 10\,I(M_{T_{0}}).
\end{equation}
\end{proposition}

\begin{proof}
Fix any $T\in[T_{0},T_{\infty})$ and let $U\subset M_{T}$ be an arbitrary Caccioppoli
set whose boundary $\partial U$ is of class $C^{1}$.  
Let $\Phi_{s}$ denote the flow generated by the future-directed vector 
field $\partial_{T}$, so that 
\(
\Phi_{T-T_{0}}\colon M_{T_{0}}\to M_{T}
\)
is a diffeomorphism.
Set
\[
      U_{0}:=\Phi_{-(T-T_{0})}(U)\subset M_{T_{0}},
      \qquad 
      U^{c}_{0}:=\Phi_{-(T-T_{0})}(U^{c}).
\]
By global hyperbolicity, the integral curves of $\partial_{T}$
intersect each $M_{T}$ exactly once, so the map
$\Phi_{T-T_{0}}$ is globally well-defined and smooth.

\medskip
\noindent Let $\nu$ denote the outward unit normal to $\partial U$ in $(M_{T},g(T))$ and 
let $\nu_{0}$ denote the outward unit normal to $\partial U_{0}$ 
in $(M_{T_{0}},g(T_{0}))$.  
By Proposition~\ref{metric2}, the differential 
$D\Phi_{T-T_{0}}|_{T_{x}(\partial U_{0})}$ satisfies the uniform
ellipticity bounds
\begin{equation}\label{eq:area-distortion}
\beta(T)
\;\le\;
\frac{d\sigma_{g(T)}}{d\sigma_{g(T_{0})}}
\Big|_{\partial U}
\;\le\;
\alpha(T),
\end{equation}
where $\beta(T)>0$ and $\alpha(T)<\infty$ are the 
eigenvalues as in Proposition~\ref{metric2} and verify the uniform estimates in the same proposition \ref{metric2}.
In particular,
\begin{equation}\label{eq:area-ineq}
   \mathrm{Area}_{g(T)}(\partial U)
   \;\ge\;
   \beta(T)\,
   \mathrm{Area}_{g(T_{0})}(\partial U_{0}).
\end{equation}

\medskip
\noindent Similarly, Proposition~\ref{metric2} implies pointwise bounds for the 
Riemannian volume form:
\begin{equation}\label{eq:vol-distortion}
     \beta(T)^{\frac{3}{2}}
     \;\le\;
     \frac{d\mathrm{vol}_{g(T)}}{d\mathrm{vol}_{g(T_{0})}}
     \;\le\;
     \alpha(T)^{\frac{3}{2}},
\end{equation}
where $\alpha(T),\beta(T)$ are likewise controlled 
by \ref{metric2}.  
Integrating \eqref{eq:vol-distortion} over $U_{0}$ and $U^{c}_{0}$ yields
\begin{equation}\label{eq:vol-ineq}
   \mathrm{Vol}_{g(T)}(U)
       \;\le\; \alpha(T)^{\frac{3}{2}}\,\mathrm{Vol}_{g(T_{0})}(U_{0}), 
   \qquad
   \mathrm{Vol}_{g(T)}(U^{c})
       \;\le\; \alpha(T)^{\frac{3}{2}}\,\mathrm{Vol}_{g(T_{0})}(U^{c}_{0}).
\end{equation}

\medskip
\noindent By definition of the isoperimetric constant on $M_{T}$,
\[
   I(M_{T})
   =\sup_{U\subset M,\partial U\in C^{1}}\frac{\min\{\text{vol}(U),\text{vol}(U^{c})\}}{[\text{Area}(\partial U)]^{\frac{3}{2}}}.
\]
Using \eqref{eq:area-ineq} and \eqref{eq:vol-ineq} gives
\[
   I(M_{T})
   \leq \frac{\alpha(T)^{\frac{3}{2}}}{\beta(T)^{\frac{3}{2}}}\sup_{U\subset M,\partial U_{0}\in C^{1}}\frac{\min\{\text{vol}(U_{0}),\text{vol}(U^{c}_{0})\}}{[\text{Area}(\partial U_{0})]^{\frac{3}{2}}}=\frac{\alpha(T)^{\frac{3}{2}}}{\beta(T)^{\frac{3}{2}}}I(M_{T_{0}}).
\]
\medskip
Now in light of the proposition \ref{metric2}, the ratio $
      \frac{\alpha(T)^{\frac{3}{2}}}{\beta(T)^{\frac{3}{2}}}$
remains contained in the interval \([1/10,\, 10]\) for all 
$T\in[T_{0},T_{\infty})$.  
Hence, $I(M_{T})\;\leq\;10\,I(M_{T_{0}})$,
which establishes the desired estimate \eqref{eq:iso-bound}.
\end{proof}

  

\begin{proposition}[Sobolev Embedding for Tensor Fields]
\label{sobolev}
Let $(M,g)$ be a closed (i.e., compact without boundary), $C^{\infty}$ $3$-dimensional Riemannian manifold, and let $\psi \in C^{\infty}(\Gamma({}^{K}TM \otimes {}^{L}T^{*}M))$ be a smooth section of a mixed tensor bundle over $M$. Then the following Sobolev-type inequalities hold:
\begin{align}
\label{eq:sobolev-L4}
[\text{Vol}_{g}(M)]^{-\frac{1}{12}}||\psi||_{L^{4}(M,g)}\leq C\bigg(\max(1,I(M,g))\bigg)^{\frac{1}{2}}\bigg(||\nabla\psi||_{L^{2}(M,g)}+[\text{Vol}_{g}(M)]^{-\frac{1}{3}}||\psi||_{L^{2}(M,g)}\bigg),\\
\label{eq:sobolevL4}
 ||\psi||_{L^{\infty}(M,g)}\leq C\bigg(\max(1,I(M,g))\bigg)^{\frac{1}{2}}[\text{Vol}_{g}(M)]^{\frac{1}{3}-\frac{1}{4}}\bigg(||\nabla\psi||_{L^{4}(M,g)}+[\text{Vol}_{g}(M)]^{-\frac{1}{3}}||\psi||_{L^{4}(M,g)}\bigg),   
\end{align}
where the constants are purely numerical.
\end{proposition}

\begin{proof}
The proof proceeds via localization, partition of unity, and a scalar reduction argument adapted to tensor fields. First, since $M$ is a closed Riemannian manifold, it admits a finite atlas $\{(U_i, \varphi_i)\}_{i=1}^N$ such that each chart $\varphi_i : U_i \to \mathbb{R}^3$ is a diffeomorphism onto its image, and the pulled-back metric components $g_{ij}$ and their derivatives are uniformly bounded on each chart. Let $\{\chi_i\}_{i=1}^N$ be a smooth partition of unity subordinate to this cover.

\noindent Let $\psi \in C^\infty(\Gamma({}^{K}TM \otimes {}^{L}T^*M))$. Define for $\epsilon > 0$ the scalar function:
\[
f_\epsilon := \sqrt{|\psi|^2 + \epsilon},
\]
where $|\psi|^2 = g^{i_1 j_1} \cdots g^{i_K j_K} g_{a_1 b_1} \cdots g_{a_L b_L} \psi_{i_1 \cdots i_K}^{a_1 \cdots a_L} \psi_{j_1 \cdots j_K}^{b_1 \cdots b_L}$ is the pointwise norm squared induced by $g$.

\noindent Since the Levi-Civita connection $\nabla$ is compatible with $g$, we compute:
\[
\nabla f_\epsilon = \frac{1}{\sqrt{|\psi|^2 + \epsilon}} \langle \psi, \nabla \psi \rangle_g,
\]
where the inner product $\langle \psi, \nabla \psi \rangle_g$ is taken pointwise with respect to $g$. This yields the pointwise bound:
\[
|\nabla f_\epsilon| \leq \frac{|\psi| \cdot |\nabla \psi|}{\sqrt{|\psi|^2 + \epsilon}} \leq |\nabla \psi|.
\]

\noindent We now apply the standard scale-invariant Sobolev inequality for scalar functions on 3-dimensional Riemannian manifolds:
\[
[\text{Vol}_{g}(M)]^{-\frac{1}{12}}||f_{\epsilon}||_{L^{4}(M,g)}\leq C\bigg(\max(1,I(M,g))\bigg)^{\frac{1}{2}}\bigg(||\nabla f_{\epsilon}||_{L^{2}(M,g)}+[\text{Vol}_{g}(M)]^{-\frac{1}{3}}||f_{\epsilon}||_{L^{2}(M,g)}\bigg).
\]
Using the inequality $||f_\epsilon||_{L^p(M)} \leq \||\psi|\|_{L^p(M)} + \epsilon^{1/2}|M|^{1/p}$ and passing to the limit as $\epsilon \to 0$ yields
\[
[\text{Vol}_{g}(M)]^{-\frac{1}{12}}||\psi||_{L^{4}(M,g)}\leq C\bigg(\max(1,I(M,g))\bigg)^{\frac{1}{2}}\bigg(||\nabla\psi||_{L^{2}(M,g)}+[\text{Vol}_{g}(M)]^{-\frac{1}{3}}||\psi||_{L^{2}(M,g)}\bigg).
\]
This proves inequality \eqref{eq:sobolev-L4}.

\noindent 
We assume the scalar version of the scale-invariant Sobolev inequality:
\begin{equation}\label{scalar-Sobolev}
 ||f||_{L^{\infty}(M,g)}\leq C\bigg(\max(1,I(M,g))\bigg)^{\frac{1}{2}}[\text{Vol}_{g}(M)]^{\frac{1}{3}-\frac{1}{4}}\bigg(||\nabla f||_{L^{4}(M,g)}+[\text{Vol}_{g}(M)]^{-\frac{1}{3}}||f||_{L^{4}(M,g)}\bigg),
\end{equation}
for all smooth real-valued functions $f$. Let $\psi\in C^{\infty}(\Gamma({}^{K}TM\otimes {}^{L}T^{*}M))$ be a smooth mixed tensor field.
Set 
\[
f := |\psi|_{g},
\]
the pointwise norm induced by $g$. Then $f$ is a smooth nonnegative scalar function, and
\[
\|\psi\|_{L^{\infty}}=\|f\|_{L^{\infty}},\qquad
\|\psi\|_{L^{4}}=\|f\|_{L^{4}}.
\]
\noindent 
For tensor fields on a Riemannian manifold, one has the pointwise Kato inequality as seen in the proof of the first inequality (\ref{eq:sobolevL4})
\begin{equation}\label{Kato}
 |\nabla f|\le |\nabla\psi|.
\end{equation}
This follows from writing $f = (\psi_{\alpha}\psi_{\alpha}+\epsilon)^{1/2}$ in a local chart and
and applying the Cauchy-Schwarz inequality.
Thus 
\[
\|\nabla f\|_{L^{4}} \le \|\nabla\psi\|_{L^{4}}.
\]
Applying \eqref{scalar-Sobolev} to the scalar function $f=|\psi|$ gives
\[
\|\psi\|_{L^{\infty}}
=\|f\|_{L^{\infty}}
\le C\Big(\max(1,I(M,g))\Big)^{\frac12}
 \Big( Vol_{g}(M)^{\frac13}\|\nabla f\|_{L^{4}}
     + {Vol}_{g}(M)^{-\frac14}\|f\|_{L^{4}} \Big).
\]
\noindent Using Kato's inequality \eqref{Kato} and the equalities
$\|f\|_{L^{4}}=\|\psi\|_{L^{4}}$, 
$\|\nabla f\|_{L^{4}}\le \|\nabla\psi\|_{L^{4}}$,
we obtain
\[
\|\psi\|_{L^{\infty}}
\le C\Big(\max(1,I(M,g))\Big)^{\frac12}
\Big( {Vol}_{g}(M)^{\frac13}\|\nabla\psi\|_{L^{4}}
    + {Vol}_{g}(M)^{-\frac14}\|\psi\|_{L^{4}} \Big),
\]
which is exactly the desired tensorial inequality \eqref{eq:sobolevL4}.
\end{proof}

\begin{corollary}
For the dynamical manifold $(M_{T},g(T))$, the Sobolev inequalities \ref{eq:sobolev-L4}-\ref{eq:sobolevL4} hold with the constant $C\bigg(\max(1,I(M,g))\bigg)$ being uniformly bounded by $O(1)$ purely numerical entity $\forall T\in [T_{0},T_{\infty})$ in light of the proposition \ref{prop:isoperimetric-control}.    
\end{corollary}

\noindent Lastly, we note that on $(M,g)$ we can obtain compactness theorems for Sobolev embeddings. This can be proven by working on local charts using the uniform estimates on the metric $g(t)$ on the interval $[T_{0},T_{\infty}]$ in light of propositions \ref{metric1}-\ref{metric2} and gluing the charts together in a compatible manner. In particular, on $(M,g)$ for a $C^{\infty}$ tensor field $\psi$, the Sobolev norm $H^{I}(M)$ is defined as
\begin{align}
 ||\psi||_{H^{I}(M)}:=\sum_{j\leq I}||\nabla^{j}\psi||_{L^{2}(M)}.   
\end{align}
\begin{proposition}
The embedding $H^{I}(M)\hookrightarrow H^{J}(M)$ is compact for $I>J$.      
\end{proposition}
\begin{theorem}[Local Well-Posedness and Continuation]
\label{local}
Let \( g_0 \) be a Riemannian metric on a compact 3-manifold \( M \), satisfying the uniform ellipticity bounds
\[
C^{-1} \xi_0 \leq g_0 \leq C \xi_0
\]
for some smooth background metric \( \xi_0 \) and constant \( C > 0 \). Suppose the initial data \((\Sigma_0, \mathfrak{T}_0)\) belong to the Sobolev space \( H^I \times H^{I-1} \) for some \( I > \frac{5}{2} \), and satisfy the constraint equations \eqref{eq:HC11}-\eqref{eq:cmc2} associated with the re-scaled evolution system \eqref{eq:curvature1}--\eqref{eq:curvature2}, supplemented by the elliptic lapse equation \eqref{eq:lapse3} in constant mean curvature (CMC) transported spatial gauge.

\noindent Then there exists a time \( t^* > 0 \), depending only on \( C \), \( \| \Sigma_0 \|_{H^I} \), and \( \| \mathfrak{T}_0 \|_{H^{I-1}} \), such that the Cauchy problem admits a unique solution
\[
(\Sigma(t), \mathfrak{T}(t)) \in C\left([0, t^*]; H^I \times H^{I-1}\right),
\]
with lapse \( N(t) \in H^{I+2} \), and the solution map
\[
(\Sigma_0, \mathfrak{T}_0) \mapsto (\Sigma(t), \mathfrak{T}(t), N(t))
\]
is continuous as a map
\[
H^I \times H^{I-1} \to C\left([0, t^*]; H^I \times H^{I-1} \times H^{I+2}\right).
\]

\noindent Let \( t^{\max} \geq t^* \) denote the maximal time of existence of the solution. Then either \( t^{\max} = \infty \), or the solution breaks down at \( t^{\max} \) in the sense that
\[
\lim_{t \to t^{\max}} \sup \max \left( \mu_1(t)^{-1}, \mu_2(t)^{-1}, \mu_3(t)^{-1}, \| \Sigma(t) \|_{L^\infty},||\nabla(\frac{N}{n}-1)||_{L^{\infty}} \right) = \infty,
\]
where \( \{ \mu_i(t) \}_{i=1}^3 \) denote the eigenvalues of the symmetric endomorphism \( \xi_0^{-1} g(t) \).
\end{theorem}



\label{intermediate}
\noindent  In this section, we derive uniform estimates for the spatial Riemann curvature tensor $\operatorname{Riem}_{ijkl}$ on each constant mean curvature (CMC) slice, under the bootstrap assumptions introduced in Section~\ref{bootstrapargument}. Our goal is to express $\operatorname{Riem}_{ijkl}$ entirely in terms of the renormalized dynamical variables, namely the trace-free part of the second fundamental form $\Sigma$ and the auxiliary tensor $\mathfrak{T}$, and to obtain estimates compatible with the Sobolev regularity used in our energy framework.

\noindent To this end, we recall that in three spatial dimensions, the Riemann curvature tensor admits the following Ricci decomposition:
\begin{align} \label{eq:ricci-decomposition}
\operatorname{Riem}_{ijkl}
= \frac{1}{2}\left( g_{ik} \operatorname{Ric}_{jl} + g_{jl} \operatorname{Ric}_{ik} - g_{il} \operatorname{Ric}_{jk} - g_{jk} \operatorname{Ric}_{il} \right)
+ \frac{1}{6} R(g) \left( g_{il} g_{jk} - g_{ik} g_{jl} \right),
\end{align}
where $\operatorname{Ric}_{ij}$ and $R(g)$ denote, respectively, the Ricci tensor and scalar curvature of the induced Riemannian metric $g$ on the CMC slice.

\noindent The Ricci tensor can be expressed algebraically in terms of the renormalized momentum variable $\mathfrak{T}$ via the relation
\begin{align} \label{eq:ricci-mathfrakT}
\operatorname{Ric}_{ij} = \mathfrak{T}_{ij} + \frac{1}{3}R(g) g_{ij},
\end{align}
Furthermore, the scalar curvature satisfies the renormalized Hamiltonian constraint, which reads
\begin{align} \label{eq:scalar-renorm}
R(g) + \frac{n-1}{n} = |\Sigma|_g^2,
\end{align}
where $|\Sigma|_g^2 = g^{ia}g^{jb} \Sigma_{ij} \Sigma_{ab}$ denotes the pointwise squared norm of the symmetric trace-free tensor $\Sigma$.

\noindent Substituting~\eqref{eq:ricci-mathfrakT} and~\eqref{eq:scalar-renorm} into~\eqref{eq:ricci-decomposition}, we conclude that the full spatial Riemann curvature tensor $\operatorname{Riem}_{ijkl}$ is entirely determined by the tensors $\mathfrak{T}$ and $\Sigma$, together with the metric $g$. That is,
\begin{align} \label{eq:riem-sigma-mathfrakT}
\operatorname{Riem}_{ijkl} = \mathcal{F}_{ijkl}(g,\Sigma,\mathfrak{T}),
\end{align}
where $\mathcal{F}_{ijkl}$ denotes a universal algebraic expression involving the metric, its inverse, and the tensors $\mathfrak{T}$ and $\Sigma$. In particular, any Sobolev or pointwise bound on $\mathfrak{T}$ and $\Sigma$ yields a corresponding bound on the full curvature tensor.

\noindent This observation is central to the curvature estimates employed throughout the remainder of the analysis. In particular, it ensures that the control of geometric quantities such as $\operatorname{Riem}$ reduces to the control of the dynamical variables governed by the Einstein evolution equations in the chosen CMC-transported spatial gauge.
 \footnote{In $n\geq 4$, one needs to estimate the additional Weyl curvature as well. In such case, one first imposes bootstrap assumption on Weyl curvature and then improves the bootstrap by means of  separate energy estimates for the Weyl curvature}.
\subsection{Elliptic estimates}
\label{ellipticestimates}
\noindent Recall the elliptic equation 
\begin{align}
-\Delta_{g}(\frac{N}{n}-1)+(|\Sigma|^{2}+\frac{1}{3})(\frac{N}{n}-1)=-|\Sigma|^{2}. 
\end{align}
We obtain the following estimate for the normalized lapse function $\frac{N}{n}-1$
\begin{proposition}\label{elliptic}
Let $(M,g)$ be a compact Riemannian $3$-manifold with metric $g$ satisfying the bootstrap assumption \eqref{eq:strap2}, and let $\Sigma$ denote a symmetric transverse-traceless (TT) $2$-tensor on $M$. Consider the lapse function $N$ solving the elliptic equation
\[
-\Delta_g \left(\frac{N}{n} - 1\right) + \left(|\Sigma|_g^2 + \frac{1}{3}\right) \left(\frac{N}{n} - 1\right) = - |\Sigma|_g^2,
\]
where $n=3$. Then for any integer $I \in \{2,3\}$, the following Sobolev estimate holds:
\begin{align}\label{eq:lapse_sobolev_estimate}
\sum_{j=0}^{I+2} \|\nabla^j\big(\frac{N}{n} - 1\big)\|_{L^2(M)} 
\lesssim (1 + \Gamma) \sum_{j=0}^I \|\nabla^j \Sigma\|_{L^2(M)}^2,
\end{align}
where $\nabla$ denotes the Levi-Civita connection of $g$ and $\Gamma$ controls curvature norms as specified in the bootstrap regime.
\end{proposition}
\begin{remark}
Note the loss of decay of $(\frac{N}{n}-1)$ at the top-most order due to loss of decay of $\Sigma$ at the top-most order.     
\end{remark}
\begin{proof}
\noindent
Under the bootstrap assumption~\eqref{eq:strap2}, we establish energy estimates for the spatial Riemann curvature tensor $\operatorname{Riem}$ of the Riemannian metric $g$ on the spatial slice $M$. More precisely, for all integers $0 \leq I \leq 2$, we claim the following uniform bound:
\begin{align} \label{eq:Riem_bound_claim}
    \|\nabla^I \operatorname{Riem} \|_{L^2(M)} \lesssim \Gamma + 1.
\end{align}

\noindent
To derive this, we exploit the classical Ricci decomposition of the Riemann tensor in dimension $n = 3$, where the Weyl tensor vanishes identically:
\begin{align}
    \operatorname{Riem}_{ijkl} = \tfrac{1}{2} \left( g_{ik} \operatorname{Ric}_{jl} + g_{il} \operatorname{Ric}_{kj} - g_{jk} \operatorname{Ric}_{il} - g_{jl} \operatorname{Ric}_{ik} \right) 
    + \tfrac{1}{6} R(g) \left( g_{il} g_{jk} - g_{ik} g_{jl} \right).
\end{align}
\noindent
From this decomposition, it follows that the full Riemann tensor is algebraically expressible in terms of the Ricci tensor and scalar curvature. Invoking the renormalized Hamiltonian constraint, we write:
\begin{align}
    R(g) = -\tfrac{n-1}{n} + |\Sigma|^2,
\end{align}
\noindent
where $\Sigma$ denotes the trace-free part of the second fundamental form in CMC gauge, and where all norms and derivatives are computed with respect to the spatial metric $g$.

\vspace{1ex}
\noindent
Let us now provide schematic bounds for $\nabla^I \operatorname{Riem}$ in terms of the fundamental dynamical fields $\mathfrak{T}$ and $\Sigma$. Using the expressions for $\operatorname{Ric}$ and $R(g)$, we obtain:
\begin{align}
    |\nabla \operatorname{Riem}| 
    &\lesssim |\nabla \mathfrak{T}| + |\nabla R(g)| 
    \lesssim |\nabla \mathfrak{T}| + |\Sigma||\nabla \Sigma|, \\
    |\nabla^2 \operatorname{Riem}| 
    &\lesssim |\nabla^2 \mathfrak{T}| + |\nabla \Sigma|^2 + |\Sigma||\nabla^2 \Sigma|.
\end{align}

\noindent
By applying the Cauchy–Schwarz inequality and standard Sobolev embedding on the compact manifold $M$, together with the bootstrap assumption and bounds derived in Propositions~\ref{metric1}--\ref{metric2}, we deduce:
\begin{align}
    \| \nabla \operatorname{Riem} \|_{L^2(M)} 
    &\lesssim \| \nabla \mathfrak{T} \|_{L^2(M)} 
    + \| \nabla \Sigma \|_{L^2(M)} \| \nabla^2 \Sigma \|_{L^2(M)} 
    \lesssim \Gamma + e^{-2\gamma T} \mathds{Y}^2 
    \lesssim \Gamma + 1, \\
    \| \nabla^2 \operatorname{Riem} \|_{L^2(M)} 
    &\lesssim \| \nabla^2 \mathfrak{T} \|_{L^2(M)} + \| \nabla^2 \Sigma \|_{L^2(M)}^2 
    \lesssim \Gamma + e^{-2\gamma T} \mathds{Y}^2 
    \lesssim \Gamma + 1.
\end{align}

\noindent
Finally, the $L^2$ bound for $\operatorname{Riem}$ without derivatives follows directly from the pointwise Ricci decomposition and the bootstrap bound on $\mathfrak{T}$ and $\Sigma$:
\begin{align}
    \| \operatorname{Riem} \|_{L^2(M)} \lesssim \Gamma + 1.
\end{align}

\noindent
Thus, we conclude that all desired curvature norms up to three derivatives are uniformly controlled by the bootstrap parameter $\Gamma$, thereby validating the estimate~\eqref{eq:Riem_bound_claim}. Now we are ready to obtain the estimates for $\frac{N}{n}-1$. Let us recall the elliptic equation 
\begin{align}
-\Delta_{g}(\frac{N}{n}-1)+(|\Sigma|^{2}+\frac{1}{3})(\frac{N}{n}-1)=-|\Sigma|^{2}.    
\end{align}
Let us denote the operator $-\Delta_{g}+(|\Sigma|^{2}+\frac{1}{3})$ by $\mathscr{L}$. We claim that an estimate of the following type holds for $u\in C^{\infty}(M)$
\begin{align}
\sum_{j=0}^{I+2}||\nabla^{j}u||_{L^{2}(M)}\lesssim \sum_{j=0}^{I}||\nabla^{j}(\mathscr{L}u)||_{L^{2}(M)}
\end{align}
with the regularity $\nabla^{I-1}\text{Riem}\in L^{2}(M)$. To do this we commute the lapse equation with $\nabla^{I}$ for $0\leq I\leq 3$
\begin{align}
\mathscr{L}\nabla^{I}(\frac{N}{n}-1)=-\Delta_{g}\nabla^{I}(\frac{N}{n}-1)+(|\Sigma|^{2}+\frac{1}{3})\nabla^{I}(\frac{N}{n}-1)\nonumber=\nabla^{I}\mathscr{L}(\frac{N}{n}-1)\\\nonumber+\sum_{J_{1}+J_{2}+J_{3}=I-1}\nabla^{J_{1}+1}\Sigma\nabla^{J_{2}}\Sigma\nabla^{J_{3}}(\frac{N}{n}-1)+\sum_{m=0}^{I-1}\sum_{J_{1}+J_{2}=m}\nabla^{J_{1}}\text{Ric}\nabla^{J_{2}+I-m}(\frac{N}{n}-1)\\\nonumber+\sum_{m=0}^{I-1}\sum_{J_{1}+J_{2}=m}\nabla^{J_{1}}\text{Riem}\nabla^{J_{2}+I-m}(\frac{N}{n}-1)+\sum_{m=0}^{I-2}\sum_{J_{1}+J_{2}=m}\nabla^{J_{1}+1}\text{Riem}\nabla^{J_{2}+I-m-1}(\frac{N}{n}-1).
\end{align}
The elliptic regularity implies 
\begin{align}
\nonumber||\nabla^{I+2}(\frac{N}{n}-1)||_{L^{2}(M)}\lesssim ||\nabla^{I}\mathscr{L}(\frac{N}{n}-1)||_{L^{2}(M)}\\\nonumber+||(|\Sigma|^{2}+\frac{1}{3})\nabla^{I}(\frac{N}{n}-1)||_{L^{2}(M)}+||\sum_{J_{1}+J_{2}+J_{3}=I-1}\nabla^{J_{1}+1}\Sigma\nabla^{J_{2}}\Sigma\nabla^{J_{3}}(\frac{N}{n}-1)||_{L^{2}(M)}\\\nonumber+||\sum_{m=0}^{I-1}\sum_{J_{1}+J_{2}=m}\nabla^{J_{1}}\text{Ric}\nabla^{J_{2}+I-m}(\frac{N}{n}-1)||_{L^{2}(M)}+||\sum_{m=}^{I-1}\sum_{J_{1}+J_{2}=m}\nabla^{J_{1}}\text{Riem}\nabla^{J_{2}+I-m}(\frac{N}{n}-1)||_{L^{2}(M)}\\\nonumber+||\sum_{m=0}^{I-2}\sum_{J_{1}+J_{2}=m}\nabla^{J_{1}+1}\text{Riem}\nabla^{J_{2}+I-m-1}(\frac{N}{n}-1)||_{L^{2}(M)}.    
\end{align}
Under the boot-strap assumption \ref{eq:strap2}, we estimate each term separately. First estimate
\begin{align}
||(|\Sigma|^{2}+\frac{1}{3})\nabla^{I}(\frac{N}{n}-1)||_{L^{2}(M)}\lesssim \nonumber||\nabla^{I}(\frac{N}{n}-1)||_{L^{2}(M)}(1+e^{-2\gamma T}\mathds{Y}^{2})\lesssim ||\nabla^{I}(\frac{N}{n}-1)||_{L^{2}(M)}.    
\end{align}
The next term reads 
\begin{align}
 ||\sum_{J_{1}+J_{2}+J_{3}=I-1}\nabla^{J_{1}+1}\Sigma\nabla^{J_{2}}\Sigma\nabla^{J_{3}}(\frac{N}{n}-1)||_{L^{2}(M)}\lesssim e^{-\gamma T}\mathds{Y}\Gamma ||\nabla^{I}(\frac{N}{n}-1)||_{L^{2}(M)}\\\nonumber\lesssim ||\nabla^{I}(\frac{N}{n}-1)||_{L^{2}(M)}.   
\end{align}
for the following terms we have 
\begin{align}
||\sum_{m=0}^{I-1}\sum_{J_{1}+J_{2}=m}\nabla^{J_{1}}\text{Ric}\nabla^{J_{2}+I-m}(\frac{N}{n}-1)||_{L^{2}(M)}\lesssim \Gamma ||\nabla^{I}(\frac{N}{n}-1)||_{L^{2}(M)},\\\nonumber 
||\sum_{m=}^{I-1}\sum_{J_{1}+J_{2}=m}\nabla^{J_{1}}\text{Riem}\nabla^{J_{2}+I-m}(\frac{N}{n}-1)||_{L^{2}(M)}\lesssim \Gamma ||\nabla^{I}(\frac{N}{n}-1)||_{L^{2}(M)},\\\nonumber
||\sum_{m=0}^{I-2}\sum_{J_{1}+J_{2}=m}\nabla^{J_{1}+1}\text{Riem}\nabla^{J_{2}+I-m-1}(\frac{N}{n}-1)||_{L^{2}(M)}\lesssim \Gamma ||\nabla^{I}(\frac{N}{n}-1)||_{L^{2}(M)}.
\end{align}
Therefore collecting all the terms together, we conclude 
\begin{align}
\label{eq:iteratenew}
||(\frac{N}{n}-1)||_{H^{I+2}(M)}\lesssim ||\mathscr{L}(\frac{N}{n}-1)||_{H^{I}(M)}+(1+\Gamma)||(\frac{N}{n}-1)||_{H^{I}(M)}.    
\end{align}
Now we prove the desired estimates in an iterative way starting from the first-order estimate. First, recall the elliptic equation 
\begin{align}
-\Delta_{g}(\frac{N}{n}-1)+(|\Sigma|^{2}+\frac{1}{3})(\frac{N}{n}-1)=-|\Sigma|^{2}.     
\end{align}
Now multiply both sides by $\frac{N}{n}-1$ and integrate by parts 
\begin{align}
\label{eq:intby}
 \int_{M}|\nabla(\frac{N}{n}-1)|^{2}\mu_{g}+\int_{M}(|\Sigma|^{2}+\frac{1}{3})(\frac{N}{n}-1)^{2}\mu_{g}=-\int_{M}|\Sigma|^{2}(\frac{N}{n}-1)\mu_{g}.   
\end{align}
Now we estimate the term on the right-hand side by Cauchy-Schwartz as follows 
\begin{align}
-\int_{M}|\Sigma|^{2}(\frac{N}{n}-1)\mu_{g}\leq (n+1)\int_{M}|\Sigma|^{4}\mu_{g}+\frac{1}{n+1}\int_{M}(\frac{N}{n}-1)^{2}\mu_{g}.    
\end{align}
Therefore, \ref{eq:intby} reads 
\begin{align}
 \int_{M}|\nabla(\frac{N}{n}-1)|^{2}\mu_{g}+\int_{M}|\Sigma|^{2}(\frac{N}{n}-1)^{2}\mu_{g}\nonumber+(\frac{1}{3}-\frac{1}{n+1})\int_{M}(\frac{N}{n}-1)^{2}\mu_{g}\leq (n+1)\int_{M}|\Sigma|^{4}\mu_{g}
\end{align}
which by means of the inequality \ref{sobolev} yields 
\begin{align}
  \int_{M}|\nabla(\frac{N}{n}-1)|^{2}\mu_{g}+\int_{M}|\Sigma|^{2}(\frac{N}{n}-1)^{2}\mu_{g}\nonumber+\frac{1}{n(n+1)}\int_{M}(\frac{N}{n}-1)^{2}\mu_{g} \\\nonumber 
  \lesssim \left(\int_{M}|\nabla\Sigma|^{2}\mu_{g}+\int_{M}|\Sigma|^{2}\mu_{g}\right)^{2}.
\end{align}
More concretely, we have the two lowest-order estimates 
\begin{align}
||\frac{N}{n}-1||_{L^{2}(M)}\lesssim ||\Sigma||^{2}_{H^{1}(M)},\\
||\frac{N}{n}-1||_{H^{1}(M)}\lesssim ||\Sigma||^{2}_{H^{1}(M)}.
\end{align}
Now iterate this using \ref{eq:iteratenew} to obtain 
\begin{align}
||\frac{N}{n}-1||_{H^{2}(M)}\lesssim |||\Sigma|^{2}||_{L^{2}(M)}+(1+\Gamma)||\frac{N}{n}-1)||_{L^{2}(M)}\lesssim |||\Sigma|^{2}||_{L^{2}(M)}\nonumber+(1+\Gamma)||\Sigma||^{2}_{H^{1}(M)},\\\nonumber
||\frac{N}{n}-1||_{H^{3}(M)}\lesssim |||\Sigma|^{2}||_{H^{1}(M)}+(1+\Gamma)||\frac{N}{n}-1)||_{H^{1}(M)}\lesssim |||\Sigma|^{2}||_{H^{1}(M)}+(1+\Gamma)||\Sigma||^{2}_{H^{1}(M)},\\\nonumber
||\frac{N}{n}-1||_{H^{4}(M)}\lesssim |||\Sigma|^{2}||_{H^{2}(M)}+(1+\Gamma)||\frac{N}{n}-1)||_{H^{2}(M)}\lesssim |||\Sigma|^{2}||_{H^{2}(M)}+(1+\Gamma)||\Sigma||^{2}_{H^{1}(M)}
\end{align}
and so on. ow use the embedding $H^{i_{1}}(M)\hookrightarrow H^{i_{2}}(M)$ for $i_{1}>i_{2}$, and the algebra property of the Sobolev spaces $H^{s}(M)$ for $s>\frac{n}{2}$ or that of $L^{\infty}(M)\cap H^{s}(M)$ for $s\geq 0$, we conclude the desired estimate 
\begin{align}
 \sum_{j=0}^{I+2}||\nabla^{j}(\frac{N}{n}-1)||_{L^{2}(M)}\lesssim (1+\Gamma)\sum_{j=0}^{I}||\nabla^{j}\Sigma||^{2}_{L^{2}(M)},~I>2.   
\end{align}


\noindent The elliptic estimate yields, together with the algebra property of Sobolev spaces for $I>2$
\begin{align}
 ||\frac{N}{n}-1||_{H^{I+2}(M)}\lesssim ||\Sigma||^{2}_{H^{I}(M)}
\end{align}
%For lower order case, the following decay estimate is obtained by 
% \begin{align}
% ||\frac{N}{n}-1||_{H^{I+2}(M)}\lesssim ||\Sigma||^{2}_{H^{I}(M)}\lesssim e^{-2\gamma T}\mathds{Y}^{2},~for~I=3.
%\end{align}
\end{proof}
%\begin{remark}
%Instead of using the elegant argument utilizing the compact imbedding and injectivity of the operator $\mathscr{L}$, one can also use the rough estimate $||(\frac{N}{n}-1)||_{H^{I+2}(M)}\lesssim ||\mathscr{L}(\frac{N}{n}-1)||_{H^{I}(M)}+(1+\Gamma)||(\frac{N}{n}-1)||_{H^{I}(M)}$ together with an iterative elliptic estimate starting with $I=0$ to prove $||(\frac{N}{n}-1)||_{H^{I+2}(M)}\lesssim ||\mathscr{L}(\frac{N}{n}-1)||_{H^{I}(M)}+(1+\Gamma)e^{-2\gamma T}$. This estimate too suffices for our purpose too.       
%\end{remark}
%Let us commute $\nabla_{i}$ 
%\begin{align}
%-\nabla_{i}\nabla_{m}\nabla^{m}(\frac{N}{n}-1)+n(|\Sigma|^{2}+\frac{1}{3})\nabla_{i}(\frac{N}{n}-1)=-n\langle \Sigma,\nabla_{i}\Sigma\rangle+n(\frac{N}{n}-1)\langle \Sigma,\nabla_{i}\Sigma\rangle   
%\end{align}
%or 
%\begin{align}
%\Delta_{g}\nabla_{i}(\frac{N}{n}-1)+ n(|\Sigma|^{2}+\frac{1}{3})\nabla_{i}(\frac{N}{n}-1)=R_{ij}\nabla^{j}(\frac{N}{n}-1)\nonumber-n\langle \Sigma,\nabla_{i}\Sigma\rangle+n(\frac{N}{n}-1)\langle \Sigma,\nabla_{i}\Sigma\rangle   
%\end{align}
%Now we obtain a recursion relation 
%\begin{align}
%\Delta_{g}\nabla^{I}(\frac{N}{n}-1)+n(|\Sigma|^{2}+\frac{1}{3})\nabla^{I}(\frac{N}{n}-1)=\nabla^{I-1}\left(R_{ij}\nabla^{j}(\frac{N}{n}-1)\nonumber-n\langle \Sigma,\nabla_{i}\Sigma\rangle+n(\frac{N}{n}-1)\langle \Sigma,\nabla_{i}\Sigma\rangle\right)\\\nonumber 
%-n\nabla^{I-1}(|\Sigma|^{2}+\frac{1}{3})\nabla_{i}(\frac{N}{n}-1)+[\Delta_{g},\nabla^{I-1}]\nabla_{i}(\frac{N}{n}-1)
%\end{align}


%Let us compute the commutator term. We write the multiple derivatives in schematic form. Define $u:=\frac{N}{n}-1$
%\begin{align}
 %   \gamma^{I}:=[\nabla_{i}\nabla^{i},\nabla^{I}]u\\\nonumber
  %  =[\nabla_{i}\nabla^{i},\nabla]\nabla^{I-1}u+\nabla[\nabla_{i}\nabla^{i},\nabla^{I-1}]u\\\nonumber 
   % =[\nabla_{i},\nabla]\nabla^{i}\nabla^{I-1}u+\nabla^{i}[\nabla_{i},\nabla]\nabla^{I-1}u+\nabla\gamma^{I-1}
%\end{align}
%Now note the computations 
%\begin{align}
 %[\nabla_{i},\nabla]\nabla^{i}\nabla^{I-1}u\sim \text{Ric}\nabla^{I}u+\text{Riem}\nabla^{I}u.   
%\end{align}
%Therefore we have the following relation 
%\begin{align}
 %\gamma^{I}=\alpha^{I}+\nabla\gamma^{I-1},   
%\end{align}
%where $\alpha^{I}:=\text{Ric}\nabla^{I}u+\text{Riem}\nabla^{I}u+\nabla(\text{Riem}\nabla^{I-1}u)$. Therefore 
%\begin{align}
 %\gamma^{I}=\alpha^{I}+\nabla\gamma^{I-1}=\alpha^{I}+\nabla(\alpha^{I-1}+\nabla\gamma^{I-2})=\sum_{m=0}^{I-1}\nabla^{m}\alpha^{I-m}   
%\end{align}
%and explicitly 
%\begin{align}
 %\sum_{m=0}^{I-1}\nabla^{m}\alpha^{I-m}=\sum_{m=0}^{I-1}\nabla^{m}\left(\text{Ric}\nabla^{I-m}u+\text{Riem}\nabla^{I-m}u+\nabla(\text{Riem}\nabla^{I-m-1}u)\right)\\\nonumber 
 %=\sum_{m=0}^{I-1}\sum_{J_{1}+J_{2}=m}\nabla^{J_{1}}\text{Ric}\nabla^{J_{2}+I-m}u+\sum_{m=}^{I-1}\sum_{J_{1}+J_{2}=m}\nabla^{J_{1}}\text{Riem}\nabla^{J_{2}+I-m}u\\\nonumber+\sum_{m=0}^{I-2}\sum_{J_{1}+J_{2}=m}\nabla^{J_{1}+1}\text{Riem}\nabla^{J_{2}+I-m-1}u
%\end{align}

\noindent The next part entails controlling the lower order norm $\mathcal{F}$ of $\Sigma$ in terms of the top order norm $\mathcal{O}$ and the initial data norm $\mathcal{I}^{0}$. 

\subsection{Estimating the lower order norm of TT tensor $\Sigma$ via transport equation for $\Sigma$}
\label{entropy}
\noindent In this section, we prove that the weighted lower order norm of the $TT$ tensor $\Sigma$ is controlled by the top order norm of $(\mathfrak{T},\Sigma)$ and the initial data. Recall the equation for $\Sigma$
\begin{align}
\frac{\partial \Sigma_{ij}}{\partial T}&=&-(n-1)\Sigma_{ij}-\frac{\varphi}{\tau}N\mathfrak{T}_{ij}+\frac{\varphi}{\tau}\nabla_{i}\nabla_{j}(\frac{N}{n}-1)+\frac{2\varphi}{\tau}N\Sigma_{ik}\Sigma^{k}_{j}\\
 &&\nonumber-\frac{\varphi}{n\tau}(\frac{N}{n}-1)g_{ij}-(n-2)(\frac{N}{n}-1)\Sigma_{ij}.    
\end{align}
To estimate the lower order norm $\mathcal{F}$, we will treat the equation for $\Sigma$ as a transport equation. 

\begin{definition}
\label{entropy2}
\begin{align} 
\mathscr{E}^{low}:=\frac{1}{2}\sum_{I\leq 2}\int_{M}|\nabla^{I}\Sigma|^{2}\mu_{g}
\end{align}
\end{definition}
\noindent First note that $\mathcal{F}^{2}\approx e^{2\gamma T}\mathscr{E}^{low}$. We have the following theorem estimating $\mathcal{F}^{2}$.
\begin{proposition}[Coercive control of the reversed entropy]\label{model2}
Let $(M^n,g(T))$ be a smooth solution of the rescaled Einstein vacuum equations with cosmological constant $\Lambda>0$ on a time interval $[T_0,T]$, written in the fixed gauge of this work.  
Let $\mathscr{E}^{\mathrm{low}}(T)\ge0$ denote the reversed entropy functional controlling the lower order norm of $\Sigma$ as defined in \ref{entropy2} and let $\mathcal{F}(T)$ be the associated re-scaled entity i.e., 
\begin{equation}\label{eq:equiv}
e^{2T}\,\mathscr{E}^{\mathrm{low}}(T)\ \lesssim \ \mathcal{F}(T)^2\ \lesssim\ e^{2T}\,\mathscr{E}^{\mathrm{low}}(T)
\end{equation}
and $\mathcal{O}$ be the top order norm \ref{eq:O}
Then the following estimate for the re-scaled reversed entropy holds
\begin{enumerate}
\item[{\normalfont(i)}] (\emph{A–priori upper bound}).  
\begin{equation}\label{eq:upper}
\mathcal{F}(T)^2\ \lesssim\ \Bigl(1+\bigl(\mathcal{I}^{0}\bigr)^{2}+\mathcal{O}^{2}\Bigr).
\end{equation}
\end{enumerate}
In particular, combining \eqref{eq:equiv} and \eqref{eq:upper} yields
\begin{equation}\label{eq:announced}
e^{2T}\,\mathscr{E}^{\mathrm{low}}(T)\ \lesssim\ \Bigl(1+\bigl(\mathcal{I}^{0}\bigr)^{2}+\mathcal{O}^{2}\Bigr),
\end{equation}
where the involved constants are of pure numerical nature i.e., dpendend on the dimension and independent of $T$. 
\end{proposition}

\begin{proof}
Recall definition \ref{entropy2}
 \begin{align}
 \mathscr{E}^{low}(T):=\frac{1}{2}\sum_{I\leq 2}\int_{M(T)}\langle\nabla^{I}\Sigma,\nabla^{I}\Sigma\rangle\mu_{g}.
\end{align}
Explicit application of the time evolution operator $\frac{d}{dT}$ together with the transport theorem \cite{marsden} yields
\begin{align}
 \frac{d\mathscr{E}^{low}}{dT}=\sum_{I\leq 2}\int_{M}\langle\nabla^{I}\partial_{T}\Sigma,\nabla^{I}\Sigma\rangle\mu_{g}+\sum_{I\leq 2}\int_{M}\langle[\partial_{T},\nabla^{I}]\Sigma,\nabla^{I}\Sigma\rangle\mu_{g}+\mathscr{ER}_{2},   
\end{align}
where the error terms $\mathscr{ER}_{2}$ reads schematically
\begin{align}
\mathscr{ER}_{2}\sim \sum_{I\leq 2}\int_{M(T)}\partial_{T}(g^{-1}....g^{-1}g.....g)\nabla^{I}\Sigma\nabla^{I}\Sigma\mu_{g}+\sum_{I\leq 2}\int_{M(T)}|\nabla^{I}\Sigma|^{2}\partial_{T}\mu_{g}.      
\end{align}
Using the equation for $\partial_{T}\Sigma$, $\frac{d\mathscr{E}^{low}}{dT}$ reads 
\begin{align}
\nonumber \frac{d\mathscr{E}^{low}}{dT}=-2(n-1)\mathscr{E}+\sum_{I\leq 2}\frac{\varphi}{\tau}\int_{M(T)}N \nabla^{I}\mathfrak{T}\nabla^{I}\Sigma\mu_{g}
+\sum_{I\leq 2} \sum_{m=0}^{I-1}\sum_{J_{1}+J_{2}=m}\int_{M(T)}\frac{\varphi}{\tau}\nabla^{J_{1}+1}(\frac{N}{n}-1)\nabla^{J_{2}+I-m-1}\Sigma\nabla^{I}\Sigma\mu_{g}\\\nonumber 
+\sum_{I\leq 2}\sum_{m=0}^{I-1}\sum_{J_{1}+J_{2}+J_{3}=m}\int_{M(T)}\frac{\varphi}{\tau}\nabla^{J_{1}}N\nabla^{J_{2}+1}\Sigma\nabla^{J_{3}+I-m-1}\Sigma\nabla^{I}\Sigma\mu_{g}
 +\sum_{I\leq 2}\sum_{m=0}^{I-1}\sum_{J_{1}+J_{2}=m}\int_{M(T)}\nabla^{J_{1}+1}(\frac{N}{n}-1)\nabla^{J_{2}+I-m-1}\Sigma\nabla^{I}\Sigma\mu_{g}\\\nonumber 
 +\frac{\varphi}{\tau}\sum_{I\leq 2}\sum_{J_{1}+J_{2}+J_{3}=I}\int_{M(T)}\nabla^{J_{1}}N\nabla^{J_{2}}\Sigma\nabla^{J_{3}}\Sigma\nabla^{I}\Sigma\mu_{g} 
 +\sum_{I\leq 2}\frac{\varphi}{\tau}\int_{M(T)}\nabla^{I+2}(\frac{N}{n}-1)\nabla^{I}\Sigma\mu_{g}\\\nonumber 
 +\sum_{I\leq 2}\sum_{J_{1}+J_{2}=I}\int_{M(T)}\nabla^{J_{1}}(\frac{N}{n}-1)\nabla^{J_{2}}\Sigma\nabla^{I}\Sigma\mu_{g} 
 +\sum_{I\leq 2}\frac{\varphi}{\tau}\int_{M}\sum_{J_{1}+J_{2}=I-1}\nabla^{J_{1}+1}(\frac{N}{n}-1)\nabla^{J_{2}}\mathfrak{T}\nabla^{I}\Sigma\mu_{g}+\mathscr{ER}_{2}.
\end{align}
Now we estimate each term separately. First, note that theorem \ref{topestimate} implies boundedness of $\sum_{I\leq 2}||\nabla^{I}\mathfrak{T}||_{L^{2}(M)}$ by $\mathcal{I}^{0}$. Therefore   
\begin{align}
 |\sum_{I\leq 2}\frac{\varphi}{\tau}\int_{M(T)}N \nabla^{I}\mathfrak{T}\nabla^{I}\Sigma\mu_{g}|\lesssim e^{-T}\mathcal{O}^{2}.   
\end{align}
The next term is estimated as follows 
\begin{align}
|\sum_{I\leq 2} \sum_{m=0}^{I-1}\sum_{J_{1}+J_{2}=m}\int_{M(T)}\frac{\varphi}{\tau}\nabla^{J_{1}+1}(\frac{N}{n}-1)\nabla^{J_{2}+I-m-1}\Sigma\nabla^{I}\Sigma\mu_{g}|   \lesssim e^{-(1+4\gamma)T}\mathds{Y}^{4}, 
\end{align}
\begin{align}
 |\sum_{I\leq 2}\sum_{m=0}^{I-1}\sum_{J_{1}+J_{2}+J_{3}=m}\int_{M(T)}\frac{\varphi}{\tau}\nabla^{J_{1}}N\nabla^{J_{2}+1}\Sigma\nabla^{J_{3}+I-m-1}\Sigma\nabla^{I}\Sigma\mu_{g}|\lesssim e^{-(1+3\gamma)T}\mathds{Y}^{3},   
\end{align}


\begin{align}
|\sum_{I\leq 2}\sum_{m=0}^{I-1}\sum_{J_{1}+J_{2}=m}\int_{M(T)}\nabla^{J_{1}+1}(\frac{N}{n}-1)\nabla^{J_{2}+I-m-1}\Sigma\nabla^{I}\Sigma\mu_{g}|\lesssim e^{-4\gamma T}\mathds{Y}^{4},    
\end{align}

\begin{align}
|\frac{\varphi}{\tau}\sum_{I\leq 2}\sum_{J_{1}+J_{2}+J_{3}=I}\int_{M(T)}\nabla^{J_{1}}N\nabla^{J_{2}}\Sigma\nabla^{J_{3}}\Sigma\nabla^{I}\Sigma\mu_{g}|\lesssim e^{-(1+3\gamma)T}\mathds{Y}^{3},
\end{align}

\begin{align}
|\sum_{I\leq 2}\frac{\varphi}{\tau}\int_{M(T)}\nabla^{I+2}(\frac{N}{n}-1)\nabla^{I}\Sigma\mu_{g}|\lesssim e^{-(1+3\gamma)T}\mathds{Y}^{3},
\end{align}

\begin{align}
 |\sum_{I\leq 2}\sum_{J_{1}+J_{2}=I}\int_{M(T)}\nabla^{J_{1}}(\frac{N}{n}-1)\nabla^{J_{2}}\Sigma\nabla^{I}\Sigma\mu_{g}| \lesssim e^{-4\gamma T} \mathds{Y}^{4}, 
\end{align}

\begin{align}
|\sum_{I\leq 2}\frac{\varphi}{\tau}\int_{M}\sum_{J_{1}+J_{2}=I-1}\nabla^{J_{1}+1}(\frac{N}{n}-1)\nabla^{J_{2}}\mathfrak{T}\nabla^{I}\Sigma\mu_{g}| \lesssim e^{-(1+3\gamma)T}\Gamma\mathds{Y}^{3}.   
\end{align}
Collection of all terms yields 
\begin{align}
\frac{d\mathscr{E}}{dT}+2(n-1)\mathscr{E}\lesssim e^{-T}\mathcal{O}^{2}+e^{-(1+4\gamma)T}\mathds{Y}^{4}+e^{-(1+3\gamma)T}\mathds{Y}^{3}\\\nonumber+e^{-4\gamma T}\mathds{Y}^{4}+e^{-(1+3\gamma)T}\Gamma\mathds{Y}^{3}
\end{align}
integration of which yields 
\begin{align}
\mathscr{E}(T) \lesssim\; & e^{-2(n-1)(T-T_{0})} \mathscr{E}(T_{0)} \notag \\
& + e^{-2(n-1)T} \left(e^{[2(n-1)-1]T} - e^{[2(n-1)-1]T_{0}}\right) \mathcal{O}^{2} \notag \\
& + e^{-2(n-1)T} \left(e^{[2(n-1)-(1+4\gamma)]T} - e^{[2(n-1)-(1+4\gamma)]T_{0}}\right) \mathds{Y}^{4} \notag \\
& + e^{-2(n-1)T} \left(e^{[2(n-1)-(1+3\gamma)]T} - e^{[2(n-1)-(1+3\gamma)]T_{0}}\right) \mathds{Y}^{3} \notag \\
& + e^{-2(n-1)T} \left(e^{[2(n-1)-4\gamma]T} - e^{[2(n-1)-4\gamma]T_{0}}\right) \mathds{Y}^{4} \notag \\
& + e^{-2(n-1)T} \left(e^{[2(n-1)-(1+3\gamma)]T} - e^{[2(n-1)-(1+3\gamma)]T_{0}}\right) \Gamma \mathds{Y}^{3}.
\end{align}

Alternatively, we can write:
\begin{align}
e^{2\gamma T} \mathscr{E}(T) \leq\; & \frac{e^{2(n-1-\gamma)T_{0}}}{e^{2(n-1-\gamma)T}} e^{2\gamma T_{0}} \mathscr{E}(T_{0}) \notag \\
& + e^{-2(n-1-\gamma)T} \left(e^{[2(n-1)-1]T} - e^{[2(n-1)-1]T_{0}}\right) \mathcal{O}^{2} \notag \\
& + e^{-2(n-1-\gamma)T} \left(e^{[2(n-1)-(1+4\gamma)]T} - e^{[2(n-1)-(1+4\gamma)]T_{0}}\right) \mathds{Y}^{4} \notag \\
& + e^{-2(n-1-\gamma)T} \left(e^{[2(n-1)-(1+3\gamma)]T} - e^{[2(n-1)-(1+3\gamma)]T_{0}}\right) \mathds{Y}^{3} \notag \\
& + e^{-2(n-1-\gamma)T} \left(e^{[2(n-1)-4\gamma]T} - e^{[2(n-1)-4\gamma)]T_{0}}\right) \mathds{Y}^{4} \notag \\
=\; & \frac{e^{2(n-1-\gamma)T_{0}}}{e^{2(n-1-\gamma)T}} e^{2\gamma T_{0}} \mathscr{E}(T_{0}) \notag \\
& + e^{-(1-2\gamma)T} \left(1 - \frac{e^{[2(n-1)-1]T_{0}}}{e^{[2(n-1)-1]T}}\right) \mathcal{O}^{2} \notag \\
& + e^{-(1+2\gamma)T} \left(1 - \frac{e^{[2(n-1)-(1+4\gamma)]T_{0}}}{e^{[2(n-1)-(1+4\gamma)]T}}\right) \mathds{Y}^{4} \notag \\
& + e^{-(1+\gamma)T} \left(1 - \frac{e^{[2(n-1)-(1+3\gamma)]T_{0}}}{e^{[2(n-1)-(1+3\gamma)]T}}\right) \mathds{Y}^{3} \notag \\
& + e^{-2\gamma T} \left(1 - \frac{e^{[2(n-1)-4\gamma]T_{0}}}{e^{[2(n-1)-4\gamma]T}}\right) \mathds{Y}^{3} \notag \\
\lesssim\; & e^{2\gamma T_{0}} \mathscr{E}(T_{0}) 
+ e^{-(1-2\gamma)T} \mathcal{O}^{2} 
+ e^{-(1+2\gamma)T_{0}} \mathds{Y}^{4} 
+ e^{-(1+\gamma)T_{0}} \mathds{Y}^{3} 
+ e^{-2\gamma T_{0}} \mathds{Y}^{3} \notag \\
\lesssim\; & e^{-(1-2\gamma)T} \mathcal{O}^{2} + (\mathcal{I}^{0})^{2} + 1.
\end{align}
We begin from this suboptimal decay estimate 
\begin{align}
\label{eq:suboptimal_recall}
e^{2\gamma T} \mathscr{E}^{low}(T) := \frac{e^{2\gamma T}}{2} \sum_{I\leq 2} \int_{M(T)} \langle \nabla^{I}\Sigma, \nabla^{I}\Sigma \rangle \, \mu_g
\lesssim e^{-(1-2\gamma)T} \mathcal{O}^2 + (\mathcal{I}^0)^2 + 1,  \forall T \in [T_0, T_\infty].
\end{align}

\noindent Observe that the bootstrap assumption \eqref{eq:strap2} allows one to control $e^{2\gamma T}\mathscr{E}^{low}(T)$ uniformly by a constant multiple of $(\mathcal{I}^0)^2 + 1$ provided that $\gamma \leq \frac{1}{2}$. We refrain from immediately fixing $\gamma$ at this threshold. Instead, we pursue an iterative improvement scheme to extend the admissible range of $\gamma$ up to the optimal value $\gamma = 1$.

\noindent By employing the decay rate for the coefficient $\frac{\varphi}{\tau} \lesssim e^{-T}$, and applying the estimate \eqref{model2}, setting $\gamma = \frac{1}{2}$ in \eqref{eq:suboptimal_recall} yields
\begin{align}
e^{T} \mathscr{E}^{low}(T) \lesssim 1 + (\mathcal{I}^0)^2 + \mathcal{O}^2.
\end{align}

\noindent Using this improved decay for the energy, we revisit the energy inequality. Incorporating the refined decay estimate for $||\nabla^{I} \Sigma||_{L^2(M)}$,
\[
||\nabla^{I} \Sigma||_{L^2(M)} \lesssim e^{-\frac{1}{2} T} \left( \mathcal{O} + (\mathcal{I}^0)^2 + 1 \right), \quad I \leq 2,
\]
corresponding to $\gamma = \frac{1}{2}$, the differential inequality governing $\mathscr{E}^{low}(T)$ takes the form
\begin{align}
\frac{d}{dT} \mathscr{E}^{low}\nonumber + 2(n-1) \mathscr{E}^{low} \lesssim e^{-(1+\frac{1}{2}) T} \mathcal{O}^2 + e^{-(1+4\gamma) T} \mathds{Y}^4 + e^{-(1+3\gamma) T} \mathds{Y}^3 + e^{-4\gamma T} \mathds{Y}^4 + e^{-(1+3\gamma) T} \mathcal{I}^0 \mathds{Y}^3.
\end{align}

\noindent Integrating this inequality, one obtains the improved bound
\begin{align}
e^{(1+\frac{1}{2}) T} \mathscr{E}^{low}(T) \lesssim 1 + (\mathcal{I}^0)^2 + \mathcal{O}^2.
\end{align}

\noindent Iterating this procedure—each time replacing $\gamma$ by the newly achieved decay rate—leads to a convergent sequence for $2 \gamma^{optimal}$ given by
\begin{align}
2 \gamma^{optimal} = 1 + \frac{1}{2} \left( 1 + \frac{1}{2} \left( 1 + \cdots \right) \right) = 2,
\quad \text{hence} \quad \gamma^{optimal} = 1.
\end{align}

\noindent Consequently, the optimal decay rate for the low-order energy is established:
\begin{align}
e^{2\gamma T} \mathscr{E}^{low}(T) \lesssim 1 + (\mathcal{I}^0)^2 + \mathcal{O}^2, \quad \gamma = 1.
\end{align}
\end{proof}
\noindent As an immediate corollary, we obtain 
\begin{corollary}
 $\mathcal{F}\lesssim 1+\mathcal{I}^{0}+\mathcal{O}$ and $\mathcal{N}^{\infty}\lesssim 1+(\mathcal{I}^{0})^{2}+\mathcal{O}^{2}$ by elliptic estimates \ref{elliptic}.   
\end{corollary}







\subsection{Energy Estimates}
\label{energy}
\noindent With the individual estimates available, we are ready to complete the energy estimates for the pair $(\Sigma,\mathfrak{T})$. Let us recall the definition of the top order energy $\mathscr{E}^{top}$
\begin{definition}
\label{toporder}
\begin{align}
\mathscr{E}^{top}:=\frac{1}{2}\sum_{I\leq 3}\int_{M}\langle\nabla^{I}\Sigma,\nabla^{I}\Sigma\rangle \mu_{g}+\frac{1}{2}\sum_{I\leq 3}\int_{M}\langle\nabla^{I-1}\mathfrak{T},\nabla^{I-1}\mathfrak{T}\rangle\mu_{g}.    
\end{align}    
\end{definition}

\noindent First, we recall the following elementary proposition 
\begin{proposition}[Integrated energy estimate for the wave system]
\label{wave}
Let $(M,g)$ be a closed Riemannian manifold and consider a time interval $[T_0,T_{\infty}] \subset \mathbb{R}$. Suppose $(\Phi,\Psi)$ is a pair of smooth tensor fields of type $(K,L)$ on $[T_0,T_{\infty}]\times M$ satisfying the coupled system
\begin{align}
\partial_T \Phi &= -N \Psi + P_{\Phi}, \\
\partial_T \Psi &= -N \Delta_g \Phi + Q_{\Psi},
\end{align}
where $N = N(T,x)$ is the lapse function, and $P_{\Phi}$, $Q_{\Psi}$ are given source terms of the same tensorial type as $\Phi$ and $\Psi$, respectively. Denote by $M(T)$ the time slice at time $T$ equipped with the induced metric $g(T)$ and volume form $\mu_g$.

\noindent Then, the following integrated energy identity holds for all $T \in [T_0,T_{\infty}]$:
\begin{align}
\int_{M(T)} \big( |\nabla \Phi|^2 + |\Psi|^2 \big) \,\mu_g
= &\int_{M(T_0)} \big( |\nabla \Phi|^2 + |\Psi|^2 \big) \,\mu_g \\
&+ 2 \int_{T_0}^T \int_{M(t)} \big( \langle \nabla P_{\Phi}, \nabla \Phi \rangle + \langle Q_{\Psi}, \Psi \rangle \big) \,\mu_g\, dt \nonumber \\
&+ \int_{T_0}^T \mathscr{E}\mathscr{R}(t)\, dt, \nonumber
\end{align}
where the error functional $\mathscr{E}\mathscr{R}(t)$ admits the bound
\begin{align}
|\mathscr{E}\mathscr{R}(t)| \lesssim &\left\| \frac{N}{n} - 1 \right\|_{L^\infty(M(t))} \big( \|\nabla \Phi\|_{L^2(M(t))}^2 + \|\Psi\|_{L^2(M(t))}^2 \big) \\
&+ \frac{\varphi}{\tau} \|\Sigma\|_{L^\infty(M(t))} \big( \|\nabla \Phi\|_{L^2(M(t))}^2 + \|\Psi\|_{L^2(M(t))}^2 \big) \nonumber \\
&+ \left[ \frac{\varphi}{\tau} \big( \|\nabla \Sigma\|_{L^4(M(t))} + \|\Sigma\|_{L^\infty(M(t))} \|\nabla (\frac{N}{n} - 1)\|_{L^4(M(t))} \big) + \|\nabla(\frac{N}{n} - 1)\|_{L^4(M(t))} \right] \nonumber \\
&\quad \times \|\Phi\|_{L^4(M(t))} \|\nabla \Phi\|_{L^2(M(t))}. \nonumber
\end{align}
Here, all norms are computed with respect to the metric $g(T)$ on the slice $M(T)$, and the constants implied by $\lesssim$ are of numerical type.

\end{proposition}
\begin{proof}
Let $(\Phi, \Psi)$ be a smooth pair of tensor fields on the compact Riemannian manifold $(M,g)$, where the metric $g = g(T)$ depends smoothly on time $T \in [T_0, T_\infty]$. We consider the energy functional
\[
\mathscr{E}(T) := \frac{1}{2} \int_{M(T)} \big( |\nabla \Phi|^2 + |\Psi|^2 \big) \, \mu_g,
\]
where $\mu_g$ denotes the Riemannian volume form induced by $g(T)$ on the hypersurface $M(T)$, and $|\cdot|$ and $\nabla$ denote the metric norm and Levi-Civita connection respectively.
 
\noindent By the transport theorem for evolving hypersurfaces (cf.~\cite{marsden}), we have the exact identity
\begin{align} \label{eq:transport_energy}
\frac{d}{dT} \int_{M(T)} \left( |\nabla \Phi|^2 + |\Psi|^2 \right) \mu_g
= \int_{M(T)} \partial_T \left( |\nabla \Phi|^2 + |\Psi|^2 \right) \mu_g
+ \int_{M(T)} \left( |\nabla \Phi|^2 + |\Psi|^2 \right) \partial_T \mu_g.
\end{align}
We rewrite \eqref{eq:transport_energy} as
\begin{align} \label{eq:energy_derivative}
\frac{d}{dT} \mathscr{E}(T) 
= \int_{M(T)} \left( \langle \nabla \partial_T \Phi, \nabla \Phi \rangle + \langle \partial_T \Psi, \Psi \rangle \right) \mu_g 
+ \mathscr{ER}(T),
\end{align}
where the \emph{error term} $\mathscr{ER}(T)$ collects the contributions
\[
 \mathscr{ER}(T)\sim \int_{M(T)}\langle[\partial_{T},\nabla]\Phi,\nabla\Phi\rangle\mu_{g}  +\int_{M(T)}\partial_{T}(g^{-1}....g^{-1}g.....g)\nabla\Phi\nabla\Phi\mu_{g}\]
 \[+\int_{M(T)}\partial_{T}(g^{-1}....g^{-1}g.....g)\Psi\Psi\mu_{g}+\int_{M(T)}(|\nabla\Phi|^{2}+|\Psi|^{2})\partial_{T}\mu_{g}  
\]
 
\noindent Recall that the time evolution of the metric $g = g(T)$
\begin{align} \label{eq:metric_evolution}
\partial_T g_{ij} = -\frac{2\varphi}{\tau} N \Sigma_{ij} - 2 \left(1 - \frac{N}{n}\right) g_{ij},
\quad\text{and}\quad
\partial_T g^{ij} = \frac{2\varphi}{\tau} N \Sigma^{ij} + 2 \left(1 - \frac{N}{n}\right) g^{ij},
\end{align}
where $\varphi, \tau, N, n$, and $\Sigma$ are as in the geometric setup.

\noindent Consequently, the time derivative of the volume form $\mu_g$ is given by
\[
\partial_T \mu_g = \frac{1}{2} \mu_g \, \mathrm{tr}_g (\partial_T g) = - n \left(1 - \frac{N}{n}\right) \mu_g,
\]
which follows from \eqref{eq:metric_evolution}.

\noindent The commutator $[\partial_T, \nabla] \Phi$ is expressible via the time derivative of the Christoffel symbols $\Gamma[g]$:
\[
[\partial_T, \nabla] \Phi = (\partial_T \Gamma) \ast \Phi,
\]
where $\ast$ denotes appropriate tensor contractions. By direct computation,
\[
|\partial_T \Gamma| \lesssim \frac{\varphi}{\tau} |\nabla (N \Sigma)| + |\nabla \left(\tfrac{N}{n} - 1\right)|,
\]
hence
\[
|[\partial_T, \nabla] \Phi| \lesssim \frac{\varphi}{\tau} \left( |\nabla \Sigma| + |\Sigma| |\nabla(\tfrac{N}{n} - 1)| \right) + |\nabla(\tfrac{N}{n} - 1)|.
\]

\noindent Applying Hölder's inequality with Sobolev embeddings yields
\begin{align*}
\left| \int_{M(T)} \langle [\partial_T, \nabla] \Phi, \nabla \Phi \rangle \mu_g \right|
&\lesssim \frac{\varphi}{\tau} \left( \|\nabla \Sigma\|_{L^4} + \|\Sigma\|_{L^\infty} \|\nabla(\tfrac{N}{n} - 1)\|_{L^4} \right) \|\Phi\|_{L^4} \|\nabla \Phi\|_{L^2} \\
&\quad + \|\nabla(\tfrac{N}{n} - 1)\|_{L^4} \|\Phi\|_{L^4} \|\nabla \Phi\|_{L^2}.
\end{align*}

\noindent Similarly, the terms involving $\partial_T \mu_g$ are controlled by
\[
\left| \int_{M(T)} \left( |\nabla \Phi|^2 + |\Psi|^2 \right) \partial_T \mu_g \right| \lesssim \left\| 1 - \frac{N}{n} \right\|_{L^\infty} \left( \|\nabla \Phi\|_{L^2}^2 + \|\Psi\|_{L^2}^2 \right).
\]

\noindent Collecting these bounds, we deduce the estimate
\begin{align} \label{eq:error_estimate}
|\mathscr{ER}(T)| \lesssim & \left( \left\| \frac{N}{n} - 1 \right\|_{L^\infty(M(T))} + \frac{\varphi}{\tau} \|\Sigma\|_{L^\infty(M(T))} \right) \left( \|\nabla \Phi\|_{L^2(M(T))}^2 + \|\Psi\|_{L^2(M(T))}^2 \right) \\
& + \left[ \frac{\varphi}{\tau} \left( \|\nabla \Sigma\|_{L^4(M(T))} + \|\Sigma\|_{L^\infty(M(T))} \|\nabla(\tfrac{N}{n} - 1)\|_{L^4(M(T))} \right) + \|\nabla(\tfrac{N}{n} - 1)\|_{L^4(M(T))} \right] \|\Phi\|_{L^4(M(T))} \|\nabla \Phi\|_{L^2(M(T))}. \nonumber
\end{align}

 \noindent We now turn to the principal terms in \eqref{eq:energy_derivative}:
\[
\int_{M(T)} \langle \nabla \partial_T \Phi, \nabla \Phi \rangle \mu_g + \int_{M(T)} \langle \partial_T \Psi, \Psi \rangle \mu_g.
\]
\noindent Substituting the evolution equations
\[
\partial_T \Phi = - N \Psi + P_{\Phi}, \quad \partial_T \Psi = - N \Delta_g \Phi + Q_{\Psi},
\]
we write
\begin{align*}
&\int_{M(T)} \langle \nabla \partial_T \Phi, \nabla \Phi \rangle \mu_g + \int_{M(T)} \langle \partial_T \Psi, \Psi \rangle \mu_g \\
&= \int_{M(T)} \langle \nabla (- N \Psi + P_{\Phi}), \nabla \Phi \rangle \mu_g + \int_{M(T)} \langle - N \Delta_g \Phi + Q_{\Psi}, \Psi \rangle \mu_g \\
&= \int_{M(T)} \langle \nabla P_{\Phi}, \nabla \Phi \rangle \mu_g + \int_{M(T)} \langle Q_{\Psi}, \Psi \rangle \mu_g + \underbrace{\int_{M(T)} \langle \nabla (- N \Psi), \nabla \Phi \rangle \mu_g + \int_{M(T)} \langle - N \Delta_g \Phi, \Psi \rangle \mu_g}_{=: I}.
\end{align*}

\noindent By integration by parts and the absence of boundary (since $M$ is closed), the integral
\[
I = \int_{M(T)} \langle \nabla (- N \Psi), \nabla \Phi \rangle \mu_g + \int_{M(T)} \langle - N \Delta_g \Phi, \Psi \rangle \mu_g
\]
vanishes exactly. This cancellation is a fundamental feature reflecting the underlying hyperbolic structure of the system and guarantees no loss of derivatives in the energy estimates, preserving regularity. This completes the proof.
\end{proof}
\begin{remark}
The terms such as $||\Phi||_{L^{4}(M)}$ can further be controlled by $\sum_{I\leq 1}||\nabla^{I}\Phi||_{L^{2}(M)}$ by means of proposition 
\ref{sobolev}.
\end{remark}


\begin{proposition}[Top-order energy bound]\label{topestimate}
Assume the hypotheses of the main theorem~\ref{main} are satisfied, and let $(\Sigma,\mathfrak{T})$ denote the coupled wave pair associated with the rescaled Einstein system in the gauge fixed framework of this work.  
Let $\mathscr{E}^{\mathrm{top}}(T)\ge0$ denote the top--order energy functional corresponding to $(\Sigma,\mathfrak{T})$ as defined in \ref{toporder}, and let $\mathcal{I}^{0}$ denote the initial size of the data at time $T_0=a\gg 1$.  

\noindent Then for every $T$ in the lifespan of the solution one has
\[
\mathscr{E}^{\mathrm{top}}\ \lesssim\ 1+\mathcal{I}^{0},
\]
where the involved constant is of purely numerical type.
\end{proposition}

\begin{proof}
The proof of the energy estimates for the top-order term is based on the proposition \ref{wave}. First, recall the definition of the top-order energy
\begin{align}
    \mathscr{E}^{top}:=\frac{1}{2}\sum_{I\leq 2}\int_{M}\langle\nabla^{I+1}\Sigma,\nabla^{I+1}\Sigma\rangle \mu_{g}+\frac{1}{2}\sum_{I\leq 2}\int_{M}\langle\nabla^{I}\mathfrak{T},\nabla^{I}\mathfrak{T}\rangle\mu_{g}.
\end{align}
We want to use the proposition \ref{wave} with $\Phi=\nabla^{I}\Sigma$ and $\Psi=\nabla^{I}\mathfrak{T}$ i.e., 
\begin{align}
\label{eq:wave1}
 \partial_{T}\nabla^{I}\Sigma=-N\frac{\varphi}{\tau}\nabla^{I}\mathfrak{T}+P^{I},\\
 \label{eq:wave2}
 \partial_{T}\nabla^{I}\mathfrak{T}=-N\frac{\varphi}{\tau}\Delta_{g}\nabla^{I}\Sigma+Q^{I},
\end{align}
where the error terms $P^{I}$ are $Q^{I}$ are expressed schematically as follows 
\begin{align}
P^{I} &:=& [\partial_{T},\nabla^{I}]\Sigma 
+ \frac{\varphi}{\tau}[\nabla^{I},N]\mathfrak{T} 
+ \nabla^{I}\Bigg( 
  -(n-1)\Sigma_{ij} 
  + \frac{\varphi}{\tau} \nabla_{i}\nabla_{j}\left(\frac{N}{n}-1\right) \nonumber \\
&&\quad 
  + \frac{2\varphi}{\tau}N \Sigma_{ik}\Sigma^{k}_{j}
  - \frac{\varphi}{n\tau}\left(\frac{N}{n}-1\right)g_{ij} 
  - (n-2)\left(\frac{N}{n}-1\right)\Sigma_{ij} 
\Bigg), \\
Q^{I} &:=& [\partial_{T},\nabla^{I}]\mathfrak{T} 
- \frac{\varphi}{\tau}[\nabla^{I},N\Delta_{g}]\Sigma 
+ \nabla^{I}\Bigg(
    n\frac{\varphi}{\tau} \nabla^{l}\nabla_{i}\left(\frac{N}{n}-1\right)\Sigma_{jl} \nonumber \\
&&\quad 
  + n\frac{\varphi}{\tau} \nabla^{l}\left(\frac{N}{n}-1\right)\nabla_{i}\Sigma_{jl}
  + (2 - n)\nabla_{i}\nabla_{j}\left(\frac{N}{n}-1\right) \nonumber \\
&&\quad 
  + n\frac{\varphi}{\tau} \nabla^{l}\nabla_{j}\left(\frac{N}{n}-1\right)\Sigma_{il}
  + n\frac{\varphi}{\tau} \nabla^{l}\left(\frac{N}{n}-1\right)\nabla_{j}\Sigma_{il} \nonumber \\
&&\quad 
  - n\frac{\varphi}{\tau} \Delta_{g}\left(\frac{N}{n}-1\right)\Sigma_{ij}
  - 2n\frac{\varphi}{\tau} \nabla^{l}\left(\frac{N}{n}-1\right)\nabla_{l}\Sigma_{ij} \nonumber \\
&&\quad 
  - \Delta_{g}\left(\frac{N}{n}-1\right)g_{ij}
  + N\frac{\varphi}{\tau}\left( \mathfrak{T}_{ki}\Sigma^{k}_{j} + \mathfrak{T}_{kj}\Sigma^{k}_{i} \right)
  + \frac{2(n-1)}{n^{2}}\left(\frac{N}{n}-1\right)g_{ij}
\Bigg).
\end{align}
A couple of important points to note here. In the expression or $P^{I}$, term $\frac{\varphi}{n\tau}\nabla^{I}(\frac{N}{n}-1)g_{ij}$ is pure trace while $\Sigma$ is transe-verse traceless. Therefore, this term does not contribute to the energy estimates. Similarly, in the expression of $Q^{I}$, the terms $\nabla^{I+2}(\frac{N}{n}-1)g_{ij}$ and $\nabla^{I}(\frac{N}{n}-1)g_{ij}$ are of pure trace type and therefore in the energy estimate for $\mathfrak{T}$ they contribute to the nonlinear term $|\Sigma|^{2}$ since $\tr_{g}\mathfrak{T}=|\Sigma|^{2}$ by the Hamiltonian constraint (\ref{eq:HC11}). 
Now apply proposition \ref{wave} to the system (\ref{eq:wave1})-(\ref{eq:wave2}) to obtain 
\begin{align}
\mathscr{E}^{top}(T)= \mathscr{E}^{top}(T_{0})\nonumber+2\sum_{I\leq 2}\int_{T_{0}}^{T}\int_{M(T)}\left(\langle\nabla P^{I},\nabla^{I+1}\Sigma\rangle+\langle Q^{I},\nabla^{I}\mathfrak{T}\rangle\right)\mu_{g}dt+\int_{T_{0}}^{T}\mathscr{ER}(t)dt,  
\end{align}
where  
\begin{align}
 |\mathscr{ER}(t)|\lesssim \left(||\frac{N}{n}-1||_{L^{\infty}(M(T))}\nonumber+\frac{\varphi}{\tau}||\Sigma||_{L^{\infty}(M(T))}\right)\sum_{I\leq 2}(||\nabla^{I+1}\Sigma||^{2}_{L^{2}(M(T))}+|\nabla^{I}\mathfrak{T}|^{2}_{L^{2}(M(T))})\\\
\nonumber+[\frac{\varphi}{\tau}(||\nabla\Sigma||_{L^{4}(M(T))}+||\Sigma||_{L^{\infty}(M(T))}||\nabla(\frac{N}{n}-1)||_{L^{4}(M(T))})+||\nabla(\frac{N}{n}-1)||_{L^{4}(M(T))}]\\\nonumber\sum_{I\leq 2}||\nabla^{I}\Sigma||_{L^{4}(M(T))}||\nabla^{I+1}\Sigma||_{L^{2}(M(T))}.   
\end{align}
Our goal is to control the error terms and the spacetime integral terms involving $P^{I}$ and $Q^{I}$. We do so schematically as follows 
\begin{align}
 &&\int_{T_{0}}^{T}\int_{M(T)}\langle\nabla P^{I},\nabla^{I+1}\Sigma\rangle\mu_{g}dt\\\nonumber 
 &&= 
 \underbrace{-2(n-1)\int_{T_{0}}^{T}\int_{M(T)}|\nabla^{I+1}\Sigma|^{2}\mu_{g}dt}_{I-decay~term}\\\nonumber
 &&+\sum_{I\leq 2}\sum_{m=0}^{I}\sum_{J_{1}+J_{2}=m}\int_{T_{0}}^{T}\frac{\varphi}{\tau}\int_{M(T)}\nonumber \nabla^{J_{1}+1}(\frac{N}{n}-1)\nabla^{J_{2}+I-m-1}\Sigma\nabla^{I+1}\Sigma\mu_{g}dt\\\nonumber&&+\sum_{I\leq 2}\sum_{m=0}^{I}\sum_{J_{1}+J_{2}+J_{3}=m}\int_{T_{0}}^{T}\frac{\varphi}{\tau}\int_{M(T)}\nabla^{J_{1}}N\nabla^{J_{2}+1}\Sigma\nabla^{J_{3}+I-m-1}\Sigma\nabla^{I+1}\Sigma\mu_{g}dt\\\nonumber 
 &&+\sum_{I\leq 2}\sum_{m=0}^{I}\sum_{J_{1}+J_{2}=m}\int_{T_{0}}^{T}\int_{M(T)}\nabla^{J_{1}+1}(\frac{N}{n}-1)\nabla^{J_{2}+I-m-1}\Sigma\nabla^{I+1}\Sigma\mu_{g}dt\\\nonumber 
 &&+\sum_{I\leq 2}\int_{T_{0}}^{T}\frac{\varphi}{\tau}\int_{M(T)}\nabla^{I+3}(\frac{N}{n}-1)\nabla^{I+1}\Sigma\mu_{g}\\\nonumber 
&&+\sum_{I\leq 2}\sum_{J_{1}+J_{2}+J_{3}=I+1}\int_{T_{0}}^{T}\frac{\varphi}{\tau}\int_{M(T)}\nabla^{J_{1}}N\nabla^{J_{2}}\Sigma\nabla^{J_{3}}\Sigma\nabla^{I+1}\Sigma\mu_{g}\\\nonumber &&+\sum_{I\leq 2}\sum_{J_{1}+J_{2}=I+1}\int_{T_{0}}^{T}\int_{M(T)}\nabla^{J_{1}}(\frac{N}{n}-1)\nabla^{J_{2}}\Sigma\nabla^{I+1}\Sigma\mu_{g}\\\nonumber 
&&+\sum_{I\leq 2}\sum_{J_{1}+J_{2}=I}\int_{T_{0}}^{T}\frac{\varphi}{\tau}\int_{M(T)}\nabla^{J_{1}+1}(\frac{N}{n}-1)\nabla^{J_{2}}\mathfrak{T}\nabla^{I+1}\Sigma\mu_{g}.
\end{align}
We can estimate each term as follows. We use the elliptic estimate whenever necessary. First note that $\frac{\varphi}{\tau}=\frac{e^{-T}}{\sqrt{e^{-2T}+3n(n+1)}}\lesssim e^{-T}$ 
\begin{align}
| \sum_{I\leq 2}\sum_{m=0}^{I}\sum_{J_{1}+J_{2}=m}\int_{T_{0}}^{T}\frac{\varphi}{\tau}\int_{M(T)}\nonumber \nabla^{J_{1}+1}(\frac{N}{n}-1)\nabla^{J_{2}+I-m}\Sigma\nabla^{I+1}\Sigma |\lesssim \int_{T_{0}}^{T}e^{-(1+3\gamma)t}\mathds{Y}^{4}dt\\\nonumber \lesssim (e^{-(1+3\gamma)T_{0}}-e^{-(1+3\gamma)T})\mathds{Y}^{3}\Gamma\lesssim  e^{-(1+3\gamma)T_{0}}(1-e^{-(1+3\gamma)(T-T_{0})})\mathds{Y}^{3}\Gamma\lesssim 1. 
\end{align}
The maximum derivative on the lapse function in the previous estimate is $I+1$ and therefore by the elliptic estimate, it is bounded by $||\Sigma||^{2}_{H^{I-1}(M)}$ which yields $e^{-2\gamma t}$ decay. The next term is estimated as 
\begin{align}
 |\sum_{I\leq 2}\sum_{m=0}^{I}\sum_{J_{1}+J_{2}+J_{3}=m}\int_{T_{0}}^{T}\frac{\varphi}{\tau}\int_{M(T)}\nabla^{J_{1}}N\nabla^{J_{2}+1}\Sigma\nabla^{J_{3}+I-m}\Sigma\nonumber\nabla^{I+1}\Sigma\mu_{g}dt|\lesssim \int_{T_{0}}^{T}e^{-(1+\gamma)t}\mathds{Y}\Gamma^{2}dt\\\nonumber 
 \lesssim e^{-(1+\gamma)T_{0}}(1-e^{-(1+\gamma)(T-T_{0})})\mathds{Y}\Gamma^{2}\lesssim 1.
\end{align}
The next term is estimated as 
\begin{align}
|\sum_{I\leq 2}\sum_{m=0}^{I}\sum_{J_{1}+J_{2}=m}\int_{T_{0}}^{T}\int_{M(T)}\nabla^{J_{1}+1}(\frac{N}{n}-1)\nabla^{J_{2}+I-m}\Sigma\nabla^{I+1}\Sigma\mu_{g}dt|\lesssim \int_{T_{0}}^{T}e^{-3\gamma t} \Gamma\mathds{Y}^{3}dt\\\nonumber \lesssim e^{-3\gamma T_{0}}(1-e^{-3\gamma(T-T_{0})})\Gamma\mathds{Y}^{3}\lesssim 1, 
\end{align}
\begin{align}
|\sum_{I\leq 2}\int_{T_{0}}^{T}\frac{\varphi}{\tau}\int_{M(T)}\nabla^{I+3}(\frac{N}{n}-1)\nabla^{I+1}\Sigma\mu_{g}|\lesssim \int_{T_{0}}^{T}e^{-t}\Gamma^{3}dt\lesssim e^{-T_{0}}(1-e^{-(T-T_{0})})\Gamma^{3}\lesssim 1,  
\end{align}
\begin{align}
|\sum_{I\leq 2}\sum_{J_{1}+J_{2}+J_{3}=I+1}\int_{T_{0}}^{T}\frac{\varphi}{\tau}\int_{M(T)}\nabla^{J_{1}}N\nabla^{J_{2}}\Sigma\nabla^{J_{3}}\Sigma\nabla^{I+1}\Sigma\mu_{g}|\lesssim   \int_{T_{0}}^{T}e^{-(1+\gamma)t}\Gamma^{2}\mathds{Y}dt\\\nonumber 
\lesssim e^{-(1+\gamma)T_{0}}(1-e^{-(1+\gamma)(T-T_{0})})\Gamma^{2}\mathds{Y}\lesssim 1.
\end{align}
The next term is estimated as follows 
\begin{align}
|\sum_{I\leq 2}\sum_{J_{1}+J_{2}=I+1}\int_{T_{0}}^{T}\int_{M(T)}\nabla^{J_{1}}(\frac{N}{n}-1)\nabla^{J_{2}}\Sigma\nabla^{I+1}\Sigma\mu_{g}|\lesssim \int_{T_{0}}^{T}e^{-2\gamma t}\Gamma^{2}\mathds{Y}^{2}dt\\\nonumber \lesssim e^{-2\gamma T_{0}}(1-e^{-2\gamma(T-T_{0})})\Gamma^{2}\mathds{Y}^{2}\lesssim 1,    
\end{align}
where we noted that the top order term in lapse $||\nabla^{4}(\frac{N}{n}-1)||_{L^{2}(M)}$ is estimated by $||\Sigma||^{2}_{H^{2}(M)}$ which exhibits $e^{-2\gamma t}$ decay. The last term in the expression of $P^{I}$ is estimated as 
\begin{align}
 |\sum_{I\leq 2}\sum_{J_{1}+J_{2}=I}\int_{T_{0}}^{T}\frac{\varphi}{\tau}\int_{M(T)}\nabla^{J_{1}+1}(\frac{N}{n}-1)\nabla^{J_{2}}\mathfrak{T}\nabla^{I+1}\Sigma\mu_{g}|\lesssim \int_{T_{0}}^{T}e^{-(1+2\gamma) t}\Gamma^{2}\mathds{Y}^{2} dt\\\nonumber 
 \lesssim e^{-(1+2\gamma) T_{0}}(1-e^{-(1+2\gamma)(T-T_{0})})\Gamma^{2}\mathds{Y}^{2}\lesssim 1.
\end{align}
Now we estimate the terms in the expressions of $Q^{I}$. First, the spacetime integral involving $Q^{I}$ is explicitly evaluated as follows (written in schematic notation) 
\begin{align}
&&\int_{T_{0}}^{T}\int_{M(T)}\langle Q^{I},\nabla^{I}\mathfrak{T}\rangle\mu_{g}dt  \\\nonumber 
&&\sim \sum_{I\leq 2}\sum_{m=0}^{I-1}\sum_{J_{1}+J_{2}=m}\int_{T_{0}}^{T}\frac{\varphi}{\tau}\int_{M}\nabla^{J_{1}+1}(\frac{N}{n}-1)\nabla^{J_{2}+I-m-1}\mathfrak{T}\nabla^{I}\mathfrak{T}\mu_{g}dt\\\nonumber 
&&+\sum_{I\leq 2}\sum_{m=0}^{I-1}\sum_{J_{1}+J_{2}+J_{3}=m}\int_{T_{0}}^{T}\frac{\varphi}{\tau}\int_{M(T)}\nabla^{J_{1}}N\nabla^{J_{2}+1}\Sigma\nabla^{J_{3}+I-m-1}\mathfrak{T}\nabla^{I}\mathfrak{T}\mu_{g}dt\\\nonumber 
&&+\sum_{I\leq 2}\sum_{m=0}^{I-1}\sum_{J_{1}+J_{2}=m}\int_{T_{0}}^{T}\int_{M(T)}\nabla^{J_{1}+1}(\frac{N}{n}-1)\nabla^{J_{2}+I-m-1}\mathfrak{T}\nabla^{I}\mathfrak{T}\mu_{g}dt\\\nonumber 
&&+\sum_{I\leq 2}\int_{T_{0}}^{T}\frac{\varphi}{\tau}\int_{M}\nabla(\frac{N}{n}-1)\nabla^{I+1}\Sigma\nabla^{I}\mathfrak{T}\mu_{g}dt\\\nonumber
&&+\sum_{I\leq 2}\sum_{m=0}^{I-1}\sum_{J_{1}+J_{2}=m}\int_{T_{0}}^{T}\frac{\varphi}{\tau}\int_{M(T)}N\nabla^{J_{1}}\text{Riem}\nabla^{J_{2}+I-m}\Sigma\nabla^{I}\mathfrak{T}\mu_{g}dt\\\nonumber 
&&+\sum_{I\leq 2}\sum_{m=0}^{I-1}\sum_{J_{1}+J_{2}=m}\int_{T_{0}}^{T}\frac{\varphi}{\tau}\int_{M(T)}\nabla^{J_{1}+1}\text{Riem}\nabla^{J_{2}+I-m-1}\Sigma,\nabla^{I}\mathfrak{T}\rangle\mu_{g}dt\\\nonumber 
&&+\sum_{I\leq 2}\sum_{J_{1}+J_{2}=I}\int_{T_{0}}^{T}\frac{\varphi}{\tau}\int_{M(T)}\nabla^{J_{1}}\text{Riem}\nabla^{J_{2}}\Sigma\nabla^{I}\mathfrak{T}\mu_{g}dt\\\nonumber 
&&+\sum_{I\leq 2}\sum_{J_{1}+J_{2}=I}\int_{T_{0}}^{T}\frac{\varphi}{\tau}\int_{M(T)}\nabla^{J_{1}+2}(\frac{N}{n}-1)\nabla^{J_{2}}\Sigma,\nabla^{I}\mathfrak{T}\rangle\mu_{g}dt\\\nonumber &&+\sum_{I\leq 2}\sum_{J_{1}+J_{2}=I}\int_{T_{0}}^{T}\frac{\varphi}{\tau}\int_{M(T)}\nabla^{J_{1}+1}(\frac{N}{n}-1)\nabla^{J_{2}+1}\Sigma\nabla^{I}\mathfrak{T}\mu_{g}dt\\\nonumber 
&&+\sum_{I\leq 2}\sum_{J_{1}+J_{2}=I}\int_{T_{0}}^{T}\frac{\varphi}{\tau}\int_{M}\langle\nabla^{J_{1}}\mathfrak{T}\nabla^{J_{2}}\Sigma \nabla^{I}\mathfrak{T}\rangle\mu_{g}dt\\\nonumber 
&&+\sum_{I\leq 2}\sum_{J_{1}+J_{2}=I-1}\int_{T_{0}}^{T}\frac{\varphi}{\tau}\int_{M(T)}\nabla^{J_{1}+1}(\frac{N}{n}-1)\nabla^{J_{2}+1}\Sigma\nabla^{I}\mathfrak{T}\mu_{g}dt.
\end{align}
Now we estimate each term separately. The first term is estimated as follows 
\begin{align}
|\sum_{I\leq 2}\sum_{m=0}^{I-1}\sum_{J_{1}+J_{2}=m}\int_{T_{0}}^{T}\frac{\varphi}{\tau}\int_{M}\nabla^{J_{1}+1}(\frac{N}{n}-1)\nabla^{J_{2}+I-m-1}\mathfrak{T}\nabla^{I}\mathfrak{T}\mu_{g}dt|\\\nonumber 
\lesssim \int_{T_{0}}^{T}e^{-(1+2\gamma)t}\mathds{Y}^{2}\Gamma^{2}dt \lesssim e^{-(1+2\gamma) T_{0}}(1-e^{-(1+2\gamma)(T-T_{0})})\Gamma^{2}\mathds{Y}^{2} \lesssim 1. 
\end{align}
\noindent Here, the elliptic estimate \ref{elliptic} for the lapse function is used. The next term is estimated as 
\begin{align}
|\sum_{I\leq 2}\sum_{m=0}^{I-1}\sum_{J_{1}+J_{2}+J_{3}=m}\int_{T_{0}}^{T}\frac{\varphi}{\tau}\int_{M(T)}\nabla^{J_{1}}N\nabla^{J_{2}+1}\Sigma\nabla^{J_{3}+I-m-1}\mathfrak{T}\nabla^{I}\mathfrak{T}\mu_{g}dt|\\\nonumber 
\lesssim \int_{T_{0}}^{T}e^{-(1+\gamma)t}\Gamma^{2}\mathds{Y}dt\lesssim e^{-(1+\gamma) T_{0}}(1-e^{-(1+\gamma)(T-T_{0})})\Gamma^{2}\mathds{Y}\lesssim 1. 
\end{align}
Notice that $\Sigma$ appears as $\nabla^{I}\Sigma$ at the top most order and therefore is estimated in $L^{2}$ yielding extra $e^{-\gamma t}$ decay. The next term is estimated as 
\begin{align}
|\sum_{I\leq 2}\sum_{m=0}^{I-1}\sum_{J_{1}+J_{2}=m}\int_{T_{0}}^{T}\int_{M(T)}\nabla^{J_{1}+1}(\frac{N}{n}-1)\nabla^{J_{2}+I-m-1}\mathfrak{T}\nabla^{I}\mathfrak{T}\mu_{g}dt|\\\nonumber 
\lesssim \int_{T_{0}}^{T}e^{-2\gamma t}\mathds{Y}^{2}\Gamma^{2}dt\lesssim e^{-2\gamma T_{0}}(1-e^{-2\gamma(T-T_{0})})\Gamma^{2}\mathds{Y}^{2}\lesssim 1.
\end{align}
The next terms are estimated as 
\begin{align}
|\sum_{I\leq 2}\int_{T_{0}}^{T}\frac{\varphi}{\tau}\int_{M}\nabla(\frac{N}{n}-1)\nabla^{I+1}\Sigma\nabla^{I}\mathfrak{T}\mu_{g}dt|\lesssim \int_{T_{0}}^{T}e^{-(1+2\gamma)t}\mathds{Y}^{2}\Gamma^{2}dt\\\nonumber \lesssim e^{-(1+2\gamma) T_{0}}(1-e^{-(1+2\gamma)(T-T_{0})})\Gamma^{2}\mathds{Y}^{2} \lesssim 1,     
\end{align}
\begin{align}
|\sum_{I\leq 2}\sum_{m=0}^{I-1}\sum_{J_{1}+J_{2}=m}\int_{T_{0}}^{T}\frac{\varphi}{\tau}\int_{M(T)}N\nabla^{J_{1}}\text{Riem}\nabla^{J_{2}+I-m}\Sigma\nabla^{I}\mathfrak{T}\mu_{g}dt|\\\nonumber 
\lesssim \int_{T_{0}}^{T}e^{-(1+\gamma)t}\Gamma^{2}\mathds{Y}dt\lesssim e^{-(1+\gamma) T_{0}}(1-e^{-(1+\gamma)(T-T_{0})})\Gamma^{2}\mathds{Y} \lesssim 1,    
\end{align}
\begin{align}
 |\sum_{I\leq 2}\sum_{m=0}^{I-1}\sum_{J_{1}+J_{2}=m}\int_{T_{0}}^{T}\frac{\varphi}{\tau}\int_{M(T)}\nabla^{J_{1}+1}\text{Riem}\nabla^{J_{2}+I-m-1}\Sigma,\nabla^{I}\mathfrak{T}\rangle\mu_{g}dt|\\\nonumber 
 \lesssim \int_{T_{0}}^{T}e^{-(1+\gamma)t}\Gamma^{2}\mathds{Y}dt\lesssim e^{-(1+\gamma) T_{0}}(1-e^{-(1+\gamma)(T-T_{0})})\Gamma^{2}\mathds{Y} \lesssim 1,
\end{align}
\begin{align}
 |\sum_{I\leq 2}\sum_{J_{1}+J_{2}=I}\int_{T_{0}}^{T}\frac{\varphi}{\tau}\int_{M(T)}\nabla^{J_{1}}\text{Riem}\nabla^{J_{2}}\Sigma\nabla^{I}\mathfrak{T}\mu_{g}dt|\\\nonumber
 \lesssim \int_{T_{0}}^{T}e^{-(1+\gamma)t}\Gamma^{2}\mathds{Y}dt\lesssim e^{-(1+\gamma) T_{0}}(1-e^{-(1+\gamma)(T-T_{0})})\Gamma^{2}\mathds{Y} \lesssim 1,
\end{align}
\begin{align}
 |\sum_{I\leq 2}\sum_{J_{1}+J_{2}=I}\int_{T_{0}}^{T}\frac{\varphi}{\tau}\int_{M(T)}\nabla^{J_{1}+2}(\frac{N}{n}-1)\nabla^{J_{2}}\Sigma\nabla^{I}\mathfrak{T}\mu_{g}dt|\\\nonumber 
 \lesssim \int_{T_{0}}^{T}e^{-(1+3\gamma)t}\mathds{Y}^{3}\Gamma dt\lesssim e^{-(1+3\gamma) T_{0}}(1-e^{-(1+3\gamma)(T-T_{0})})\Gamma\mathds{Y}^{3} \lesssim 1,
\end{align}
\begin{align}
 |\sum_{I\leq 2}\sum_{J_{1}+J_{2}=I}\int_{T_{0}}^{T}\frac{\varphi}{\tau}\int_{M(T)}\nabla^{J_{1}+1}(\frac{N}{n}-1)\nabla^{J_{2}+1}\Sigma\nabla^{I}\mathfrak{T}\mu_{g}dt|\\\nonumber 
 \lesssim \int_{T_{0}}^{T}e^{-(1+2\gamma)t}\mathds{Y}^{2}\Gamma^{2}dt\lesssim e^{-(1+2\gamma) T_{0}}(1-e^{-(1+2\gamma)(T-T_{0})})\Gamma^{2}\mathds{Y}^{2} \lesssim 1,
\end{align}
\begin{align}
|\sum_{I\leq 2}\sum_{J_{1}+J_{2}=I}\int_{T_{0}}^{T}\frac{\varphi}{\tau}\int_{M}\nabla^{J_{1}}\mathfrak{T}\nabla^{J_{2}}\Sigma \nabla^{I}\mathfrak{T}\mu_{g}dt|\lesssim \int_{T_{0}}^{T}e^{-(1+\gamma)t}\mathds{Y}\Gamma^{2}dt\\\nonumber \lesssim e^{-(1+\gamma) T_{0}}(1-e^{-(1+\gamma)(T-T_{0})})\Gamma^{2}\mathds{Y}\lesssim 1,
\end{align}
and 
\begin{align}
|\sum_{I\leq 2}\sum_{J_{1}+J_{2}=I-1}\int_{T_{0}}^{T}\frac{\varphi}{\tau}\int_{M(T)}\nabla^{J_{1}+1}(\frac{N}{n}-1)\nabla^{J_{2}+1}\Sigma\nabla^{I}\mathfrak{T}\mu_{g}dt|\\\nonumber \lesssim \int_{T_{0}}^{T}e^{-(1+2\gamma)t}\Gamma^{2}\mathds{Y}^{2}dt\lesssim e^{-(1+2\gamma) T_{0}}(1-e^{-(1+2\gamma)(T-T_{0})})\Gamma^{2}\mathds{Y}^{2}\lesssim 1.     
\end{align}
Now note that the first term in the expression involving $P^{I}$ contains term $I$ which is negative definite. Therefore, the collection of every term yields 
\begin{align}
 \mathscr{E}^{top}(T)\lesssim \mathscr{E}^{top}(T_{0})+1\lesssim \mathcal{I}^{0}+1   
\end{align}
which completes the proof of the theorem. 
\end{proof}
\begin{corollary}
 $\mathcal{O}\lesssim 1+\mathcal{I}^{0}$   
\end{corollary}

\begin{remark}
 Notice that the estimates are uniform in $T$, so one can take the limit $T_{\infty}=\infty$.   
\end{remark}

\section{Proof of the main results}
\label{mainproof}
\subsection{Proof of theorem \ref{main}}
\label{proof1}
\noindent In this section, we prove the main theorem with the aid of the uniform estimates obtained in section \ref{energy}. Note that we have already presented the idea of the proof in section \ref{bootstrapargument}. Here we will use a contradiction argument. First, recall the statement of the main theorem 
\begin{theorem-non} [Global existence, $\Lambda>0, \sigma(M)\leq 0$]
Let $(\widehat{M}^{3+1}, \widehat{g})$ be a globally hyperbolic Lorentzian spacetime satisfying the Einstein vacuum equations with positive cosmological constant $\Lambda > 0$, and suppose that $\widehat{M}$ admits a constant mean curvature (CMC) foliation by compact spacelike hypersurfaces diffeomorphic to a closed $3$-manifold $M$. Assume furthermore that $M$ is of negative Yamabe type, i.e., $\sigma(M) \leq 0$.

\noindent Fix a smooth background Riemannian metric $\xi_0$ on $M$ and a constant $C > 1$. For any sufficiently large initial energy quantity $\mathcal{I}^{0} > 0$, there exists a constant $a = a(\mathcal{I}^{0}) > 0$, sufficiently large so that $\mathcal{I}^{0} e^{-a/10} < 1$.

\noindent Let $(g_0, \Sigma_0)$ be an initial data set verifying the Einstein constraint equations at initial CMC time $T_0 = a$, written in CMC-transported spatial coordinates and satisfying:
\begin{align}
\label{eq:1}
C^{-1} \xi_0 \leq g_0 \leq C \xi_0,
\end{align}
\begin{align}
\label{eq:2}
\sum_{I = 0}^{3} \| \nabla^I \Sigma_0 \|_{L^2(M)} + \sum_{I = 0}^{2} \left( \| \nabla^I \mathfrak{T}[g_0] \|_{L^2(M)} + \| e^a \nabla^I \Sigma_0 \|_{L^2(M)} \right) \leq \mathcal{I}^0,
\end{align}
where $\mathfrak{T}[g_0]$ denotes the renormalized trace-free spatial Ricci curvature tensor of $g_0$.

\noindent Then, the Einstein-$\Lambda$ evolution equations admit a unique classical solution
\[
T \mapsto (g(T), \Sigma(T)) \in \mathcal{C}^\infty([T_0, \infty) \times M)
\]
in CMC-transported spatial coordinates, satisfying the constraint equations at each slice $T$ and obeying the following uniform a priori estimates for all $T \in [T_0, \infty)$:
\begin{align}
\sum_{I = 0}^{3} \| \nabla^I \Sigma(T) \|_{L^2(M)} + \sum_{I = 0}^{2} \left( \| \nabla^I \mathfrak{T}[g(T)] \|_{L^2(M)} + \| e^T \nabla^I \Sigma(T) \|_{L^2(M)} \right) \leq C_1(1 + \mathcal{I}^0),
\end{align}
\begin{align}
C_2^{-1} g_0 \leq g(T) \leq C_2 g_0.
\end{align}
Here, $C_1, C_2 > 0$ are numerical constants depending only on the universal geometric and analytic data of the problem (e.g., Sobolev constants of $(M, g_0)$ and the constants in the structure equations), but independent of $T$. The developed spacetime is future geodesically complete.

\noindent Moreover, the solution $(g(T), \Sigma(T))$ converges in the $C^\infty$ topology, as $T \to \infty$, to a limiting Riemannian metric $\widetilde{g}$ on $M$ of pointwise constant negative scalar curvature, in the sense that
\[
\Sigma(T) \to 0 \quad \text{and} \quad g(T) \to \widetilde{g}, \quad \text{as } T \to \infty,
\]
with convergence holding in all Sobolev norms. In particular, the spacetime $(\widehat{M}, \widehat{g})$ admits a future-complete CMC foliation asymptotic to a constant negative scalar curvature slice.  
\end{theorem-non}

\noindent We begin by establishing global existence of the CMC Einstein-$\Lambda$ flow under the stated assumptions. The proof proceeds by a standard continuation argument, combining local well-posedness for the elliptic-hyperbolic system with a contradiction argument based on a priori estimates and weak compactness.

\noindent Let us first observe that the coupled Einstein-$\Lambda$ system, expressed in CMC-transported spatial coordinates, reduces to a manifestly elliptic-hyperbolic formulation for the unknowns $(g, \Sigma)$, consisting of a hyperbolic evolution equation for the metric $g$, a transport-type equation for $\Sigma$, and elliptic constraint equations for the lapse and shift. The system satisfies the structural conditions of the classical theory developed for mixed elliptic-hyperbolic systems (see, e.g., \cite{elliptichyperbolic}). In particular, the associated initial value problem admits a unique classical solution on a time interval $[T_0, T)$ for some $T > T_0$, with the initial data $(g_0, \Sigma_0)$ prescribed at time $T_0 = a$.

\noindent We assume, for contradiction, that the maximal existence interval is bounded above by some finite time $T < \infty$, i.e., the solution ceases to exist beyond time $T$. By the uniform a priori estimates derived in Section~\ref{energy}, combined with Sobolev persistence of regularity, the fields $g(T'), \Sigma(T'), \mathfrak{T}(T')$ remain uniformly bounded in the Sobolev spaces $H^3(M)$, $H^3(M)$, and $H^2(M)$ respectively, for all $T' \in [T_0, T)$.

\noindent Using reflexivity and the Banach-Alaoglu theorem, we extract a weakly convergent subsequence $\{T_n'\}_{n \in \mathbb{N}} \subset [T_0, T)$ with $T_n' \to T^-$ such that:
\begin{align*}
\Sigma(T_n') &\rightharpoonup \Sigma(T) \quad \text{weakly in } H^3(M), \\
\mathfrak{T}(T_n') &\rightharpoonup \mathfrak{T}(T) \quad \text{weakly in } H^2(M).
\end{align*}
Moreover, by weak lower semicontinuity of the Sobolev norms and the uniform estimates, we have:
\begin{align*}
\sum_{I = 0}^{3} \| \nabla^I \Sigma(T) \|_{L^2(M)} &\leq \liminf_{n \to \infty} \sum_{I = 0}^{3} \| \nabla^I \Sigma(T_n') \|_{L^2(M)} \lesssim 1 + \mathcal{I}^0, \\
\sum_{I = 0}^{2} \| \nabla^I \mathfrak{T}(T) \|_{L^2(M)} &\leq \liminf_{n \to \infty} \sum_{I = 0}^{2} \| \nabla^I \mathfrak{T}(T_n') \|_{L^2(M)} \lesssim 1 + \mathcal{I}^0.
\end{align*}
Since the constraint equations are preserved along the flow and depend smoothly on the variables, it follows that $(g(T), \Sigma(T))$ defines a compatible initial data set at time $T$. By the same local well-posedness theory, we may extend the solution to a strictly larger time interval $[T, T + \epsilon)$ for some $\epsilon > 0$ depending only on the bounds for $\Sigma(T)$ and $\mathfrak{T}(T)$. This contradicts the assumption that $[T_0, T)$ was the maximal interval of existence. We therefore conclude that the solution extends globally in time, i.e., $T = \infty$.

\noindent We now address the asymptotic behavior as $T \to \infty$. From the a priori energy estimates derived in Section~\ref{energy}, and in particular from exponential decay of the appropriate norms, the time derivative of the metric satisfies:
\[
\|\partial_T g(T')\|_{H^2(M)} \lesssim \mathcal{I}^0 e^{-(1 + \gamma) T'} + (\mathcal{I}^0)^2 e^{-2\gamma T'}
\]
for all $T' \geq T_0$. Integrating this differential inequality from $T$ to infinity yields
\begin{align*}
\| g(T) - g(\infty) \|_{H^2(M)} &\leq \int_T^{\infty} \| \partial_{T'} g(T') \|_{H^3(M)} \, dT' \\
&\lesssim \mathcal{I}^0 e^{-(1 + \gamma) T} + (\mathcal{I}^0)^2 e^{-2\gamma T},
\end{align*}
which shows that $g(T)$ converges strongly in $H^2(M)$ to a limiting Riemannian metric $\widetilde{g} := g(\infty)$ as $T \to \infty$. Similarly,
\[
\| \Sigma(T) \|_{H^2(M)} \lesssim \mathcal{I}^0 e^{-\gamma T} \to 0,
\]
so that $\Sigma(T) \to 0$ strongly in $H^2(M)$, and hence in $C^0(M)$ by Sobolev embedding.

\noindent Finally, we consider the scalar curvature of the limiting metric. The Hamiltonian constraint equation takes the form
\[
R(g(T)) + \frac{n - 1}{n} = |\Sigma(T)|^2_g,
\]
and from the above convergence and the Sobolev algebra property of $H^s(M)$ for $s > \frac{3}{2}$, it follows that:
\begin{align*}
\| R(g(T)) + \tfrac{n - 1}{n} \|_{H^2(M)} &= \| |\Sigma(T)|^2_g \|_{H^2(M)} \\
&\lesssim \| \Sigma(T) \|_{H^2(M)}^2 \to 0.
\end{align*}
Therefore, in the limit $T \to \infty$, we obtain:
\[
R(\widetilde{g}) = \lim_{T \to \infty} R(g(T)) = - \frac{n - 1}{n},
\]
with convergence strongly in $H^2(M)$ and hence pointwise. Thus, the limiting Riemannian manifold $(M, \widetilde{g})$ is a smooth manifold equipped with a metric of constant negative scalar curvature, and the full spacetime $(\widetilde{M}, \widehat{g})$ admits a future-complete CMC foliation asymptotic to this constant curvature geometry.
Let $\mathcal{C}(\lambda)$ be a future-directed causal geodesic (timelike or null), with affine parameter $\lambda$ and tangent vector $\alpha^\mu = \frac{d \mathcal{C}^\mu}{d\lambda}$ satisfying $\widehat{g}(\alpha, \alpha) = -1$ (timelike) or $0$ (null). In CMC-transported coordinates $(T, x^i)$, we write $\mathcal{C}(\lambda) = (T(\lambda), x^i(\lambda))$, and define $\alpha^0 := \frac{dT}{d\lambda}$. To prove future completeness, it suffices to show
\[
\lambda(T) \to \infty \quad \text{as} \quad T \to \infty \quad \Longleftrightarrow \quad \int_{T_0}^\infty \frac{1}{\alpha^0(T)} \, dT = \infty.
\]

\noindent Let $N$ denote the rescaled lapse in the CMC gauge so that the spacetime metric reads
\[
\widehat{g} = -N^2 dT^2 + g_{ij}(T) dx^i dx^j.
\]
Introduce the future-directed co-vector field $Z := N \partial_T$, and decompose the geodesic tangent as
\[
\alpha = N \alpha^0 Z + W,
\]
where $W$ is tangent to the constant-$T$ hypersurfaces. Then
\[
\widehat{g}(\alpha, \alpha) = -N^2 (\alpha^0)^2 + |W|^2_{g(T)} = \begin{cases}
-1 & \text{(timelike)} \\
0 & \text{(null)}
\end{cases},
\]
so that $|W|^2_{g(T)} = N^2 (\alpha^0)^2 + \epsilon$, where $\epsilon = 1$ or $0$ respectively. Thus,
\[
|W|^{2}_{g(T)}\geq N^2 (\alpha^0)^2, \quad \text{and in all cases} \quad f := N^2 (\alpha^0)^2 \geq 0.
\]

\noindent Differentiating $f$ along $\mathcal{C}$ and using the geodesic equation $\nabla[\widehat{g}]_{\alpha} \alpha = 0$, we obtain
\[
\frac{d}{dT} f = \frac{2}{\alpha^0} \widehat{g}(\alpha, Z) \cdot \widehat{g}(\alpha, \nabla[\widehat{g}]_{\alpha} Z),
\]
and using $\widehat{g}(\alpha, Z) = -N \alpha^0$, we compute
\[
\widehat{g}(\alpha, \nabla_\alpha Z) = -\alpha^0 \nabla_W N + K_{ij} W^i W^j,
\]
where $K$ is the second fundamental form of the constant-$T$ slices. Thus,
\[
\frac{d}{dT} f = -2 N \left( \alpha^0 \nabla_W N - K_{ij} W^i W^j \right).
\]

\noindent Decompose $K = \Sigma + \tfrac{\tau}{n} g$, with $\tau = \mathrm{tr}_g K < 0$ and $\Sigma$ trace-free. Using Cauchy--Schwarz and Sobolev embedding, we estimate
\[
\left| \frac{d}{dT} \log f \right| \leq \| \nabla N \|_{L^\infty} + \| N \Sigma \|_{L^\infty}.
\]
By the main decay estimates (Theorem~\ref{main}), we have
\[
\| \nabla N \|_{L^\infty} + \| N \Sigma \|_{L^\infty} \lesssim e^{-T},
\]
and thus
\[
\left| \frac{d}{dT} \log f \right| \lesssim e^{-T}, \quad \Rightarrow \quad f(T) \to C > 0 \quad \text{as } T \to \infty.
\]

\noindent Therefore, $N^2 (\alpha^0)^2 \leq C$ uniformly for large $T$, and
\[
\int_{T_0}^\infty \frac{1}{\alpha^0} \, dT \geq \int_{T_0}^\infty \frac{N}{\sqrt{f}} \, dT \gtrsim \int_{T_0}^\infty N \, dT.
\]
Since $N = n + \mathcal{O}(e^{-T})$, the lapse remains uniformly bounded below by a positive constant, implying
\[
\int_{T_0}^\infty N \, dT = \infty.
\]
It follows that $\lambda(T) \to \infty$ as $T \to \infty$, establishing future completeness for all causal geodesics.




\hfill$\square$
\begin{remark}[Non-convergence to an Einstein Metric]
\label{noconvergence}
We emphasize that in general, the trace-free renormalized Ricci tensor $\mathfrak{T}$ does not decay to zero as $T \to \infty$. This indicates that the limiting metric $\widetilde{g} := \lim_{T \to \infty} g(T)$ fails to be Einstein, even though it has constant scalar curvature. The underlying obstruction is geometric: the compact manifold $M$ may not admit any Einstein metric (e.g., hyperbolic metric in the case $n=3$), and hence the evolution governed by the Einstein equations with $\Lambda > 0$ does not necessarily drive the geometry toward an Einstein configuration.

\noindent To substantiate this, consider the evolution of the squared $L^2$-norm of $\mathfrak{T}$,
\begin{align}
S(T) := \| \mathfrak{T}(T) \|_{L^2(M)}^2.
\end{align}
Using the evolution equation for $\mathfrak{T}$ and the uniform high-order energy bounds obtained in Section~\ref{proof1}, along with standard Sobolev inequalities and Grönwall-type arguments, one derives the differential inequality
\begin{align}
\frac{d}{dT} S(T) \geq -C (\mathcal{I}^0)^2 e^{-2T},
\end{align}
where $C>0$ is a constant depending only on the background geometry and the constants in the elliptic-hyperbolic structure of the system. Integrating this inequality from $T_0$ to $T$, we obtain
\begin{align}
S(T) \geq S(T_0) - C (\mathcal{I}^0)^2 e^{-2T_0} \left(1 - e^{-2(T - T_0)}\right).
\end{align}
Letting $T \to \infty$, we find
\begin{align}
\liminf_{T \to \infty} \| \mathfrak{T}(T) \|_{L^2(M)}^2 \geq S(T_0) - C (\mathcal{I}^0)^2 e^{-2T_0}.
\end{align}
Hence, for sufficiently large initial energy $\mathcal{I}^0, T_{0}$, and correspondingly large $S(T_0)$ (which is permitted by our assumptions), one may ensure that
\begin{align}
\liminf_{T \to \infty} \| \mathfrak{T}(T) \|_{L^2(M)}^2 \gtrsim 1,
\end{align}
i.e., $\mathfrak{T}$ remains bounded away from zero as $T \to \infty$. This implies that the limiting metric $\widetilde{g}$ fails to satisfy the Einstein condition $\operatorname{Ric}_{\widetilde{g}} = \lambda \widetilde{g}$ for any $\lambda \in \mathbb{R}$, although it satisfies the weaker condition of constant scalar curvature. Notably, the evolution does not asymptote to a hyperbolic geometry even when $M$ admits one. Thus, the asymptotic geometry is generally not Einstein in the presence of a positive cosmological constant.
\end{remark}

\subsection{Proof of Corollary \ref{nocollapse}}
\label{noncollapse}
\noindent \begin{proof}
\noindent Recall that the space of Riemannian metrics on a closed 3-manifold \(M\) satisfying the uniform bounds
\[
|\mathrm{Riem}(g)| \leq C, \quad \operatorname{Vol}(M,g) \geq v > 0, \quad \operatorname{diam}(M,g) \leq D,
\]
for fixed constants \(C, v, D > 0\), is precompact in the \(C^{1,\alpha}\) and \(L^{2,p}\) topologies for appropriate \(\alpha \in (0,1)\) and \(p > 3\).

\noindent For any \(\epsilon > 0\), define the \(\epsilon\)-thick and \(\epsilon\)-thin parts of \((M,g)\) by
\[
M^{\epsilon} := \{ x \in M \mid \operatorname{Vol}(B_x(1)) \geq \epsilon \}, \quad M_{\epsilon} := \{ x \in M \mid \operatorname{Vol}(B_x(1)) < \epsilon \},
\]
where \(B_x(1)\) is the geodesic ball of radius 1 centered at \(x\) with respect to \(g\). The Bishop-Gromov volume comparison theorem \cite{bg} implies that for any sequence of metrics \(\{g_i\}\) satisfying the uniform bounds above, the thin part is empty for sufficiently small \(\epsilon\), i.e., \(M = M^{\epsilon}\).

\noindent By Theorem \ref{main}, the uniform \(L^2\)-based energy estimates on curvature and the metric components imply uniform control on \(|\mathrm{Riem}(g(T))|\), diameter, and volume lower bounds for all \(T \geq T_0\).

\noindent Furthermore, consider the variation of the volume of a fixed geodesic ball \(B(1)\) under the CMC evolution:
\[
\left| \operatorname{Vol}(B(1))_T - \operatorname{Vol}(B(1))_{T_0} \right| = \left| \int_{T_0}^T \frac{d}{dT'} \operatorname{Vol}(B(1))_{T'} \, dT' \right|.
\]
By the first variation formula for volume under the flow, and using the uniform exponential decay estimates for the lapse function \(N\) (cf. Section \ref{energy}), we obtain
\[
\left| \operatorname{Vol}(B(1))_T - \operatorname{Vol}(B(1))_{T_0} \right| \lesssim \operatorname{Vol}(B(1))_{T_0} \int_{T_0}^T \left\| \frac{N}{n} - 1 \right\|_{L^\infty(B(t'))} dt' \lesssim e^{-2T_0}(1 - e^{-2(T - T_0)}).
\]
Thus, for sufficiently large initial time \(T_0\), the local volume \(\operatorname{Vol}(B(1))_T\) remains uniformly comparable to \(\operatorname{Vol}(B(1))_{T_0}\), establishing the uniform boundedness claimed in Corollary \ref{nocollapse}.
\end{proof}

\begin{theorem-non}[Ringström's Conjecture in the $\Lambda > 0$, $\sigma(M) \leq 0$ Setting]
Let $(\widehat{M}^{3+1}, \widehat{g})$ be a globally hyperbolic Einstein vacuum spacetime with positive cosmological constant $\Lambda > 0$, admitting a global CMC foliation by compact spacelike hypersurfaces diffeomorphic to a closed $3$-manifold $M$ of negative Yamabe type, and initial data satisfying Theorem~\ref{main}. Then there exists a CMC slice $M_{T}$ such that for every future-directed inextendible causal curve $\gamma$,
\[
M_{T} \not\subset J^{-}(\gamma).
\]
\end{theorem-non}

\begin{proof}
Let $\gamma(s) = (T(s),x^i(s))$ be any future-directed inextendible causal curve. In CMC-transported coordinates, the rescaled spacetime metric takes the form
\[
\widehat{g} = -N^2 dT^2 + g_{ij}(T) dx^i dx^j,
\]
with lapse $N$ uniformly bounded above by the estimates $\|\frac{N}{n}-1\|_{L^{\infty}(M)}\lesssim e^{-2T},~T\gg 1$. Causality implies
\[
g_{ij}(T) \frac{dx^i}{dT} \frac{dx^j}{dT} \leq N^2(T) \lesssim |\tau|^{-2}.
\]
In light of the estimates of theorem~\ref{main} and the scaling property of \ref{eq:scaling} in section \ref{mainsection}, the physical metric $\widetilde{g}(T)\approx \varphi^{-2}g(T)$ for the re-scaled nonphysical metric $g(T)$. Hence,
\[
\left\| \frac{dx^i}{dT} \right\|_{g(T)} \lesssim \bigg|\frac{\varphi}{\tau}\bigg|.
\]
Note that for the pure vacuum i.e., $\Lambda=0$, one has $\varphi=\tau$ and therefore the right hand side does not have a decay structure. For $\Lambda>0$, one observes $|\varphi(T)|\lesssim e^{-T}$ and $|\tau|= \sqrt{e^{-2T}+3\Lambda}=\sqrt{e^{-2T}+3n(n+1)}$ since we fixed $\Lambda=n(n+1)$ in our framework. 
The length of the spatial projection $\pi_M \circ \gamma$ of $\gamma$ in $(M, g(T))$ satisfies
\[
\text{Length}_{\bar{g}}[\pi_M \circ \gamma] \leq \int_{T}^\infty Cn^{-1}(n+1)^{-1} e^{-T^{'}} dT^{'} = Cn^{-1}(n+1)^{-1} e^{-T}\lesssim n^{-1}(n+1)^{-1}=\Lambda^{-1}.
\]
Thus, $\gamma$'s spatial projection $\pi_M \circ \gamma$ cannot intersect all of $M_{T}$. Therefore, $M_{T} \not\subset J^{-}(\gamma)$.
\end{proof}



\section{Funding and/or Conflicts of interests/Competing interests}
\noindent The authors have no conflicts to disclose. 

\section{Data Availability}
\noindent Data sharing is not applicable to this article as no new data were created or analyzed in this study.

\begin{thebibliography}{}

\bibitem{YCBY} A. Abrahams, A. Anderson, Y. Choquet-Bruhat, J.W. York, Geometrical hyperbolic systems for general relativity and gauge theories, \textit{Classical and Quantum Gravity}, vol. 14, 1997.      

\bibitem{andersson2004future} L. Andersson, V. Moncrief, Future complete vacuum spacetimes, \textit{The Einstein equations and the large scale behavior of gravitational fields}, 299-330, 2004, Springer.

\bibitem{andersson2011einstein} L. Andersson, V. Moncrief, Einstein spaces as attractors for the Einstein flow, \textit{Journal of differential geometry}, vol. 89, 1-47, 2011.

\bibitem{andersson2020nonlinear} L. Andersson, D. Fajman, Nonlinear stability of the Milne model with matter, \textit{Communications in Mathematical Physics}, vol. 378, 261-198, 2020.   





\bibitem{anderson2001long} M.T. Anderson, On Long-Time Evolution in General Relativity and Geometrization of 3-Manifolds, Communications in Mathematical Physics, vol. 222, 533-567, 2001.   




\bibitem{brendle} S. Brendle, Convergence of the Yamabe flow for arbitrary initial energy, Journal of Differential Geometry, vol. 69, 217-278, 2005.  



\bibitem{cheeger1} J. Cheeger, M. Gromov, Collapsing Riemannian manifolds while keeping their curvature bounded: I, Mathematical Sciences Research Institute, 1985  


\bibitem{cheeger2} J. Cheeger, M. Gromov, Collapsing Riemannian manifolds while keeping their curvature bounded: II, Journal of Differential Geometry, vol. 32, 269-298, 1990.

\bibitem{C09} D. Christodoulou, \textit{The formation of black holes in general relativity}. Monographs in Mathematics, European Mathematical Society (2009).


\bibitem{christodoulou1993global} D. Christodoulou, S. Klainerman, The global nonlinear stability of the Minkowski space, \textit{S\'eminaire Equations aux d\'eriv\'ees partielles (Polytechnique)}, pages 1-29, 1993.   

\bibitem{elliptichyperbolic} L. Andersson, V. Moncrief, Elliptic-hyperbolic systems and the Einstein equations, \textit{Annales Henri Poincar\'e}, vol. 4, pages 1-34, 2003.

\bibitem{fajman2018cmc} D. Fajman, K. Kr{\"o}ncke, On the CMC-Einstein-$\Lambda$ flow, \textit{Classical and Quantum Gravity}, vol. 35, 195005, 2018.  

\bibitem{fajman2020stable} D. Fajman, K. Kr{\"o}ncke, Stable fixed points of the Einstein flow with positive cosmological constant, \textit{Communications in Analysis and Geometry}, vol. 28, 1533-1576, 2020. 


\bibitem{fajman2020future} D. Fajman, Future attractors in 2+ 1 dimensional $\Lambda$ gravity, vol. 125, page 121102, 2020.



\bibitem{fajman2021attractors} D. Fajman, Z. Wyatt, Attractors of the Einstein-Klein-Gordon system, \textit{Communications in Partial Differential Equations}, vol. 46, pages 1-30, 2021.  

\bibitem{fajman2021stability} D. Fajman, J. Joudioux, J. Smulevici, The stability of the Minkowski space for the Einstein--Vlasov system, \textit{Analysis \& PDE}, vol. 14, 425-531, 2021.  

\bibitem{fajman2022cosmic} D. Fajman, L. Urban, Cosmic Censorship near FLRW spacetimes with negative spatial curvature, arXiv preprint arXiv:2211.08052, 2022.  




\bibitem{fischer1996quantum} A. Fischer, V. Moncrief, Quantum conformal superspace, \textit{General Relativity and Gravitation}, vol. 28, 221-237, 1996.



\bibitem{fischer2000reduced} A. Fischer, V. Moncrief, The reduced Hamiltonian of general relativity and the $\sigma$-constant of conformal geometry, Mathematical and Quantum Aspects of Relativity and Cosmology: Proceeding of the Second Samos Meeting on Cosmology, Geometry and Relativity Held at Pythagoreon, Samos, Greece, 31 August--4 September 1998, 70-101, 2000   

\bibitem{fischer2001reduced} A. Fischer, V. Moncrief, The reduced Einstein equations and the conformal volume collapse of 3-manifolds, Classical and Quantum Gravity, vol. 18, 4493-4516, 2001.   




\bibitem{friedrich1986existence1} H. Friedrich, Existence and structure of past asymptotically simple solutions of Einstein's field equations with positive cosmological constant, \textit{Journal of Geometry and Physics}, vol. 3, 101-117, 1986. 





\bibitem{friedrich1986existence} H. Friedrich, On the existence of n-geodesically complete or future complete solutions of Einstein's field equations with smooth asymptotic structure, \textit{Communications in Mathematical Physics}, vol. 107, 587-609, 1986.  






\bibitem{hadvzic2015global} M. Had{\v{z}}i{\'c}, J. Speck, The global future stability of the FLRW solutions to the Dust-Einstein system with a positive cosmological constant, \textit{Journal of Hyperbolic Differential Equations}, vol. 12, 87-188, 2015.

\bibitem{KK} V. Branding, D. Fajman, K. Kr{\"o}ncke, Stable cosmological Kaluza--Klein spacetimes, v, vol. 368, 1087-1120, 2019.    


\bibitem{klainerman} S. Klainerman, The null condition and global existence to nonlinear wave equations, In \textit{Nonlinear Systems of partial differential equations in applied mathematics, Part 1 (Santa Fe, N.M., 1984), volume 23
of Lectures in Appl. Math}., pages 293–326. Amer. Math. Soc., Providence, RI, 1986.

\bibitem{koiso1979decomposition} N. Koiso, A decomposition of the space $\mathcal{M}$ of Riemannian metrics on a manifold, \textit{Osaka Journal of Mathematics}, vol. 16, 423-429, 1979. 

\bibitem{lefloch2016global} P. G. LeFloch, Y. Ma, The global nonlinear stability of Minkowski space for self-gravitating massive fields, \textit{Communications in Mathematical Physics}, vol. 346, 603-665, 2016.  

\bibitem{lefloch2021nonlinear} P.G. LeFloch, C. Wei, Nonlinear stability of self-gravitating irrotational Chaplygin fluids in a FLRW geometry, \textit{Annales de l'Institut Henri Poincar\'e C, Analyse non lin\'eaire}, vol. 38, 787-814, 2021.  



\bibitem{ringstrom2008future} H. Ringstr{\"o}m, Future stability of the Einstein-non-linear scalar field system, Inventiones mathematicae, vol. 173, 123-208, 2008.   

\bibitem{lindblad2010global} H. Lindblad, I. Rodnianski, The global stability of Minkowski space-time in harmonic gauge, \textit{Annals of Mathematics}, pages 1401-1477, 2010.    

\bibitem{marsden1980maximal} J.E. Marsden, F.J. Tipler, Maximal hypersurfaces and foliations of constant mean curvature in general relativity, \textit{Physics Reports}, vol. 66, 109-139, 1980.  

\bibitem{marsden} R. Abraham, J. E. Marsden, T. Ratiu, Manifolds, tensor analysis, and applications, vol. 75, 2012.  

\bibitem{MM} V. Moncrief, P. Mondal, Could the universe have an exotic topology, Pure and Applied Mathematics Quarterly, vol. 15, 921-966, 2019.

\bibitem{mondal2019attractors} P. Mondal, Attractors of the `$n+1$' dimensional Einstein-$\Lambda$ flow, \textit{Classical and Quantum Gravity}, vol. 37, pages 235002, 2020. 

\bibitem{mondal2024global} P. Mondal, S-T Yau, Global exterior stability of the Minkowski space: Coupled Einstein--Yang--Mills perturbations, \textit{Journal of Mathematical Physics}, vol. 65, 2024.  

\bibitem{mondal2024nonlinear} P. Mondal, The nonlinear stability of (n+ 1)-dimensional FLRW spacetimes, \textit{Journal of Hyperbolic Differential Equations}, vol. 21, 329-422, 2024.  

\bibitem{oliynyk2016future} T.A. Oliynyk, Future stability of the FLRW fluid solutions in the presence of a positive cosmological constant, \textit{Communications in Mathematical Physics}, vol. 346, 293-312, 2016.  

\bibitem{oliynyk2021future} T. A. Oliynyk, Future Global Stability for Relativistic Perfect Fluids with Linear Equations of State $p=K\rho$ where $1/3<K<1/2$, \textit{SIAM Journal on Mathematical Analysis}, vol. 53, 4118-4141.  

\bibitem{penrose1999question} R. Penrose, The question of cosmic censorship, \textit{Journal of Astrophysics and Astronomy}, vol. 20, 233-248, 1999.  




\bibitem{perelman2002ricci} G. Perelman, The entropy formula for the Ricci flow and its geometric applications, arXiv preprint math/0211159, 2002   


\bibitem{perelman2003ricci} G. Perelman, Ricci flow with surgery on three-manifolds, arXiv preprint math/0303109, 2003   

\bibitem{perelman2003ricci2} G. Perelman, Finite extinction time for the solutions to the Ricci flow on certain three-manifolds, arXiv preprint math/0307245, 2003 

\bibitem{bg} P. Petersen, Riemannian Geometry, 1998

\bibitem{porti2008geometrization} J. Porti, Geometrization of three manifolds and Perelman’s proof, \textit{Rev. R. Acad. Cienc. Exactas F{\'\i}s. Nat. Ser. A Math. RACSAM}, vol. 102, 101-125, 2008.  

\bibitem{rodnianski2018regime} I. Rodnianski, J. Speck, A regime of linear stability for the Einstein-scalar field system with applications to nonlinear Big Bang formation, \textit{Annals of Mathematics}, pages 65-156, 2018.  

\bibitem{schoenyau} R. Schoen, S.-T. Yau, On the Structure of Manifolds with Positive Scalar Curvature, \textit{Manuscripta Mathematica}, vol. 28, 159–183, 1979.  

\bibitem{syau} R. Schoen, S.-T. Yau, Existence of incompressible minimal surfaces and the topology of three-dimensional manifolds with non-negative scalar curvature, \textit{Annals of Mathematics}, vol. 110, 127-142, 1979.  

\bibitem{wald} R. Wald, Asymptotic behavior of homogeneous cosmological models in the presence of a positive cosmological constant, Physical Review D, vol. 28, 2118, 1983.  



\bibitem{taylor} M. Taylor, The global nonlinear stability of Minkowski space for the massless Einstein-Vlasov system, \textit{Anal. PDE}, vol. 3, 9, 2017.  

\bibitem{wang} J. Wang, Future stability of the $3+1$ Milne model for the Einstein-Klein-Gordon system, \textit{Classical and Quantum Gravity}, vol. 36(22):225010, 2019.  

\bibitem{YM} K. Uhlenbeck, S. T. Yau, Heat flow for Yang-Mills-Higgs fields, Part I, \textit{Communications in Analysis and Geometry}, vol. 4, 1–33, 1996.


\bibitem{ye1993ricci} R. Ye, Ricci flow, Einstein metrics and space forms, Transactions of the American Mathematical Society, vol. 338, 871-896, 1993.   

\bibitem{york} J. W. York Jr., Conformally invariant orthogonal decomposition of symmetric tensors on Riemannian manifolds and the initial-value problem of general relativity, Journal of Mathematical Physics, vol. 14, 456-464, 1973.



\end{thebibliography}

\end{document}

\end{document}
\grid
\grid
